% Generated mechanically from the frozen cotangent.tex target.
% Do not edit this generated file; edit the bound locale source instead.

% OK, start here.
%

\setcounter{chapter}{91}
\chapter{余切复形}



\phantomsection
\label{cotangent-section-phantom}


\section{引言}
\label{cotangent-section-introduction}

\noindent
本章的目标是构造环同态、概形态射以及代数空间态射的余切复形。
可供参考的文献包括讲义 \cite{quillenhomology}、论文
\cite{quillencohomology}，以及专著 \cite{Andre} 和 \cite{cotangent}。




\section{给读者的建议}
\label{cotangent-section-advice-reader}

\noindent
撰写本章时，我们尽量减少单纯形技巧的使用。对环 $A$ 上的环 $B$，
我们把选择一个预解 $P_\bullet$ 看作计算范畴 $\mathcal{C}_{B/A}$
上阿贝尔层同调的工具，见注 \ref{cotangent-remark-resolution}。这类似于用
{\v C}ech 复形计算上同调时“良覆盖”所起的作用。若要简要了解
范畴上的同调，请参阅《位点上的上同调》第
\ref{sites-cohomology-section-homology} 节。对于同调而言，导出下惊叹号
函子 $L\pi_!$ 所扮演的角色，正如
$R\Gamma(\mathcal{C}_{B/A}, -)$ 对上同调所扮演的角色。第
\ref{cotangent-section-compute-L-pi-shriek} 节研究的范畴
$\mathcal{C}_{B/A}$，是由分解 $A \to P \to B$（其中 $P$ 是 $A$
上的多项式代数）组成的范畴的反范畴。该范畴上有环层同态
$$
\underline{A} \longrightarrow \mathcal{O} \longrightarrow \underline{B}
$$
并且在对象 $U = (P \to B)$ 上有 $\mathcal{O}(U)=P$。最终可得 $B$
在 $A$ 上的余切复形为
$$
L_{B/A} =
L\pi_!(\Omega_{\mathcal{O}/\underline{A}} \otimes_\mathcal{O} \underline{B})
$$
见引理 \ref{cotangent-lemma-compute-cotangent-complex}。我们始终尽力采用这一观点
证明环同态余切复形的基本性质。尤其是，所有结果原则上都可在不依赖
标准预解存在性的情况下证明，尽管我们并未这样做。这一理论相当令人
满意；唯一可能的例外是基本三角（命题 \ref{cotangent-proposition-triangle}）
的证明需要稍多一点关于导出下惊叹号函子的理论。为向读者提供另一种
途径，我们在注 \ref{cotangent-remark-triangle} 和 \ref{cotangent-remark-explicit-map}
中给出一种仅依赖标准预解简单性质的、相当完整的证明纲要。

\medskip\noindent
对于环化拓扑斯态射、概形态射、代数空间态射等的余切复形，我们的做法是
尽可能从上面讨论的“普通环同态”情形推出结论。





\section{环同态的余切复形}
\label{cotangent-section-cotangent-ring-map}

\noindent
设 $A$ 为环，以 $\textit{Alg}_A$ 记 $A$-代数范畴。考虑一对伴随函子
$(U,V)$，其中 $V:\textit{Alg}_A\to\textit{Sets}$ 是遗忘函子，
$U:\textit{Sets}\to\textit{Alg}_A$ 则把集合 $E$ 送到 $A$ 上以
$E$ 为变量集的多项式代数 $A[E]$。令 $X_\bullet$ 为《单纯形》第
\ref{simplicial-section-standard} 节所构造的
$\text{Fun}(\textit{Alg}_A,\textit{Alg}_A)$ 中的单纯形对象。

\medskip\noindent
取一个 $A$-代数 $B$，以 $P_\bullet=X_\bullet(B)$ 记所得的单纯形
$A$-代数。回忆 $P_0=A[B]$、$P_1=A[A[B]]$，依此类推。特别地，
每一项 $P_n$ 都是多项式 $A$-代数。还记得存在增广
$$
\epsilon : P_\bullet \longrightarrow B
$$
这里把 $B$ 看作常值单纯形 $A$-代数。

\begin{definition}
\label{cotangent-definition-standard-resolution}
设 $A\to B$ 为环同态。$B$ 在 $A$ 上的{\it 标准预解}是增广
$\epsilon:P_\bullet\to B$，其各项为
$$
P_0 = A[B],\quad P_1 = A[A[B]],\quad \ldots
$$
而同态按《单纯形》例 \ref{simplicial-example-polynomial-algebra-maps}
中的方式构造。
\end{definition}

\noindent
我们将看到，在某些情形下可以利用标准预解计算左导出函子。

\begin{definition}
\label{cotangent-definition-cotangent-complex-ring-map}
环同态 $A\to B$ 的{\it 余切复形} $L_{B/A}$，是与单纯形
$B$-模
$$
\Omega_{P_\bullet/A} \otimes_{P_\bullet, \epsilon} B
$$
相伴的 $B$-模复形，其中 $\epsilon:P_\bullet\to B$ 是 $B$
在 $A$ 上的标准预解。
\end{definition}

\noindent
在《单纯形》第 \ref{simplicial-section-complexes} 节中，我们为单纯形模
配上链复形；但这里使用的是上链复形。因此，次数 $-n$ 上的项
$L_{B/A}^{-n}$ 是 $B$-模
$\Omega_{P_n/A}\otimes_{P_n,\epsilon_n}B$，且当 $m>0$ 时
$L_{B/A}^m=0$。

\begin{remark}
\label{cotangent-remark-variant-cotangent-complex}
设 $A\to B$ 为环同态。令 $\mathcal{A}$ 为 $A$-代数箭头
$\psi:C\to B$ 所成的范畴，令 $\mathcal{S}$ 为同态 $E\to B$
所成的范畴，其中 $E$ 是集合。存在伴随函子
$V:\mathcal{A}\to\mathcal{S}$（遗忘函子）与
$U:\mathcal{S}\to\mathcal{A}$；后者把 $E\to B$ 送到
$A[E]\to B$。令 $X_\bullet$ 为《单纯形》第
\ref{simplicial-section-standard} 节所构造的
$\text{Fun}(\mathcal{A},\mathcal{A})$ 中的单纯形对象。图
$$
\xymatrix{
\mathcal{A} \ar[d] \ar[r] & \mathcal{S} \ar@<1ex>[l] \ar[d] \\
\textit{Alg}_A \ar[r] & \textit{Sets} \ar@<1ex>[l]
}
$$
交换。因此，$X_\bullet(\text{id}_B:B\to B)$ 等于 $B$
在 $A$ 上的标准预解。
\end{remark}

\begin{lemma}
\label{cotangent-lemma-colimit-cotangent-complex}
设 $A_i\to B_i$ 是有向指标集 $I$ 上的一族环同态。那么
$\colim L_{B_i/A_i}=L_{\colim B_i/\colim A_i}$。
\end{lemma}

\begin{proof}
这是因为遗忘函子 $V:A\textit{-Alg}\to\textit{Sets}$ 及其伴随函子
$U:\textit{Sets}\to A\textit{-Alg}$ 都与滤余极限交换。此外，函子
$B/A\mapsto\Omega_{B/A}$ 也具有这一性质（《交换代数》引理
\ref{algebra-lemma-colimit-differentials}）。
\end{proof}





\section{单纯形预解与导出下惊叹号函子}
\label{cotangent-section-compute-L-pi-shriek}

\noindent
设 $A\to B$ 为环同态。考虑如下范畴：其对象是 $A$-代数同态
$\alpha:P\to B$，其中 $P$ 是 $A$ 上以某个集合\footnote{只需考虑
基数不超过 $B$ 的基数的集合。}为变量集的多项式代数；其态射
$s:(\alpha:P\to B)\to(\alpha':P'\to B)$ 是满足
$\alpha'\circ s=\alpha$ 的 $A$-代数同态 $s:P\to P'$。以
$\mathcal{C}=\mathcal{C}_{B/A}$ 记此范畴的{\bf 反范畴}。取反范畴的
原因是，我们希望把对象 $(P,\alpha)$ 看作对应于仿射概形图
$$
\xymatrix{
\Spec(B) \ar[d] \ar[r] & \Spec(P) \ar[ld] \\
\Spec(A)
}
$$
我们赋予 $\mathcal{C}$ 混沌拓扑（《位点》例
\ref{sites-example-indiscrete}）；换言之，把 $\mathcal{C}$ 视为一个
位点，其覆盖均由恒等态射给出，因而所有预层都是层。此外，在
$\mathcal{C}$ 上赋予两个环层。第一个是把对象 $(P,\alpha)$ 送到 $P$
的层 $\mathcal{O}$；第二个是常值层 $B$，记作 $\underline{B}$。
于是得到如下环化拓扑斯态射图
\begin{equation}
\label{cotangent-equation-pi}
\vcenter{
\xymatrix{
(\Sh(\mathcal{C}), \underline{B}) \ar[r]_i \ar[d]_\pi &
(\Sh(\mathcal{C}), \mathcal{O}) \\
(\Sh(*), B)
}
}
\end{equation}
态射 $i$ 在底层拓扑斯上是恒等态射，而
$i^\sharp:\mathcal{O}\to\underline{B}$ 是显然的同态。同态 $\pi$
如《位点上的上同调》例
\ref{sites-cohomology-example-category-to-point} 所述。下列导出函子将在
后文中起重要作用：
$
Li^* : D(\mathcal{O}) \longrightarrow D(\underline{B})
$
它左伴随于 $Ri_*=i_*:D(\underline{B})\to D(\mathcal{O})$；以及
$
L\pi_! : D(\underline{B}) \longrightarrow D(B)
$
它左伴随于 $\pi^*=\pi^{-1}:D(B)\to D(\underline{B})$。

\begin{lemma}
\label{cotangent-lemma-identify-pi-shriek}
沿用上述记号，设 $P_\bullet$ 是带有增广
$\epsilon:P_\bullet\to B$ 的单纯形 $A$-代数。假设每个 $P_n$
都是 $A$ 上的多项式代数，并且 $\epsilon$ 在底层单纯形集上是平凡
Kan 纤维化。那么
$$
L\pi_!(\mathcal{F}) = \mathcal{F}(P_\bullet, \epsilon)
$$
在 $D(\textit{Ab})$ 或 $D(B)$ 中成立，并且分别关于
$\textit{Ab}(\mathcal{C})$ 或 $\textit{Mod}(\underline{B})$ 中的
$\mathcal{F}$ 具有函子性。
\end{lemma}

\begin{proof}
我们将使用《位点上的上同调》引理
\ref{sites-cohomology-lemma-compute-by-cosimplicial-resolution} 的判据。
给定 $\mathcal{C}$ 的对象 $U=(Q,\beta)$，须证明
$$
S_\bullet = \Mor_\mathcal{C}((P_\bullet, \epsilon), (Q, \beta))
$$
与单点集同伦等价。把 $Q$ 写成 $A[E]$，其中 $E$ 为某个集合（由
$\mathcal{C}$ 的选取可如此表示）。于是
$$
S_\bullet = \Mor_{\textit{Sets}}((E, \beta|_E), (P_\bullet, \epsilon))
$$
令 $*$ 为单点集上的常值单纯形集。对 $b\in B$，以
$F_{b,\bullet}$ 记由笛卡儿图
$$
\xymatrix{
F_{b, \bullet} \ar[r] \ar[d] & P_\bullet \ar[d]_\epsilon \\
{*} \ar[r]^b & B
}
$$
定义的单纯形集。按此记号，
$S_\bullet=\prod_{e\in E}F_{\beta(e),\bullet}$。由于假设
$\epsilon$ 是平凡 Kan 纤维化，$F_{b,\bullet}\to *$ 也是平凡 Kan
纤维化（《单纯形》引理
\ref{simplicial-lemma-trivial-kan-base-change}）。因此
$S_\bullet\to *$ 是平凡 Kan 纤维化（《单纯形》引理
\ref{simplicial-lemma-product-trivial-kan}），从而 $S_\bullet$ 与 $*$
同伦等价（《单纯形》引理
\ref{simplicial-lemma-trivial-kan-homotopy}）。
\end{proof}

\noindent
特别地，可以用 $B$ 在 $A$ 上的标准预解计算导出下惊叹号函子。

\begin{lemma}
\label{cotangent-lemma-pi-shriek-standard}
设 $A\to B$ 为环同态，$\epsilon:P_\bullet\to B$ 为 $B$ 在 $A$
上的标准预解，并令 $\pi$ 如 (\ref{cotangent-equation-pi}) 所示。那么
$$
L\pi_!(\mathcal{F}) = \mathcal{F}(P_\bullet, \epsilon)
$$
在 $D(\textit{Ab})$ 或 $D(B)$ 中成立，并且分别关于
$\textit{Ab}(\mathcal{C})$ 或 $\textit{Mod}(\underline{B})$ 中的
$\mathcal{F}$ 具有函子性。
\end{lemma}

\begin{proof}[第一种证明]
应用引理 \ref{cotangent-lemma-identify-pi-shriek}。由于各项 $P_n$ 都是多项式
代数，该引理的第一个假设成立。第二个假设证明如下。由《单纯形》引理
\ref{simplicial-lemma-standard-simplicial-homotopy}，同态 $\epsilon$
在底层单纯形集上是同伦等价。由《单纯形》引理
\ref{simplicial-lemma-homotopy-equivalence}，这说明 $\epsilon$ 在相伴的
阿贝尔群复形之间诱导拟同构。再由《单纯形》引理
\ref{simplicial-lemma-qis-simplicial-abelian-groups}，$\epsilon$
在底层单纯形集上是平凡 Kan 纤维化。
\end{proof}

\begin{proof}[第二种证明]
我们将使用《位点上的上同调》引理
\ref{sites-cohomology-lemma-compute-by-cosimplicial-resolution} 的判据。
设 $U=(Q,\beta)$ 为 $\mathcal{C}$ 的对象。须证明
$$
S_\bullet = \Mor_\mathcal{C}((P_\bullet, \epsilon), (Q, \beta))
$$
与单点集同伦等价。把 $Q$ 写成 $A[E]$，其中 $E$ 是某个集合（由
$\mathcal{C}$ 的选取可如此表示）。使用注
\ref{cotangent-remark-variant-cotangent-complex} 的记号可得
$$
S_\bullet = \Mor_\mathcal{S}((E \to B), i(P_\bullet \to B))
$$
由《单纯形》引理
\ref{simplicial-lemma-standard-simplicial-homotopy}，同态
$i(P_\bullet\to B)\to i(B\to B)$ 在 $\mathcal{S}$ 中是同伦等价。
因此 $S_\bullet$ 与
$$
\Mor_\mathcal{S}((E \to B), (B \to B)) = \{*\}
$$
同伦等价，正合所需。
\end{proof}

\begin{lemma}
\label{cotangent-lemma-compute-cotangent-complex}
设 $A\to B$ 为环同态，并令 $\pi$ 与 $i$ 如 (\ref{cotangent-equation-pi})
所示。存在典范同构
$$
L_{B/A} = L\pi_!(Li^*\Omega_{\mathcal{O}/A}) =
L\pi_!(i^*\Omega_{\mathcal{O}/A}) =
L\pi_!(\Omega_{\mathcal{O}/A} \otimes_\mathcal{O} \underline{B})
$$
属于 $D(B)$。
\end{lemma}

\begin{proof}
对范畴 $\mathcal{C}$ 的对象 $\alpha:P\to B$，模 $\Omega_{P/A}$
是自由 $P$-模。因此 $\Omega_{\mathcal{O}/A}$ 是平坦
$\mathcal{O}$-模。于是
$Li^*\Omega_{\mathcal{O}/A}=i^*\Omega_{\mathcal{O}/A}$ 是这样的
$\underline{B}$-模层：它把 $\alpha:P\to A$ 送到 $B$-模
$\Omega_{P/A}\otimes_{P,\alpha}B$。
% P11-E0018: the frozen source says alpha:P->A although objects of C are alpha:P->B; retained here. See control/ERRATA_CANDIDATES.jsonl.
由引理 \ref{cotangent-lemma-pi-shriek-standard}，右端由该层在标准预解上的值
计算；而依定义 \ref{cotangent-definition-cotangent-complex-ring-map}，这正是左端
的定义。
\end{proof}

\begin{lemma}
\label{cotangent-lemma-pi-lower-shriek-constant-sheaf}
若 $A\to B$ 是环同态，且 $\pi$ 如 (\ref{cotangent-equation-pi}) 所示，则
$L\pi_!(\pi^{-1}M)=M$。
\end{lemma}

\begin{proof}
由引理 \ref{cotangent-lemma-identify-pi-shriek}，$L\pi_!(\pi^{-1}M)$ 可由
$(\pi^{-1}M)(P_\bullet,\epsilon)$ 计算，而后者是 $M$ 上的常值
单纯形对象。
\end{proof}

\begin{lemma}
\label{cotangent-lemma-identify-H0}
若 $A\to B$ 是环同态，则 $H^0(L_{B/A})=\Omega_{B/A}$。
\end{lemma}

\begin{proof}
我们直接计算。使用引理 \ref{cotangent-lemma-compute-cotangent-complex} 中的等同。
显然存在从 $\Omega_{\mathcal{O}/A}\otimes\underline{B}$ 到取值为
$\Omega_{B/A}$ 的常值层的同态。因此，它诱导同态
$$
H^0(L_{B/A}) = H^0(L\pi_!(\Omega_{\mathcal{O}/A} \otimes \underline{B}))
= \pi_!(\Omega_{\mathcal{O}/A} \otimes \underline{B}) \to \Omega_{B/A}
$$
在 $\mathcal{C}_{B/A}$ 中选取对象 $P\to B$，使 $P\to B$ 满射；
由《交换代数》引理 \ref{algebra-lemma-differential-surjective} 可知上述
同态满射。为证其单射，设 $P\to B$ 是 $\mathcal{C}_{B/A}$ 的对象，
且 $\xi\in\Omega_{P/A}\otimes_PB$ 在 $\Omega_{B/A}$ 中的像为零。
先选取分解 $P\to P'\to B$，使 $P'\to B$ 满射且 $P'$ 是 $A$
上的多项式代数。可以用 $P'$ 替换 $P$。若 $B=P/I$，则
$\Omega_{P/A}\otimes_PB\to\Omega_{B/A}$ 的核是 $I/I^2$ 的像
（《交换代数》引理 \ref{algebra-lemma-differential-seq}）。设 $\xi$
是 $f\in I$ 的像。考虑两个同态 $a,b:P'=P[x]\to P$：第一个把
$x$ 映到 $0$，第二个把 $x$ 映到 $f$（在两种情形中，
$P[x]\to B$ 都把 $x$ 映到零）。于是 $\xi$ 与 $0$ 都是
$\text{d}x\otimes1$ 在 $\Omega_{P'/A}\otimes_{P'}B$ 中的像。
故 $\xi$ 与 $0$ 在余极限
$\pi_!(\Omega_{\mathcal{O}/A}\otimes\underline{B})$ 中有相同的像
（见《位点上的上同调》例
\ref{sites-cohomology-example-category-to-point}），正合所需。
\end{proof}

\begin{lemma}
\label{cotangent-lemma-pi-lower-shriek-polynomial-algebra}
若 $B$ 是环 $A$ 上的多项式代数，且 $\pi$ 如
(\ref{cotangent-equation-pi}) 所示，则 $\pi_!$ 正合，且
$\pi_!\mathcal{F}=\mathcal{F}(B\to B)$。
\end{lemma}

\begin{proof}
由引理 \ref{cotangent-lemma-identify-pi-shriek}，可用 $B$ 上的常值单纯形代数
计算 $L\pi_!$。
\end{proof}

\begin{lemma}
\label{cotangent-lemma-cotangent-complex-polynomial-algebra}
若 $B$ 是环 $A$ 上的多项式代数，则 $L_{B/A}$ 与
$\Omega_{B/A}[0]$ 拟同构。
\end{lemma}

\begin{proof}
立即由引理 \ref{cotangent-lemma-compute-cotangent-complex} 和
\ref{cotangent-lemma-pi-lower-shriek-polynomial-algebra} 得出。
\end{proof}





\section{构造预解}
\label{cotangent-section-polynomial}

\noindent
在诺特有限型情形下，可以为有限型环同态构造一个“小”的单纯形预解。

\begin{lemma}
\label{cotangent-lemma-polynomial}
设 $A$ 是诺特环，$A\to B$ 是有限型环同态。令 $\mathcal{A}$ 为
$A$-代数同态 $C\to B$ 所成的范畴。设 $n\geq0$，并设
$P_\bullet$ 是 $\mathcal{A}$ 中满足下列条件的单纯形对象：
\begin{enumerate}
\item $P_\bullet\to B$ 是单纯形集的平凡 Kan 纤维化；
\item 当 $k\leq n$ 时，$P_k$ 在 $A$ 上有限型；
\item 作为 $\mathcal{A}$ 中的单纯形对象，
$P_\bullet=\text{cosk}_n\text{sk}_nP_\bullet$。
\end{enumerate}
那么 $P_{n+1}$ 是有限型 $A$-代数。
\end{lemma}

\begin{proof}
我们给出的证明虽然直接，却略显繁琐。为说明思路，先解释较小 $n$ 时的
情形，再给出一般证明。例如，若 $n=0$，则 (3) 意味着
$P_1=P_0\times_BP_0$。由于环同态 $P_0\to B$ 满射，由《更多交换代数》
引理 \ref{more-algebra-lemma-fibre-product-finite-type}，该代数在 $A$
上有限型。

\medskip\noindent
若 $n=1$，则 (3) 意味着
$$
P_2 = \{(f_0, f_1, f_2) \in P_1^3 \mid
d_0f_0 = d_0f_1,\ d_1f_0 = d_0f_2,\ d_1f_1 = d_1f_2 \}
$$
其中各等式均在 $P_0$ 中成立。注意，三元组
$$
(d_0f_0, d_1f_0, d_1f_1) = (d_0f_1, d_0f_2, d_1f_2)
$$
是 $B$ 上纤维积 $P_0\times_BP_0\times_BP_0$ 的元素，因为同态
$d_i:P_1\to P_0$ 都是 $B$ 上的态射。因此得到同态
$$
\psi : P_2 \longrightarrow P_0 \times_B P_0 \times_B P_0
$$
$\psi$ 在元素
$(g_0,g_1,g_2)\in P_0\times_BP_0\times_BP_0$ 上的纤维，是满足
$(d_0,d_1)(f_0)=(g_0,g_1)$、$(d_0,d_1)(f_1)=(g_0,g_2)$ 以及
$(d_0,d_1)(f_2)=(g_1,g_2)$ 的 $1$-单形三元组
$(f_0,f_1,f_2)$ 所成的集合。由于 $P_\bullet\to B$ 是平凡 Kan
纤维化，同态 $(d_0,d_1):P_1\to P_0\times_BP_0$ 满射。因此
$P_2$ 位于笛卡儿图
$$
\xymatrix{
P_2 \ar[d] \ar[r] & P_1^3 \ar[d] \\
P_0 \times_B P_0 \times_B P_0 \ar[r] & (P_0 \times_B P_0)^3
}
$$
中。由《更多交换代数》引理
\ref{more-algebra-lemma-formal-consequence} 即得结论。一般情形类似，
但需要稍多一些记号。

\medskip\noindent
现在处理 $n>1$ 的情形。由《单纯形》引理
\ref{simplicial-lemma-cosk-above-object}，条件
$P_\bullet=\text{cosk}_n\text{sk}_nP_\bullet$ 蕴含同一等式在单纯形
$A$-代数范畴中成立，因而也在集合范畴中成立（因为从 $A$-代数到集合
的遗忘函子与极限交换）。所以
$$
P_{n + 1} =
\Mor(\Delta[n + 1], P_\bullet) =
\Mor(\text{sk}_n \Delta[n + 1], \text{sk}_n P_\bullet)
$$
这里使用了《单纯形》引理 \ref{simplicial-lemma-simplex-map} 和等式
(\ref{simplicial-equation-cosk})。我们将对
$1\leq k<m\leq n+1$ 归纳证明环
$$
Q_{k, m} = \Mor(\text{sk}_k \Delta[m], \text{sk}_k P_\bullet)
$$
在 $A$ 上有限型。当 $k=1$ 且 $1<m\leq n+1$ 时，这与上面对
$n=1$ 的讨论完全类似。具体而言，存在笛卡儿图
$$
\xymatrix{
Q_{1, m} \ar[d] \ar[r] & P_1^N \ar[d] \\
P_0 \times_B \ldots \times_B P_0 \ar[r] & (P_0 \times_B P_0)^N
}
$$
其中 $N={m+1\choose2}$。同前可得结论。

\medskip\noindent
设 $1\leq k_0\leq n$，并假设对所有满足
$1\leq k\leq k_0$ 与 $k<m\leq n+1$ 的指标，$Q_{k,m}$ 在 $A$
上有限型。对 $k_0+1<m\leq n+1$，我们断言存在笛卡儿方块
$$
\xymatrix{
Q_{k_0 + 1, m} \ar[d] \ar[r] & P_{k_0 + 1}^N \ar[d] \\
Q_{k_0, m} \ar[r] & Q_{k_0, k_0 + 1}^N
}
$$
其中 $N$ 是 $\Delta[m]$ 的非退化 $(k_0+1)$-单形的个数。为看出这一点，
把 $Q_{k_0+1,m}$ 的元素看作从 $\Delta[m]$ 的 $(k_0+1)$-骨架到
$P_\bullet$ 的同态 $f$。把 $f$ 限制到 $k_0$-骨架，得到图中的左侧
竖直同态；把它限制到每个非退化 $(k_0+1)$-单形，则得到上方水平
同态。给定这样的 $f$，等价于给定它在 $k_0$-骨架和每个非退化
$(k_0+1)$-面上的限制，并要求这些限制在交集上一致；这正是该图所
表达的内容。此外，由 $P_\bullet\to B$ 是平凡 Kan 纤维化可知，同态
$$
P_{k_0} \to Q_{k_0, k_0 + 1} = \Mor(\partial \Delta[k_0 + 1], P_\bullet)
$$
满射，因为当 $k_0\geq1$ 时，每个同态
$\partial\Delta[k_0+1]\to B$ 都可延拓为
$\Delta[k_0+1]\to B$（这里略去关于常值单纯形集的一小段论证）。
% P11-E0020: the frozen source uses P_{k_0} for the boundary-filling map; retained here. See control/ERRATA_CANDIDATES.jsonl.
由归纳假设，环 $Q_{k_0,m}$ 和 $Q_{k_0,k_0+1}$ 都是有限型
$A$-代数。再应用《更多交换代数》引理
\ref{more-algebra-lemma-formal-consequence}，可知
$Q_{k_0+1,m}$ 也如此。
\end{proof}

\begin{proposition}
\label{cotangent-proposition-polynomial}
设 $A$ 是诺特环，$A\to B$ 是有限型环同态。存在带有增广
$\epsilon:P_\bullet\to B$ 的单纯形 $A$-代数 $P_\bullet$，使每个
$P_n$ 都是 $A$ 上的有限型多项式代数，并且 $\epsilon$ 是单纯形集的
平凡 Kan 纤维化。
\end{proposition}

\begin{proof}
令 $\mathcal{A}$ 为 $A$-代数同态 $C\to B$ 所成的范畴。本证明中的
单纯形对象、骨架函子和余骨架函子均取在此范畴中。

\medskip\noindent
选取 $A$ 上的有限型多项式代数 $P_0$ 以及满射 $P_0\to B$。作为
第一次近似，取 $P_\bullet=\text{cosk}_0(P_0)$。换言之，
$P_\bullet$ 是各项为
$P_n=P_0\times_A\ldots\times_AP_0$ 的单纯形 $A$-代数。
% P11-E0019: the frozen source uses fibre products over A in this coskeleton; retained here. See control/ERRATA_CANDIDATES.jsonl.
（在证明的最后一段，这个单纯形对象将记作 $P^0_\bullet$。）由
《单纯形》引理 \ref{simplicial-lemma-cosk-minus-one-equivalence}，同态
$P_\bullet\to B$ 是单纯形集的平凡 Kan 纤维化。还要注意
$P_\bullet=\text{cosk}_0\text{sk}_0P_\bullet$。

\medskip\noindent
假设对某个 $n\geq0$ 已构造出 $P_\bullet$（在证明的最后一段，它将
记作 $P^n_\bullet$），满足
\begin{enumerate}
\item[(a)] $P_\bullet\to B$ 是单纯形集的平凡 Kan 纤维化；
\item[(b)] 当 $0\leq k\leq n$ 时，$P_k$ 是有限生成多项式代数；
以及
\item[(c)] $P_\bullet = \text{cosk}_n \text{sk}_n P_\bullet$
\end{enumerate}
由引理 \ref{cotangent-lemma-polynomial}，可以找到 $A$ 上的有限生成多项式代数
$Q$ 以及满射 $Q\to P_{n+1}$。由于 $P_n$ 是多项式代数，$A$-代数
同态 $s_i:P_n\to P_{n+1}$ 可提升为同态 $s'_i:P_n\to Q$。令
$d'_j:Q\to P_n$ 等于 $Q\to P_{n+1}$ 与
$d_j:P_{n+1}\to P_n$ 的复合。规定当 $k\leq n$ 时
$P'_k=P_k$，并令 $P'_{n+1}=Q$；在次数 $k\leq n-1$ 上取态射
$d'_i=d_i$ 和 $s'_i=s_i$，在次数 $n$ 上使用态射 $d'_j$ 与
$s'_i$，由此得到 $\mathcal{A}$ 中的截断单纯形对象 $P'_\bullet$。
利用 $\text{cosk}_{n+1}$ 把它延拓为 $\mathcal{A}$ 中的完整单纯形
对象。由余骨架函子的函子性，存在单纯形对象态射
$P'_\bullet\to P_\bullet$，延拓给定的 $(n+1)$-截断单纯形对象态射。
（在证明的最后一段，此态射将记作
$P^{n+1}_\bullet\to P^n_\bullet$。）

\medskip\noindent
注意，把 $n$ 替换为 $n+1$ 后，$P'_\bullet$ 满足条件 (b) 与 (c)。
我们断言同态 $P'_\bullet\to P_\bullet$ 满足《单纯形》引理
\ref{simplicial-lemma-section} 中把 $n$ 替换为 $n+1$ 后的假设
(1)、(2)、(3) 与 (4)。条件 (1) 和 (2) 由构造成立。由《单纯形》
引理 \ref{simplicial-lemma-cosk-above-object}，有
$P_\bullet = \text{cosk}_{n + 1}\text{sk}_{n + 1}P_\bullet$
且
$P'_\bullet = \text{cosk}_{n + 1}\text{sk}_{n + 1}P'_\bullet$
这些等式不仅在 $\mathcal{A}$ 中成立，也在 $A$-代数范畴中成立，
从而在集合范畴中成立（因为从 $A$-代数到集合的遗忘函子与所有极限
交换）。这证明了 (3) 与 (4)。因此该引理可用，且
$P'_\bullet\to P_\bullet$ 是平凡 Kan 纤维化。由《单纯形》引理
\ref{simplicial-lemma-trivial-kan-composition}，可知
$P'_\bullet\to B$ 是平凡 Kan 纤维化，故 (a) 也成立。

\medskip\noindent
为完成证明，取单纯形代数序列
$$
\ldots \to P^2_\bullet \to P^1_\bullet \to P^0_\bullet
$$
的逆极限 $P_\bullet=\lim P^n_\bullet$。由《单纯形》引理
\ref{simplicial-lemma-limit-trivial-kan}，同态 $P_\bullet\to B$ 是
平凡 Kan 纤维化。另一方面，上述构造在每个次数上最终稳定为一个固定的
有限生成多项式代数，正合所需。
\end{proof}

\begin{lemma}
\label{cotangent-lemma-pi-shriek-finite}
设 $A$ 是诺特环，$A\to B$ 是有限型环同态，并令 $\pi$、
$\underline{B}$ 如 (\ref{cotangent-equation-pi}) 所示。若 $\mathcal{F}$ 是
$\underline{B}$-模，并且对每个
$\alpha:P=A[x_1,\ldots,x_n]\to B$，$\mathcal{F}(P,\alpha)$ 都是
有限 $B$-模，则 $L\pi_!(\mathcal{F})$ 的各上同调模都是有限 $B$-模。
\end{lemma}

\begin{proof}
由引理 \ref{cotangent-lemma-identify-pi-shriek} 与命题
\ref{cotangent-proposition-polynomial}，可以用一个由 $\mathcal{F}$ 在有限型
多项式代数上的取值构成的复形来计算 $L\pi_!(\mathcal{F})$。
\end{proof}

\begin{lemma}
\label{cotangent-lemma-cotangent-finite}
设 $A$ 是诺特环，$A\to B$ 是有限型环同态。那么对所有
$n\in\mathbf{Z}$，$H^n(L_{B/A})$ 都是有限 $B$-模。
\end{lemma}

\begin{proof}
应用引理 \ref{cotangent-lemma-compute-cotangent-complex} 和
\ref{cotangent-lemma-pi-shriek-finite}。
\end{proof}

\begin{remark}[预解]
\label{cotangent-remark-resolution}
设 $A\to B$ 为任意环同态。若增广单纯形 $A$-代数
$\epsilon:P_\bullet\to B$ 的每个 $P_n$ 都是多项式代数，并且
$\epsilon$ 是单纯形集的平凡 Kan 纤维化，则称它为 $B$ 在 $A$ 上的
{\it 预解}。若 $P_\bullet\to B$ 是单纯形 $A$-代数的增广，每个
$P_n$ 都是满射到 $B$ 的多项式代数，则下列条件等价：
\begin{enumerate}
\item $\epsilon:P_\bullet\to B$ 是 $B$ 在 $A$ 上的预解；
\item $\epsilon:P_\bullet\to B$ 在相伴复形上是拟同构；
\item $\epsilon:P_\bullet\to B$ 诱导单纯形集的同伦等价。
\end{enumerate}
为此应用《单纯形》中的引理
\ref{simplicial-lemma-trivial-kan-homotopy},
\ref{simplicial-lemma-homotopy-equivalence}，以及
\ref{simplicial-lemma-qis-simplicial-abelian-groups}.
$B$ 在 $A$ 上的预解 $P_\bullet$ 给出 $\mathcal{C}_{B/A}$ 中的余单纯形
对象 $U_\bullet$，如《位点上的上同调》引理
\ref{sites-cohomology-lemma-compute-by-cosimplicial-resolution} 所述，
并且
$$
L\pi_!\mathcal{F} = \mathcal{F}(P_\bullet)
$$
关于 $\mathcal{F}$ 具有函子性，见引理
\ref{cotangent-lemma-identify-pi-shriek}。命题 \ref{cotangent-proposition-polynomial} 证明的
形式部分表明预解存在。在引理 \ref{cotangent-lemma-pi-shriek-standard} 的第一种
证明中，我们还看到 $B$ 在 $A$ 上的标准预解确实是一个预解（故此术语
不会造成冲突）。不过，命题 \ref{cotangent-proposition-polynomial} 的证明无需
借助《单纯形》第 \ref{simplicial-section-standard} 节中的单纯形计算，
便已说明预解存在。而且，对{\it 任意}预解选择，均有典范同构
$$
L_{B/A} = \Omega_{P_\bullet/A} \otimes_{P_\bullet, \epsilon} B
$$
属于 $D(B)$；这是引理 \ref{cotangent-lemma-compute-cotangent-complex} 的结论。
任意选择预解的自由往往很有用。
\end{remark}

\begin{lemma}
\label{cotangent-lemma-O-homology-B-homology}
设 $A\to B$ 为环同态，并令 $\pi$、$\mathcal{O}$、$\underline{B}$
如 (\ref{cotangent-equation-pi}) 所示。对任意 $\mathcal{O}$-模 $\mathcal{F}$，有
$$
L\pi_!(\mathcal{F}) = L\pi_!(Li^*\mathcal{F}) =
L\pi_!(\mathcal{F} \otimes_\mathcal{O}^\mathbf{L} \underline{B})
$$
属于 $D(\textit{Ab})$。
\end{lemma}

\begin{proof}
只需验证《位点上的上同调》引理
\ref{sites-cohomology-lemma-O-homology-qis} 的假设对
$\mathcal{C}_{B/A}$ 上的 $\mathcal{O}\to\underline{B}$ 成立。以下不再
另行说明地使用注 \ref{cotangent-remark-resolution} 的结果。选取 $B$ 在 $A$
上的一个预解 $P_\bullet$，从而得到 $\mathcal{C}_{B/A}$ 中适当的余单纯形
对象 $U_\bullet$。由于 $P_\bullet\to B$ 在相伴的阿贝尔群复形上诱导
拟同构，故 $L\pi_!\mathcal{O}=B$。另一方面，
$L\pi_!\underline{B}$ 由 $\underline{B}(U_\bullet)=B$ 计算。这验证了
《位点上的上同调》引理
\ref{sites-cohomology-lemma-O-homology-qis} 的第二个假设，证明完成。
\end{proof}

\begin{lemma}
\label{cotangent-lemma-apply-O-B-comparison}
设 $A\to B$ 为环同态，并令 $\pi$、$\mathcal{O}$、$\underline{B}$
如 (\ref{cotangent-equation-pi}) 所示。则
$$
L\pi_!(\mathcal{O}) = L\pi_!(\underline{B}) = B
\quad\text{且}\quad
L_{B/A} = L\pi_!(\Omega_{\mathcal{O}/A} \otimes_\mathcal{O} \underline{B}) =
L\pi_!(\Omega_{\mathcal{O}/A})
$$
属于 $D(\textit{Ab})$。
\end{lemma}

\begin{proof}
这只是引理 \ref{cotangent-lemma-O-homology-B-homology} 的应用（而右侧第一个
等式是引理 \ref{cotangent-lemma-compute-cotangent-complex}）。
\end{proof}

\noindent
下面是基本三角的一个容易证明的特例。

\begin{lemma}
\label{cotangent-lemma-special-case-triangle}
设 $A\to B\to C$ 为环同态。若 $B$ 是 $A$ 上的多项式代数，则
$D(C)$ 中存在可区别三角
$L_{B/A} \otimes_B^\mathbf{L} C \to L_{C/A} \to L_{C/B} \to
L_{B/A} \otimes_B^\mathbf{L} C[1]$。
\end{lemma}

\begin{proof}
以下不再另行说明地使用注 \ref{cotangent-remark-resolution} 的观察。选取 $C$
在 $B$ 上的预解 $\epsilon:P_\bullet\to C$（例如标准预解）。因为
$B$ 是 $A$ 上的多项式代数，$P_\bullet$ 也是 $C$ 在 $A$ 上的预解。
因此 $L_{C/A}$ 由
$\Omega_{P_\bullet/A}\otimes_{P_\bullet,\epsilon}C$ 计算，而
$L_{C/B}$ 由 $\Omega_{P_\bullet/B}\otimes_{P_\bullet,\epsilon}C$
计算。对每个 $n$ 都有短正合列
$0 \to \Omega_{B/A} \otimes_B P_n \to \Omega_{P_n/A} \to \Omega_{P_n/B} \to 0$
（《交换代数》引理 \ref{algebra-lemma-ses-formally-smooth}），并且
$L_{B/A}=\Omega_{B/A}[0]$（引理
\ref{cotangent-lemma-cotangent-complex-polynomial-algebra}），故结论成立。
\end{proof}

\begin{example}
\label{cotangent-example-resolution-length-two}
设 $A\to B$ 为环同态。本例构造一个“显式”的、长度为 $2$ 的 $B$
在 $A$ 上的预解 $P_\bullet$。我们遵循命题
\ref{cotangent-proposition-polynomial} 证明中的步骤；另见注
\ref{cotangent-remark-resolution} 的讨论。

\medskip\noindent
选取满射 $P_0=A[u_i]\to B$，其中 $u_i$ 是一族变量。选取理想
$\Ker(P_0\to B)$ 的生成元 $f_t\in P_0$，$t\in T$。令
$P_1=A[u_i,x_t]$；面同态 $d_0$ 和 $d_1$ 是满足
$d_j(u_i)=u_i$、$d_0(x_t)=0$ 以及 $d_1(x_t)=f_t$ 的唯一
$A$-代数同态。同态 $s_0:P_0\to P_1$ 是满足 $s_0(u_i)=u_i$
的唯一 $A$-代数同态。显然
$$
P_1 \xrightarrow{d_0 - d_1} P_0 \to B \to 0
$$
正合；特别地，同态 $(d_0,d_1):P_1\to P_0\times_BP_0$ 满射。因此，
若以 $P_\bullet$ 记由 $P_0$、$P_1$、$d_0$、$d_1$ 与 $s_0$ 给出的
$1$-截断单纯形 $A$-代数，则增广
$\text{cosk}_1(P_\bullet)\to B$ 是平凡 Kan 纤维化。命题
\ref{cotangent-proposition-polynomial} 证明步骤的下一步，是选取多项式代数
$P_2$ 以及满射
$$
P_2 \longrightarrow \text{cosk}_1(P_\bullet)_2
$$
回忆
$$
\text{cosk}_1(P_\bullet)_2 =
\{(g_0, g_1, g_2) \in P_1^3 \mid d_0(g_0) = d_0(g_1),
d_1(g_0) = d_0(g_2), d_1(g_1) = d_1(g_2)\}
$$
把 $g_i\in P_1$ 看作关于 $x_t$ 的多项式，则这些条件为
$$
g_0(0) = g_1(0),\quad
g_0(f_t) = g_2(0),\quad
g_1(f_t) = g_2(f_t)
$$
因此 $\text{cosk}_1(P_\bullet)_2$ 包含元素
$y_t=(x_t,x_t,f_t)$ 和 $z_t=(0,x_t,x_t)$。
$\text{cosk}_1(P_\bullet)_2$ 中每个元素 $G$ 都可写成
$G=H+(0,0,g)$，其中 $H$ 属于同态
$A[u_i,y_t,z_t]\to\text{cosk}_1(P_\bullet)_2$ 的像，而
$g\in P_1$ 是常数项为零且在 $P_0$ 中满足 $g(f_t)=0$ 的多项式。
注意，
\begin{enumerate}
\item $g = x_t x_{t'} - f_t x_{t'}$；以及
\item 若 $r_t \in P_0$ 且 $P_0$ 中有 $\sum r_t f_t = 0$，则
$g = \sum r_t x_t$。
\end{enumerate}
都是 $P_1$ 中具有所需形式的元素。令
$$
Rel = \Ker(\bigoplus\nolimits_{t \in T} P_0 \longrightarrow P_0),\quad
(r_t) \longmapsto \sum r_tf_t
$$
令 $P_2=A[u_i,y_t,z_t,v_r,w_{t,t'}]$，其中 $r=(r_t)\in Rel$，
并取同态
$$
P_2 \longrightarrow \text{cosk}_1(P_\bullet)_2
$$
它由下列各式给出：$y_t\mapsto(x_t,x_t,f_t)$、
$z_t\mapsto(0,x_t,x_t)$、$v_r\mapsto(0,0,\sum r_tx_t)$，以及
$w_{t,t'}\mapsto(0,0,x_tx_{t'}-f_tx_{t'})$。一次计算（略去）表明
该同态满射。上面显示的同态选择确定了
$d_0,d_1,d_2:P_2\to P_1$。最后，命题
\ref{cotangent-proposition-polynomial} 证明中的步骤要求我们选取同态
$s_0,s_1:P_1\to P_2$，它们提升两个同态
$P_1\to\text{cosk}_1(P_\bullet)_2$。显然，可以取 $s_i$ 为由
$s_0(x_t)=y_t$ 与 $s_1(x_t)=z_t$ 确定的唯一 $A$-代数同态。
\end{example}








\section{函子性}
\label{cotangent-section-functoriality}

\noindent
本节考虑环同态的交换方块
\begin{equation}
\label{cotangent-equation-commutative-square}
\vcenter{
\xymatrix{
B \ar[r] & B' \\
A \ar[u] \ar[r] & A' \ar[u]
}
}
\end{equation}
我们断言，与此图相伴有典范的 $B$-线性复形同态
$$
L_{B/A} \longrightarrow L_{B'/A'}
$$
事实上，若 $P_\bullet\to B$ 是 $B$ 在 $A$ 上的标准预解，且
$P'_\bullet\to B'$ 是 $B'$ 在 $A'$ 上的标准预解，则存在与增广
$P_\bullet\to B$ 和 $P'_\bullet\to B'$ 相容的典范单纯形
$A$-代数同态 $P_\bullet\to P'_\bullet$。这可从《单纯形》第
\ref{simplicial-section-standard} 节对标准预解的构造看出；不过在当前
特例中，只需指出同态
$$
P_0 = A[B] \longrightarrow A'[B'] = P'_0,\quad
P_1 = A[A[B]] \longrightarrow A'[A'[B']] = P'_1,
$$
等等均由给定同态 $A\to A'$ 和 $B\to B'$ 给出。所需同态
$L_{B/A}\to L_{B'/A'}$ 随即来自相伴同态
$\Omega_{P_n/A}\to\Omega_{P'_n/A'}$。

\medskip\noindent
函子性同态还可描述如下。令
$\mathcal{C}=\mathcal{C}_{B/A}$ 与
$\mathcal{C}'=\mathcal{C}_{B'/A}'$ 为第
\ref{cotangent-section-compute-L-pi-shriek} 节所考虑的范畴。
% P11-E0021: the frozen source places the prime outside the B'/A subscript; retained here. See control/ERRATA_CANDIDATES.jsonl.
存在函子
$$
u : \mathcal{C} \longrightarrow \mathcal{C}',\quad
(P, \alpha) \longmapsto (P \otimes_A A', c \circ (\alpha \otimes 1))
$$
其中 $c:B\otimes_AA'\to B'$ 是显然的同态。按《位点上的上同调》例
\ref{sites-cohomology-example-morphism-categories} 的讨论，得到拓扑斯态射
$g:\Sh(\mathcal{C})\to\Sh(\mathcal{C}')$ 以及环化拓扑斯态射的交换图
\begin{equation}
\label{cotangent-equation-double-square}
\vcenter{
\xymatrix{
(\Sh(\mathcal{C}'), \underline{B}) \ar[d]_{\pi'} &
(\Sh(\mathcal{C}'), \underline{B'}) \ar[d]_{\pi'} \ar[l]^h &
(\Sh(\mathcal{C}), \underline{B'}) \ar[d]_\pi \ar[l]^g \\
(\Sh(*), B) &
(\Sh(*), B') \ar[l]_f &
(\Sh(*), B') \ar[l]
}
}
\end{equation}
这里 $h$ 在底层拓扑斯上是恒等态射，并在环层上由环同态 $B\to B'$
给出。由《位点上的上同调》注
\ref{sites-cohomology-remark-morphism-fibred-categories}，给定
$\mathcal{C}$ 上的 $\mathcal{F}$、$\mathcal{C}'$ 上的
$\mathcal{F}'$ 以及变换 $t:\mathcal{F}\to g^{-1}\mathcal{F}'$，
便得到典范同态
$L\pi_!(\mathcal{F})\to L\pi'_!(\mathcal{F}')$。把它应用于层
$$
\mathcal{F} : (P, \alpha) \mapsto \Omega_{P/A} \otimes_P B,\quad
\mathcal{F}' : (P', \alpha') \mapsto \Omega_{P'/A'} \otimes_{P'} B',
$$
以及由典范同态
$$
\Omega_{P/A} \otimes_P B \longrightarrow
\Omega_{P \otimes_A A'/A'} \otimes_{P \otimes_A A'} B'
$$
给出的变换 $t$，得到典范同态
$$
L\pi_!(\Omega_{\mathcal{O}/A} \otimes_\mathcal{O} \underline{B})
\longrightarrow
L\pi'_!(\Omega_{\mathcal{O}'/A'} \otimes_{\mathcal{O}'} \underline{B'})
$$
由引理 \ref{cotangent-lemma-compute-cotangent-complex}，这给出
$L_{B/A}\to L_{B'/A'}$。我们略去对此同态与上面通过单纯形预解定义的
同态一致性的验证。

\begin{lemma}
\label{cotangent-lemma-flat-base-change}
假设 (\ref{cotangent-equation-commutative-square}) 诱导拟同构
$B\otimes_A^\mathbf{L}A'=B'$。沿用
(\ref{cotangent-equation-double-square}) 的记号，并设
$\mathcal{F}'\in\textit{Ab}(\mathcal{C}')$，则
$L\pi_!(g^{-1}\mathcal{F}')=L\pi'_!(\mathcal{F}')$。
\end{lemma}

\begin{proof}
以下不再另行说明地使用注 \ref{cotangent-remark-resolution} 的结果。应用
《位点上的上同调》引理 \ref{sites-cohomology-lemma-get-it-now}。
设 $P_\bullet\to B$ 为预解。只要证明
$u(P_\bullet)=P_\bullet\otimes_AA'\to B'$ 是拟同构，结论即成立。
把 $P_\bullet$ 看作单纯形 $A$-模，则其相伴 $A$-模复形
$s(P_\bullet)$ 是 $B$ 的自由 $A$-模预解：确实，每个 $P_n$ 都是
自由 $A$-模，且 $s(P_\bullet)\to B$ 是拟同构。因此
$B\otimes_A^\mathbf{L}A'$ 由
$s(P_\bullet)\otimes_AA'=s(P_\bullet\otimes_AA')$ 计算。故本引理的
假设意味着 $\epsilon':P_\bullet\otimes_AA'\to B'$ 是拟同构。
\end{proof}

\noindent
特别地，当 $A\to A'$ 平坦且 $B'=B\otimes_AA'$ 时（平坦基变换），
下述引理适用。

\begin{lemma}
\label{cotangent-lemma-flat-base-change-cotangent-complex}
若 (\ref{cotangent-equation-commutative-square}) 诱导拟同构
$B\otimes_A^\mathbf{L}A'=B'$，则函子性同态诱导同构
$$
L_{B/A} \otimes_B^\mathbf{L} B' \longrightarrow L_{B'/A'}
$$
\end{lemma}

\begin{proof}
使用等式 (\ref{cotangent-equation-double-square}) 中引入的记号。我们有
$$
L_{B/A} \otimes_B^\mathbf{L} B' =
L\pi_!(\Omega_{\mathcal{O}/A} \otimes_\mathcal{O} \underline{B})
\otimes_B^\mathbf{L} B' =
L\pi_!(Lh^*(\Omega_{\mathcal{O}/A} \otimes_\mathcal{O} \underline{B}))
$$
第一个等式来自引理 \ref{cotangent-lemma-compute-cotangent-complex}，第二个来自
《位点上的上同调》引理 \ref{sites-cohomology-lemma-change-of-rings}。
由于 $\Omega_{\mathcal{O}/A}$ 是平坦 $\mathcal{O}$-模，
$\Omega_{\mathcal{O}/A}\otimes_\mathcal{O}\underline{B}$ 是平坦
$\underline{B}$-模。因此
$Lh^*(\Omega_{\mathcal{O}/A} \otimes_\mathcal{O} \underline{B}) =
\Omega_{\mathcal{O}/A} \otimes_\mathcal{O} \underline{B'}$
它等于
$g^{-1}(\Omega_{\mathcal{O}'/A'} \otimes_{\mathcal{O}'} \underline{B'})$
这可直接检验。再由引理 \ref{cotangent-lemma-flat-base-change}，并利用
$L_{B'/A'}$ 由
$L\pi'_!(\Omega_{\mathcal{O}'/A'}\otimes_{\mathcal{O}'}\underline{B'})$
计算，即得结论。
\end{proof}

\begin{remark}
\label{cotangent-remark-homotopy-triangle}
设给定方块 (\ref{cotangent-equation-commutative-square})，且存在使下图交换的
箭头 $\kappa:B\to A'$：
$$
\xymatrix{
B \ar[r]_\beta \ar[rd]_\kappa & B' \\
A \ar[u] \ar[r]^\alpha & A' \ar[u]
}
$$
我们断言，在此情形函子性同态 $P_\bullet\to P'_\bullet$ 与复合
$P_\bullet\to B\to A'\to P'_\bullet$ 同伦。事实上，利用 $\kappa$，
函子性同态分解为
$$
P_\bullet \to P_{A'/A', \bullet} \to P'_\bullet
$$
其中 $P_{A'/A',\bullet}$ 是 $A'$ 在 $A'$ 上的标准预解。由于 $A'$
是 $A'$ 上以空集为变量集的多项式代数，由《单纯形》引理
\ref{simplicial-lemma-standard-simplicial-homotopy} 可知，增广
$\epsilon_{A'/A'}:P_{A'/A',\bullet}\to A'$ 是单纯形环的同伦等价。
注意，该引理证明中构造的同伦逆同态
$c:A'\to P_{A'/A',\bullet}$ 正是结构同态。因此，下列两个复合
$$
\xymatrix{
P_\bullet \ar[r] &
P_{A'/A', \bullet} \ar@<1ex>[rr]^{\text{id}}
\ar@<-1ex>[rr]_{c \circ \epsilon_{A'/A'}} & &
P_{A'/A', \bullet} \ar[r] &
P'_\bullet
}
$$
就是上面讨论的两个同态，并且二者同伦（《单纯形》注
\ref{simplicial-remark-homotopy-pre-post-compose}），故得到所需结论。
由于第二个同态 $P_\bullet\to P'_\bullet$ 在
$\Omega_{P_\bullet/A}\to\Omega_{P'_\bullet/A'}$ 上诱导零同态，
可知在此情形函子性同态 $L_{B/A}\to L_{B'/A'}$ 零同伦。
\end{remark}

\begin{lemma}
\label{cotangent-lemma-cotangent-complex-product}
设 $A\to B$ 与 $A\to C$ 为环同态。那么同态
$L_{B\times C/A}\to L_{B/A}\oplus L_{C/A}$ 在
$D(B\times C)$ 中是同构。
\end{lemma}

\begin{proof}
虽然本引理可由基本三角推出，我们仍给出一个直接且初等的证明。把环同态
$A\to B\times C$ 分解为 $A\to A[x]\to B\times C$，其中
$x\mapsto(1,0)$。由引理 \ref{cotangent-lemma-special-case-triangle}，有可区别三角
$$
L_{A[x]/A} \otimes_{A[x]}^\mathbf{L} (B \times C) \to L_{B \times C/A} \to
L_{B \times C/A[x]} \to L_{A[x]/A} \otimes_{A[x]}^\mathbf{L} (B \times C)[1]
$$
属于 $D(B\times C)$。类似地，有可区别三角
$$
\begin{matrix}
L_{A[x]/A} \otimes_{A[x]}^\mathbf{L} B \to L_{B/A} \to L_{B/A[x]}
\to L_{A[x]/A} \otimes_{A[x]}^\mathbf{L} B[1] \\
L_{A[x]/A} \otimes_{A[x]}^\mathbf{L} C \to L_{C/A} \to L_{C/A[x]}
\to L_{A[x]/A} \otimes_{A[x]}^\mathbf{L} C[1]
\end{matrix}
$$
因此，只需证明 $A[x]$ 上的 $B\times C$ 的结论。注意，
$A[x]\to A[x,x^{-1}]$ 平坦，且
$(B\times C)\otimes_{A[x]}A[x,x^{-1}]
=B\otimes_{A[x]}A[x,x^{-1}]$，并且
$C\otimes_{A[x]}A[x,x^{-1}]=0$。所以由基变换（引理
\ref{cotangent-lemma-flat-base-change-cotangent-complex}），同态
$L_{B\times C/A[x]}\to L_{B/A[x]}\oplus L_{C/A[x]}$ 在逆转 $x$
后成为同构。同理可证，该同态在逆转 $x-1$ 后也成为同构。
本引理得证。
\end{proof}




\section{基本三角}
\label{cotangent-section-triangle}

\noindent
本节考虑一列环同态 $A\to B\to C$。我们的目标是证明，该三角诱导
可区别三角
\begin{equation}
\label{cotangent-equation-triangle}
L_{B/A} \otimes_B^\mathbf{L} C \to L_{C/A} \to L_{C/B} \to
L_{B/A} \otimes_B^\mathbf{L} C[1]
\end{equation}
属于 $D(C)$。这将在命题 \ref{cotangent-proposition-triangle} 中证明。另一种方法
见注 \ref{cotangent-remark-triangle}。

\medskip\noindent
考虑范畴 $\mathcal{C}_{C/B/A}$，它是如下范畴的{\bf 反范畴}：其对象
是 $(P\to B,Q\to C)$，其中
\begin{enumerate}
\item $P$ 是 $A$ 上的多项式代数；
\item $P\to B$ 是 $A$-代数同态；
\item $Q$ 是 $P$ 上的多项式代数；
\item $Q\to C$ 是 $P$-代数同态。
\end{enumerate}
取反范畴是因为我们希望把 $(P\to B,Q\to C)$ 看作对应于交换图
$$
\xymatrix{
\Spec(C) \ar[d] \ar[r] & \Spec(Q) \ar[d] \\
\Spec(B) \ar[d] \ar[r] & \Spec(P) \ar[dl] \\
\Spec(A)
}
$$
令 $\mathcal{C}_{B/A}$、$\mathcal{C}_{C/A}$、$\mathcal{C}_{C/B}$
为第 \ref{cotangent-section-compute-L-pi-shriek} 节所考虑的范畴。存在函子
$$
\begin{matrix}
u_1 : \mathcal{C}_{C/B/A} \to \mathcal{C}_{B/A}, &
(P \to B, Q \to C) \mapsto (P \to B) \\
u_2 : \mathcal{C}_{C/B/A} \to \mathcal{C}_{C/A}, &
(P \to B, Q \to C) \mapsto (Q \to C) \\
u_3 : \mathcal{C}_{C/B/A} \to \mathcal{C}_{C/B}, &
(P \to B, Q \to C) \mapsto (Q \otimes_P B \to C)
\end{matrix}
$$
这些函子诱导相应的拓扑斯态射 $g_i$。记
$\mathcal{O}_i=g_i^{-1}\mathcal{O}$，从而得到环化拓扑斯态射
\begin{equation}
\label{cotangent-equation-three-maps}
\begin{matrix}
g_1 : (\Sh(\mathcal{C}_{C/B/A}), \mathcal{O}_1)
\longrightarrow (\Sh(\mathcal{C}_{B/A}), \mathcal{O}) \\
g_2 : (\Sh(\mathcal{C}_{C/B/A}), \mathcal{O}_2)
\longrightarrow (\Sh(\mathcal{C}_{C/A}), \mathcal{O}) \\
g_3 : (\Sh(\mathcal{C}_{C/B/A}), \mathcal{O}_3)
\longrightarrow (\Sh(\mathcal{C}_{C/B}), \mathcal{O})
\end{matrix}
\end{equation}
记
$\pi : \Sh(\mathcal{C}_{C/B/A}) \to \Sh(*)$,
$\pi_1 : \Sh(\mathcal{C}_{B/A}) \to \Sh(*)$,
$\pi_2 : \Sh(\mathcal{C}_{C/A}) \to \Sh(*)$，以及
$\pi_3 : \Sh(\mathcal{C}_{C/B}) \to \Sh(*)$,
于是 $\pi=\pi_i\circ g_i$。利用引理
\ref{cotangent-lemma-compute-cotangent-complex} 中对余切复形的等同以及下列
引理，我们将得到所需的可区别三角。

\begin{lemma}
\label{cotangent-lemma-triangle-ses}
沿用 (\ref{cotangent-equation-three-maps}) 的记号，令
$$
\begin{matrix}
\Omega_1 = \Omega_{\mathcal{O}/A} \otimes_\mathcal{O} \underline{B}
\text{，定义在 }\mathcal{C}_{B/A}\text{ 上} \\
\Omega_2 = \Omega_{\mathcal{O}/A} \otimes_\mathcal{O} \underline{C}
\text{，定义在 }\mathcal{C}_{C/A}\text{ 上} \\
\Omega_3 = \Omega_{\mathcal{O}/B} \otimes_\mathcal{O} \underline{C}
\text{，定义在 }\mathcal{C}_{C/B}\text{ 上}
\end{matrix}
$$
那么在 $\mathcal{C}_{C/B/A}$ 上存在典范的 $\underline{C}$-模层短正合列
$$
0 \to g_1^{-1}\Omega_1 \otimes_{\underline{B}} \underline{C} \to
g_2^{-1}\Omega_2 \to
g_3^{-1}\Omega_3 \to 0
$$
它位于 $\mathcal{C}_{C/B/A}$ 上。
\end{lemma}

\begin{proof}
回忆 $g_i^{-1}$ 只需与 $u_i$ 预复合即可得到。给定对象
$U=(P\to B,Q\to C)$，有分裂短正合列
$$
0 \to \Omega_{P/A} \otimes Q \to \Omega_{Q/A} \to \Omega_{Q/P} \to 0
$$
例如这可由《交换代数》引理
\ref{algebra-lemma-ses-formally-smooth} 得出。在 $Q$ 上与 $C$
作张量积，得到短正合列
$$
0 \to \Omega_{P/A} \otimes C \to \Omega_{Q/A} \otimes C \to
\Omega_{Q/P} \otimes C \to 0
$$
有 $\Omega_{P/A}\otimes C=\Omega_{P/A}\otimes B\otimes C$，故这是
$g_1^{-1}\Omega_1\otimes_{\underline{B}}\underline{C}$ 在 $U$ 上的值。
模 $\Omega_{Q/A}\otimes C$ 是 $g_2^{-1}\Omega_2$ 在 $U$ 上的值。
由《交换代数》引理 \ref{algebra-lemma-differentials-base-change}，有
$\Omega_{Q/P}\otimes C=\Omega_{Q\otimes_PB/B}\otimes C$，故这是
$g_3^{-1}\Omega_3$ 在 $U$ 上的值。因此，本引理的短正合列来自把
最后显示的短正合列赋给每个 $U$。
\end{proof}

\begin{lemma}
\label{cotangent-lemma-polynomial-on-top}
沿用 (\ref{cotangent-equation-three-maps}) 的记号，假设 $C$ 是 $B$ 上的
多项式代数。那么对 $\mathcal{C}_{C/B}$ 上任意阿贝尔层
$\mathcal{F}$，有
$L\pi_!(g_3^{-1}\mathcal{F}) = L\pi_{3, !}\mathcal{F} = \pi_{3, !}\mathcal{F}$
\end{lemma}

\begin{proof}
把 $C$ 写成 $B[E]$，其中 $E$ 为某个集合。选取 $B$ 在 $A$ 上的
预解 $P_\bullet\to B$。对每个 $n$，考虑
$\mathcal{C}_{C/B/A}$ 的对象
$U_n=(P_n\to B,P_n[E]\to C)$。那么 $U_\bullet$ 是
$\mathcal{C}_{C/B/A}$ 中的余单纯形对象。注意，
$u_3(U_\bullet)$ 是 $\mathcal{C}_{C/B}$ 中取值为 $(C\to C)$ 的常值
余单纯形对象。我们将证明 $\mathcal{C}_{C/B/A}$ 的对象 $U_\bullet$
满足《位点上的上同调》引理
\ref{sites-cohomology-lemma-compute-by-cosimplicial-resolution} 的假设。
这便蕴含本引理，因为它说明 $L\pi_!(g_3^{-1}\mathcal{F})$ 由取值为
$\mathcal{F}(C\to C)$ 的常值单纯形阿贝尔群计算；而由引理
\ref{cotangent-lemma-pi-lower-shriek-polynomial-algebra}，这正是
$L\pi_{3,!}\mathcal{F}=\pi_{3,!}\mathcal{F}$ 的值。

\medskip\noindent
设 $U=(\beta:P\to B,\gamma:Q\to C)$ 为 $\mathcal{C}_{C/B/A}$ 的
对象。由此范畴的定义，可把 $P$ 写成 $A[S]$，把 $Q$ 写成
$A[S\amalg T]$。须证明
$$
\Mor_{\mathcal{C}_{C/B/A}}(U_\bullet, U)
$$
与单点单纯形集 $*$ 同伦等价。注意，该单纯形集是乘积
$$
\prod\nolimits_{s \in S} F_s \times \prod\nolimits_{t \in T} F'_t
$$
其中 $F_s$ 是对应于
$U_s = (A[\{s\}] \to B, A[\{s\}] \to C)$
的单纯形集，$F'_t$ 是对应于
$U_t=(A\to B,A[\{t\}]\to C)$ 的单纯形集。确实，在
$\mathcal{C}_{C/B/A}$ 中，对象 $U$ 是乘积
$\prod U_s\times\prod U_t$。只需每个 $F_s$ 与 $F'_t$ 都与 $*$
同伦等价，见《单纯形》引理
\ref{simplicial-lemma-products-homotopy}。对 $F_s$，结论来自
$P_\bullet\to B$ 作为预解是平凡 Kan 纤维化，而 $F_s$ 是该同态在
$\beta(s)$ 上的纤维。（使用《单纯形》引理
\ref{simplicial-lemma-trivial-kan-base-change} 和
\ref{simplicial-lemma-trivial-kan-homotopy}。）$F'_t$ 的情形更有意思。
这里须说明同态
$$
P_\bullet[E] \longrightarrow C = B[E]
$$
在 $\gamma(t)\in C$ 上的纤维与一点同伦等价。事实上，我们将证明该
同态是平凡 Kan 纤维化。确实，$P_\bullet\to B$ 是平凡 Kan 纤维化。
% P11-E0022: the frozen English source says “trivial can fibration”; Chinese renders the standard term Kan and the defect is logged. See control/ERRATA_CANDIDATES.jsonl.
对任意环 $R$，有
$$
R[E] =
\colim_{\Sigma \subset \text{Map}(E, \mathbf{Z}_{\geq 0})\text{ 有限}}
\prod\nolimits_{I \in \Sigma} R
$$
（滤余极限）。所以，上面显示的单纯形集同态是平凡 Kan 纤维化的滤余
极限；由《单纯形》引理
\ref{simplicial-lemma-filtered-colimit-trivial-kan}，它仍是平凡 Kan
纤维化。
\end{proof}

\begin{lemma}
\label{cotangent-lemma-triangle-compute-lower-shriek}
沿用 (\ref{cotangent-equation-three-maps}) 的记号，对 $i=1,2,3$ 有
$Lg_{i,!}\circ g_i^{-1}=\text{id}$，因而对 $i=1,2,3$ 也有
$L\pi_!\circ g_i^{-1}=L\pi_{i,!}$。
\end{lemma}

\begin{proof}
先证 $i=1$。我们断言函子 $\mathcal{C}_{C/B/A}$ 是
$\mathcal{C}_{B/A}$ 上的纤维化范畴。设给定
$(P\to B,Q\to C)$ 以及 $\mathcal{C}_{B/A}$ 中的态射
$(P'\to B)\to(P\to B)$。回忆，这意味着存在与到 $B$ 的同态相容的
$A$-代数同态 $P\to P'$。令 $Q'=Q\otimes_PP'$，并赋予其诱导的到
$C$ 的同态。于是态射
$$
(P' \to B, Q' \to C) \longrightarrow (P \to B, Q \to C)
$$
在 $\mathcal{C}_{C/B/A}$ 中（再次注意箭头反向），并且它是
$\mathcal{C}_{C/B/A}$ 中位于 $\mathcal{C}_{B/A}$ 上的强笛卡儿态射。
还要注意，$u_1$ 在 $P\to B$ 上的纤维范畴是
$\mathcal{C}_{C/P}$。令 $\mathcal{F}$ 为 $\mathcal{C}_{B/A}$ 上的
阿贝尔层。由于这是纤维化范畴，可以应用《位点上的上同调》引理
\ref{sites-cohomology-lemma-compute-left-derived-pi-shriek}。因此，
$L_ng_{1,!}g_1^{-1}\mathcal{F}$ 是把
$U\in\Ob(\mathcal{C}_{B/A})$ 送到 $g_1^{-1}\mathcal{F}$ 在 $U$
上纤维范畴的限制之第 $n$ 个同调的（预）层。由于这些限制均为常值，
结合上面对纤维范畴的等同，所需结论由引理
\ref{cotangent-lemma-pi-lower-shriek-constant-sheaf} 得出。

\medskip\noindent
现在处理 $i=2$。我们断言 $\mathcal{C}_{C/B/A}$ 是
$\mathcal{C}_{C/A}$ 上的纤维化范畴。设给定
$(P\to B,Q\to C)$ 以及 $\mathcal{C}_{C/A}$ 中的态射
$(Q'\to C)\to(Q\to C)$。回忆，这意味着存在与到 $C$ 的同态相容的
$B$-代数同态 $Q\to Q'$。于是
% P11-E0023: the frozen source calls Q->Q' a B-algebra map although C_{C/A} morphisms are A-algebra maps; retained here. See control/ERRATA_CANDIDATES.jsonl.
$$
(P \to B, Q' \to C) \longrightarrow (P \to B, Q \to C)
$$
是 $\mathcal{C}_{C/B/A}$ 中位于 $\mathcal{C}_{C/A}$ 上的强笛卡儿
态射。注意，$u_2$ 在 $Q\to C$ 上的纤维范畴有终对象（留意箭头反向），
即 $(A\to B,Q\to C)$。令 $\mathcal{F}$ 为
$\mathcal{C}_{C/A}$ 上的阿贝尔层。由于这是纤维化范畴，可以应用
《位点上的上同调》引理
\ref{sites-cohomology-lemma-compute-left-derived-pi-shriek}。因此，
$L_ng_{2,!}g_2^{-1}\mathcal{F}$ 是把
$U\in\Ob(\mathcal{C}_{C/A})$ 送到 $g_1^{-1}\mathcal{F}$ 在 $U$
上纤维范畴的限制之第 $n$ 个同调的（预）层。
% P11-E0024: frozen source uses g_1^{-1} in the i=2 case; retained here. See control/ERRATA_CANDIDATES.jsonl.
由于这些限制均为常值，且所有纤维范畴都有终对象，所需结论由
《位点上的上同调》引理 \ref{sites-cohomology-lemma-initial-final} 得出。

\medskip\noindent
最后处理 $i=3$。在此情形，把《位点上的上同调》引理
\ref{sites-cohomology-lemma-compute-left-derived-g-shriek} 应用于
$u=u_3:\mathcal{C}_{C/B/A}\to\mathcal{C}_{C/B}$ 以及
$\mathcal{F}'=g_3^{-1}\mathcal{F}$，其中 $\mathcal{F}$ 是
$\mathcal{C}_{C/B}$ 上的阿贝尔层。设
$U=(\overline{Q}\to C)$ 为 $\mathcal{C}_{C/B}$ 的对象。那么
$\mathcal{I}_U=\mathcal{C}_{\overline{Q}/B/A}$（再次留意箭头反向）。
层 $\mathcal{F}'_U$ 由规则
$(P\to B,Q\to\overline{Q})\mapsto
\mathcal{F}(Q\otimes_PB\to C)$ 给出。换言之，该层是
$\mathcal{C}_{\overline{Q}/C}$ 上一个层沿态射
$\Sh(\mathcal{C}_{\overline{Q}/B/A})
\to\Sh(\mathcal{C}_{\overline{Q}/B})$ 的拉回。
% P11-E0025: frozen source says the sheaf is on C_{Qbar/C} although the displayed morphism and application require C_{Qbar/B}; retained here. See control/ERRATA_CANDIDATES.jsonl.
因此，引理 \ref{cotangent-lemma-polynomial-on-top} 表明，当 $n>0$ 时
$H_n(\mathcal{I}_U,\mathcal{F}'_U)=0$，而当 $n=0$ 时它等于
$\mathcal{F}(\overline{Q}\to C)$。上述《位点上的上同调》引理
\ref{sites-cohomology-lemma-compute-left-derived-g-shriek} 蕴含
$Lg_{3,!}(g_3^{-1}\mathcal{F})=\mathcal{F}$，证明完成。
\end{proof}

\begin{proposition}
\label{cotangent-proposition-triangle}
设 $A\to B\to C$ 为环同态。存在典范可区别三角
$$
L_{B/A} \otimes_B^\mathbf{L} C \to L_{C/A} \to L_{C/B} \to
L_{B/A} \otimes_B^\mathbf{L} C[1]
$$
属于 $D(C)$。
\end{proposition}

\begin{proof}
考虑引理 \ref{cotangent-lemma-triangle-ses} 的层短正合列，并应用导出函子
$L\pi_!$，得到可区别三角
$$
L\pi_!(g_1^{-1}\Omega_1 \otimes_{\underline{B}} \underline{C}) \to
L\pi_!(g_2^{-1}\Omega_2) \to
L\pi_!(g_3^{-1}\Omega_3) \to
L\pi_!(g_1^{-1}\Omega_1 \otimes_{\underline{B}} \underline{C})[1]
$$
属于 $D(C)$。由引理 \ref{cotangent-lemma-triangle-compute-lower-shriek} 和
\ref{cotangent-lemma-compute-cotangent-complex}，第二项与第三项分别等于
$L_{C/A}$ 与 $L_{C/B}$，而第一项等于
$$
L\pi_{1, !}(\Omega_1 \otimes_{\underline{B}} \underline{C}) =
L\pi_{1, !}(\Omega_1) \otimes_B^\mathbf{L} C =
L_{B/A} \otimes_B^\mathbf{L} C
$$
第一个等式来自《位点上的上同调》引理
\ref{sites-cohomology-lemma-change-of-rings}（以及 $\Omega_1$ 作为
$\underline{B}$-模层的平坦性），第二个等式来自引理
\ref{cotangent-lemma-compute-cotangent-complex}。
\end{proof}

\begin{remark}
\label{cotangent-remark-triangle}
下面概述基本三角存在性的另一种、或许更简单的证明。设
$A\to B\to C$ 为环同态，并假设 $B\to C$ 单射。令
$P_\bullet\to B$ 为 $B$ 在 $A$ 上的标准预解，令
$Q_\bullet\to C$ 为 $C$ 在 $B$ 上的标准预解。图示如下：
% P11-E0026: the frozen Q-row below uses A-polynomial towers although Q is declared the standard resolution over B; retained here. See control/ERRATA_CANDIDATES.jsonl.
$$
\xymatrix{
P_\bullet : &
A[A[A[B]]] \ar[d]
\ar@<2ex>[r]
\ar@<0ex>[r]
\ar@<-2ex>[r]
&
A[A[B]] \ar[d]
\ar@<1ex>[r]
\ar@<-1ex>[r]
\ar@<1ex>[l]
\ar@<-1ex>[l]
&
A[B] \ar[d] \ar@<0ex>[l] \ar[r] &
B \\
Q_\bullet : &
A[A[A[C]]]
\ar@<2ex>[r]
\ar@<0ex>[r]
\ar@<-2ex>[r]
&
A[A[C]]
\ar@<1ex>[r]
\ar@<-1ex>[r]
\ar@<1ex>[l]
\ar@<-1ex>[l]
&
A[C] \ar@<0ex>[l] \ar[r] &
C
}
$$
注意，由于 $B\to C$ 单射，对每个 $n$，环 $Q_n$ 都是 $P_n$ 上的
多项式代数。因此得到 $\mathcal{C}_{C/B/A}$ 中的余单纯形对象（留意
箭头反向）。现在令
$\overline{Q}_\bullet=Q_\bullet\otimes_{P_\bullet}B$。命题
\ref{cotangent-proposition-triangle} 证明的关键，是说明
$\overline{Q}_\bullet$ 是 $C$ 在 $B$ 上的预解。把《位点上的上同调》
引理 \ref{sites-cohomology-lemma-O-homology-qis} 应用于
$\mathcal{C}=\Delta$、$\mathcal{O}=P_\bullet$、
$\mathcal{O}'=B$ 与 $\mathcal{F}=Q_\bullet$，即可得到此结论（这里
用到 $Q_n$ 在 $P_n$ 上平坦；关于单纯形模与层的关系，见《位点上的
上同调》注 \ref{sites-cohomology-remark-simplicial-modules}）。这一关键
事实表明，命题 \ref{cotangent-proposition-triangle} 的可区别三角，就是与下列
单纯形 $C$-模短正合列相伴的可区别三角：
$$
0 \to
\Omega_{P_\bullet/A} \otimes_{P_\bullet} C \to
\Omega_{Q_\bullet/A} \otimes_{Q_\bullet} C \to
\Omega_{\overline{Q}_\bullet/B} \otimes_{\overline{Q}_\bullet} C \to 0
$$
该短正合列由短正合列
$0 \to \Omega_{P_n/A} \otimes_{P_n} Q_n \to \Omega_{Q_n/A} \to
\Omega_{Q_n/P_n} \to 0$ 得到，见《交换代数》引理
\ref{algebra-lemma-ses-formally-smooth}。具体地，由注
\ref{cotangent-remark-resolution} 及上述关键事实，右端复形在 $D(C)$ 中表示
$L_{C/B}$。

\medskip\noindent
若 $B\to C$ 不单射，则可用上述方法为
$A\to B\to B\times C$ 得到基本三角。由于
$L_{B \times C/B} \to L_{B/B} \oplus L_{C/B}$ 以及
$L_{B \times C/A} \to L_{B/A} \oplus L_{C/A}$
在 $D(B\times C)$ 中都是拟同构（引理
\ref{cotangent-lemma-cotangent-complex-product}），再与平坦环同态
$B\times C\to C$ 作张量积，便在 $D(C)$ 中诱导所需的可区别三角。
\end{remark}

\begin{remark}
\label{cotangent-remark-explicit-map}
设 $A\to B\to C$ 为环同态，且 $B\to C$ 单射。回忆注
\ref{cotangent-remark-triangle} 中的记号
$P_\bullet$、$Q_\bullet$、$\overline{Q}_\bullet$。令
$R_\bullet$ 为 $C$ 在 $B$ 上的标准预解。本注说明如何得到
$\Omega_{\overline{Q}_\bullet/B}
\otimes_{\overline{Q}_\bullet}C$ 与
$L_{C/B}=\Omega_{R_\bullet/B}\otimes_{R_\bullet}C$ 的典范等同。
令 $S_\bullet\to B$ 为 $B$ 在 $B$ 上的标准预解。注意，由于
$B\to C$ 单射，函子性同态 $S_\bullet\to R_\bullet$ 把 $R_n$
表为 $S_n$ 上的多项式代数。例如，在次数 $0$ 上有同态
$B[B]\to B[C]$，在次数 $1$ 上有同态
$B[B[B]]\to B[B[C]]$，依此类推。因此
$\overline{R}_\bullet=R_\bullet\otimes_{S_\bullet}B$ 也是 $B$
上的单纯形多项式代数。与注 \ref{cotangent-remark-triangle} 中相同，由
《位点上的上同调》引理
\ref{sites-cohomology-lemma-O-homology-qis} 可知
$\overline{R}_\bullet\to C$ 是预解。由于存在交换图
$$
\xymatrix{
Q_\bullet \ar[r] & R_\bullet \\
P_\bullet \ar[u] \ar[r] & S_\bullet \ar[u] \ar[r] & B
}
$$
得到典范同态
$\overline{Q}_\bullet = Q_\bullet \otimes_{P_\bullet} B \to
\overline{R}_\bullet$。所以，同态
$$
L_{C/B} = \Omega_{R_\bullet/B} \otimes_{R_\bullet} C
\longrightarrow
\Omega_{\overline{R}_\bullet/B} \otimes_{\overline{R}_\bullet} C
\longleftarrow
\Omega_{\overline{Q}_\bullet/B} \otimes_{\overline{Q}_\bullet} C
$$
都是拟同构（注 \ref{cotangent-remark-resolution}）；把其中一个与另一个的逆
复合，便得到所需等同。
\end{remark}





\section{局部化与 \'etale 环同态}
\label{cotangent-section-localization}

\noindent
本节研究环局部化后的情形。设 $A\to A'\to B$ 为环同态，满足
$B=B\otimes_A^\mathbf{L}A'$。例如，当 $A'=S^{-1}A$ 是 $A$ 关于
乘法子集 $S\subset A$ 的局部化时，这一条件成立。在此情形，对
$\mathcal{C}_{B/A'}$ 上的阿贝尔层 $\mathcal{F}'$，
$g^{-1}\mathcal{F}'$ 在 $\mathcal{C}_{B/A}$ 上的同调，与
$\mathcal{F}'$ 在 $\mathcal{C}_{B/A'}$ 上的同调一致；精确陈述见
引理 \ref{cotangent-lemma-flat-base-change}。

\begin{lemma}
\label{cotangent-lemma-localize-at-bottom}
设 $A\to A'\to B$ 为环同态，满足
$B=B\otimes_A^\mathbf{L}A'$。那么在 $D(B)$ 中有
$L_{B/A}=L_{B/A'}$。
\end{lemma}

\begin{proof}
根据上面的讨论（即应用引理 \ref{cotangent-lemma-flat-base-change}）以及引理
\ref{cotangent-lemma-compute-cotangent-complex}，须证明
$\mathcal{C}_{B/A}$ 上由规则
$(P\to B)\mapsto\Omega_{P/A}\otimes_PB$ 给出的层，是由规则
$(P\to B)\mapsto\Omega_{P/A'}\otimes_PB$ 给出的层的拉回。拉回函子
$g^{-1}$ 由预复合函子
$u:\mathcal{C}_{B/A}\to\mathcal{C}_{B/A'}$ 给出，其中
$(P\to B)\mapsto(P\otimes_AA'\to B)$。因此只须证明
$$
\Omega_{P/A} \otimes_P B =
\Omega_{P \otimes_A A'/A'} \otimes_{(P \otimes_A A')} B
$$
由《交换代数》引理 \ref{algebra-lemma-differentials-base-change}，右端等于
$$
(\Omega_{P/A} \otimes_A A') \otimes_{(P \otimes_A A')} B
$$
由于 $P$ 是 $A$ 上的多项式代数，模 $\Omega_{P/A}$ 自由，故此等式
显然成立。
\end{proof}

\begin{lemma}
\label{cotangent-lemma-derived-diagonal}
设 $A\to B$ 为环同态，满足 $B=B\otimes_A^\mathbf{L}B$。那么在
$D(B)$ 中有 $L_{B/A}=0$。
\end{lemma}

\begin{proof}
由引理 \ref{cotangent-lemma-localize-at-bottom} 和
\ref{cotangent-lemma-cotangent-complex-polynomial-algebra}，有
$L_{B/A}=L_{B/B}=0$。
\end{proof}

\begin{lemma}
\label{cotangent-lemma-bootstrap}
设 $A\to B$ 为环同态，且当 $i>0$ 时
$\text{Tor}^A_i(B,B)=0$，并且 $L_{B/B\otimes_AB}=0$。那么在
$D(B)$ 中有 $L_{B/A}=0$。
\end{lemma}

\begin{proof}
由引理 \ref{cotangent-lemma-flat-base-change-cotangent-complex}，有
$L_{B/A}\otimes_B^\mathbf{L}(B\otimes_AB)=L_{B\otimes_AB/B}$。
现在对环同态 $B\to B\otimes_AB\to B$ 使用可区别三角
(\ref{cotangent-equation-triangle})：
$$
L_{B \otimes_A B/B} \otimes^\mathbf{L}_{(B \otimes_A B)} B \to
L_{B/B} \to L_{B/B \otimes_A B} \to
L_{B \otimes_A B/B} \otimes^\mathbf{L}_{(B \otimes_A B)} B[1]
$$
再利用 $L_{B/B}$ 的消失（引理
\ref{cotangent-lemma-cotangent-complex-polynomial-algebra}）以及假设中的
$L_{B/B\otimes_AB}$ 的消失，可得
$$
0 =
L_{B \otimes_A B/B} \otimes^\mathbf{L}_{(B \otimes_A B)} B =
L_{B/A} \otimes_B^\mathbf{L} (B \otimes_A B)
\otimes^\mathbf{L}_{(B \otimes_A B)} B = L_{B/A}
$$
正合所需。
\end{proof}

\begin{lemma}
\label{cotangent-lemma-when-zero}
在下列每种情形中，余切复形 $L_{B/A}$ 均为零：
\begin{enumerate}
\item $A\to B$ 与 $B\otimes_AB\to B$ 均平坦，即 $A\to B$ 弱
\'etale（《更多交换代数》定义
\ref{more-algebra-definition-weakly-etale}）；
\item $A\to B$ 是环的平坦满态射；
\item 对某个乘法子集 $S\subset A$，有 $B=S^{-1}A$；
\item $A\to B$ 非分歧且平坦；
\item $A\to B$ 是 \'etale 同态；
\item $A\to B$ 是一族余切复形消失的环同态的滤余极限；
\item $B$ 是 $A$ 的某个局部环的 Hensel 化；
\item $B$ 是 $A$ 的某个局部环的严格 Hensel 化；以及
\item 在此增加更多情形。
% P11-E0027: frozen source contains the editorial placeholder “add more here”; retained here. See control/ERRATA_CANDIDATES.jsonl.
\end{enumerate}
\end{lemma}

\begin{proof}
在情形 (1) 中，把引理 \ref{cotangent-lemma-derived-diagonal} 应用于满射平坦
环同态 $B\otimes_AB\to B$，得到 $L_{B/B\otimes_AB}=0$；再应用
引理 \ref{cotangent-lemma-bootstrap} 即得结论。情形 (2)--(5) 都是 (1) 的
特例。(6) 来自引理 \ref{cotangent-lemma-colimit-cotangent-complex}。(7) 与
(8) 来自如下事实：严格或非严格 Hensel 化是 $A$ 的 \'etale 环扩张
的滤余极限；见《交换代数》引理
\ref{algebra-lemma-henselization-different} 和
\ref{algebra-lemma-strict-henselization-different}。
\end{proof}

\begin{lemma}
\label{cotangent-lemma-localize-on-top}
设 $A\to B\to C$ 为环同态，且 $L_{C/B}=0$。那么
$L_{C/A}=L_{B/A}\otimes_B^\mathbf{L}C$。
\end{lemma}

\begin{proof}
这是可区别三角 (\ref{cotangent-equation-triangle}) 的直接推论。
\end{proof}

\begin{lemma}
\label{cotangent-lemma-localize}
设 $A\to B$ 为环同态，$S\subset A$ 与 $T\subset B$ 为乘法子集，
且 $S$ 的像包含在 $T$ 中。那么在 $D(T^{-1}B)$ 中有
$L_{T^{-1}B/S^{-1}A}=L_{B/A}\otimes_BT^{-1}B$。
\end{lemma}

\begin{proof}
引理 \ref{cotangent-lemma-localize-on-top} 给出
$L_{T^{-1}B/A}=L_{B/A}\otimes_BT^{-1}B$，引理
\ref{cotangent-lemma-localize-at-bottom} 给出
$L_{T^{-1}B/A}=L_{T^{-1}B/S^{-1}A}$。
\end{proof}

\begin{lemma}
\label{cotangent-lemma-cotangent-complex-henselization}
设 $A\to B$ 是局部环之间的局部同态。令 $A^h\to B^h$ 以及
$A^{sh}\to B^{sh}$ 分别为诱导的 Hensel 化同态与严格 Hensel 化同态。
那么
$$
L_{B^h/A^h} = L_{B^h/A} = L_{B/A} \otimes_B^\mathbf{L} B^h
\quad\text{分别}\quad
L_{B^{sh}/A^{sh}} = L_{B^{sh}/A} = L_{B/A} \otimes_B^\mathbf{L} B^{sh}
$$
分别在 $D(B^h)$ 与 $D(B^{sh})$ 中成立。
\end{lemma}

\begin{proof}
由引理 \ref{cotangent-lemma-when-zero}，复形 $L_{A^h/A}$、$L_{A^{sh}/A}$、
$L_{B^h/B}$ 与 $L_{B^{sh}/B}$ 均为零。对
$A\to B\to B^h$ 使用基本可区别三角
(\ref{cotangent-equation-triangle})，得到
$L_{B^h/A}=L_{B/A}\otimes_B^\mathbf{L}B^h$。再对
$A\to A^h\to B^h$ 使用基本三角，得到
$L_{B^h/A^h}=L_{B^h/A}$。严格 Hensel 化的情形同理。
\end{proof}




\section{光滑环同态}
\label{cotangent-section-smooth}

\noindent
设 $C\to B$ 为核是 $I$ 的环满射。若 $I/I^2$ 是平坦 $B$-模，且
$\text{Tor}_*^C(B,B)$ 是 $I/I^2$ 上的外代数，则称这样的环同态
“弱拟正则”。若要把引理 \ref{cotangent-lemma-when-zero} 对“\'etale 环同态”
所作的事情推广到“光滑环同态”，应考察满足乘法同态
$B\otimes_AB\to B$ 弱拟正则的平坦环同态 $A\to B$。目前我们只处理
光滑环同态。

\begin{lemma}
\label{cotangent-lemma-when-projective}
若 $A\to B$ 是光滑环同态，则 $L_{B/A}=\Omega_{B/A}[0]$。
\end{lemma}

\begin{proof}
由引理 \ref{cotangent-lemma-identify-H0}，次数 $0$ 的上同调已经一致。因此只需
证明其余上同调群均为零。由引理 \ref{cotangent-lemma-localize-on-top}，对
$g\in B$ 有 $L_{B_g/A}=(L_{B/A})_g$，故只需在 $\Spec(B)$ 上局部地
证明。于是可假设 $A\to B$ 标准光滑（《交换代数》引理
\ref{algebra-lemma-smooth-syntomic}），即可以把 $A\to B$ 分解为
$A\to A[x_1,\ldots,x_n]\to B$，其中
$A[x_1,\ldots,x_n]\to B$ 是 \'etale 同态。在此情形，引理
\ref{cotangent-lemma-when-zero} 与 \ref{cotangent-lemma-localize-on-top} 给出
$L_{B/A}=L_{A[x_1,\ldots,x_n]/A}\otimes B$，再由引理
\ref{cotangent-lemma-cotangent-complex-polynomial-algebra} 得出结论。
\end{proof}





\section{正特征}
\label{cotangent-section-positive-characteristic}

\noindent
本节固定一个素数 $p$。若环 $A$ 中 $p=0$，则以 $F_A:A\to A$
记 Frobenius 自同态 $a\mapsto a^p$。

\begin{lemma}
\label{cotangent-lemma-frobenius-homotopy}
设 $A\to B$ 为环同态，且在 $A$ 中 $p=0$。令 $P_\bullet$ 为 $B$
在 $A$ 上的标准预解。由图
$$
\xymatrix{
B \ar[r]_{F_B} & B \\
A \ar[u] \ar[r]^{F_A} & A \ar[u]
}
$$
按第 \ref{cotangent-section-functoriality} 节的方式诱导的同态
$P_\bullet\to P_\bullet$，与在每个 $P_n$ 上由 Frobenius 给出的自同态
$P_\bullet\to P_\bullet$ 同伦。
\end{lemma}

\begin{proof}
令 $\mathcal{A}$ 为 $\mathbf{F}_p$-代数同态 $A\to B$ 所成的范畴。
令 $\mathcal{S}$ 为二元组 $(A,E)$ 所成的范畴，其中 $A$ 是
$\mathbf{F}_p$-代数，$E$ 是集合。考虑伴随函子
$$
V : \mathcal{A} \to \mathcal{S}, \quad (A \to B) \mapsto (A, B)
$$
以及
$$
U : \mathcal{S} \to \mathcal{A}, \quad (A, E) \mapsto (A \to A[E])
$$
令 $X$ 为《单纯形》第 \ref{simplicial-section-standard} 节所构造的、
从 $\mathcal{A}$ 到 $\mathcal{A}$ 的函子范畴中的单纯形对象。
显然 $P_\bullet=X(A\to B)$，因为若固定 $A$，则……。
% P11-E0028: frozen source contains duplicated “in” and an incomplete sentence ending “because if we fix A then.”; retained as an explicit ellipsis. See control/ERRATA_CANDIDATES.jsonl.

\medskip\noindent
令 $Y=U\circ V$。回忆 $X$ 由 $Y$ 及若干同态构造，其各项
$X_n=Y\circ\ldots\circ Y$ 含 $n+1$ 个因子；该构造见《单纯形》例
\ref{simplicial-example-godement}，细节见《单纯形》引理
\ref{simplicial-lemma-standard-simplicial} 的证明。

\medskip\noindent
令 $f:\text{id}_\mathcal{A}\to\text{id}_\mathcal{A}$ 为恒等函子的
Frobenius 自同态。换言之，令
$f_{A\to B}=(F_A,F_B):(A\to B)\to(A\to B)$。于是
$X(A\to B)$ 上的两个同态分别由自然变换
$f\star1_X$ 与 $1_X\star f$ 给出。细节略去。结论遂由《单纯形》引理
\ref{simplicial-lemma-godement-before-after} 得出。
\end{proof}

\begin{lemma}
\label{cotangent-lemma-frobenius-acts-as-zero}
设 $p$ 为素数，$A\to B$ 为环同态，并假设在 $A$ 中 $p=0$。由
Frobenius 同态 $F_A$ 与 $F_B$ 按第 \ref{cotangent-section-functoriality} 节
诱导的同态 $L_{B/A}\to L_{B/A}$ 零同伦。
\end{lemma}

\begin{proof}
令 $P_\bullet$ 为 $B$ 在 $A$ 上的标准预解。由引理
\ref{cotangent-lemma-frobenius-homotopy}，$F_A$ 与 $F_B$ 诱导的同态
$P_\bullet\to P_\bullet$，与在每一项上由 Frobenius 给出的同态
$F_{P_\bullet}:P_\bullet\to P_\bullet$ 同伦。显然
$F_{P_\bullet}$ 在 $\Omega_{P_n/A}\to\Omega_{P_n/A}$ 上诱导零同态
（因为 $p$ 次幂的导数为零），故得到所需结论。
\end{proof}

\begin{lemma}
\label{cotangent-lemma-perfect-zero}
设 $p$ 为素数，$A\to B$ 为环同态，并假设在 $A$ 中 $p=0$。若
$A$ 与 $B$ 都是完美环，则 $L_{B/A}$ 在 $D(B)$ 中为零。
\end{lemma}

\begin{proof}
同态 $(F_A,F_B):(A\to B)\to(A\to B)$ 是同构，故在 $L_{B/A}$
上诱导同构；另一方面，由引理 \ref{cotangent-lemma-frobenius-acts-as-zero}，
它在 $L_{B/A}$ 上诱导零同态。
\end{proof}





\section{与朴素余切复形的比较}
\label{cotangent-section-surjections}

\noindent
朴素余切复形已在《代数》第 \ref{algebra-section-netherlander} 节中引入。

\begin{remark}
\label{cotangent-remark-make-map}
设 $A \to B$ 为环同态。依照第 \ref{cotangent-section-compute-L-pi-shriek} 节，
在 $\mathcal{C}_{B/A}$ 上工作，并令
$\mathcal{J} \subset \mathcal{O}$ 为
$\mathcal{O} \to \underline{B}$ 的核。由引理
\ref{cotangent-lemma-apply-O-B-comparison}，注意到 $L\pi_!(\mathcal{J}) = 0$。令
$\Omega =  \Omega_{\mathcal{O}/A} \otimes_\mathcal{O} \underline{B}$
，则由引理 \ref{cotangent-lemma-compute-cotangent-complex}，
$L_{B/A} = L\pi_!(\Omega)$。于是
$L\pi_!(\mathcal{J} \to \Omega) = L\pi_!(\Omega) = L_{B/A}$。
因此，对 $\mathcal{C}_{B/A}$ 的任一对象 $U = (P \to B)$，我们得到态射
\begin{equation}
\label{cotangent-equation-comparison-map-A}
(J \to \Omega_{P/A} \otimes_P B) \longrightarrow L_{B/A}
\end{equation}
其中 $J = \Ker(P \to B)$；这是 $D(A)$ 中的态射，参见
《站点上的上同调》评注
\ref{sites-cohomology-remark-map-evaluation-to-derived}。
继续以此方式进行。由引理
\ref{cotangent-lemma-O-homology-B-homology}，注意到
$L\pi_!(\mathcal{J} \otimes_\mathcal{O}^\mathbf{L} \underline{B}) =
L\pi_!(\mathcal{J}) = 0$。由于 $\text{Tor}_0^\mathcal{O}(\mathcal{J}, \underline{B}) =
\mathcal{J}/\mathcal{J}^2$
，谱序列
$$
H_p(\mathcal{C}_{B/A}, \text{Tor}_q^\mathcal{O}(\mathcal{J}, \underline{B}))
\Rightarrow 
H_{p + q}(\mathcal{C}_{B/A},
\mathcal{J} \otimes_\mathcal{O}^\mathbf{L} \underline{B}) = 0
$$
（《导出范畴》引理 \ref{derived-lemma-two-ss-complex-functor} 的对偶）
蕴含
$H_0(\mathcal{C}_{B/A}, \mathcal{J}/\mathcal{J}^2) = 0$
且 $H_1(\mathcal{C}_{B/A}, \mathcal{J}/\mathcal{J}^2) = 0$。
由此，$\underline{B}$-模复形
$\mathcal{J}/\mathcal{J}^2 \to \Omega$ 满足
$\tau_{\geq -1}L\pi_!(\mathcal{J}/\mathcal{J}^2 \to \Omega) =
\tau_{\geq -1}L_{B/A}$。
因此，对 $\mathcal{C}_{B/A}$ 的任一对象 $U = (P \to B)$，我们得到态射
\begin{equation}
\label{cotangent-equation-comparison-map}
(J/J^2 \to \Omega_{P/A} \otimes_P B) \longrightarrow \tau_{\geq -1}L_{B/A}
\end{equation}
；这是 $D(B)$ 中的态射，参见《站点上的上同调》评注
\ref{sites-cohomology-remark-map-evaluation-to-derived}。
\end{remark}

\noindent
第一种情形是环的满射。

\begin{lemma}
\label{cotangent-lemma-surjection}
\begin{slogan}
满环同态的余切复形在零次的上同调为零；其 $-1$ 次上同调是核模去其平方。
\end{slogan}
设 $A \to B$ 为核是 $I$ 的满环同态。则
$H^0(L_{B/A}) = 0$ 且 $H^{-1}(L_{B/A}) = I/I^2$。
此同构来自 $\mathcal{C}_{B/A}$ 的对象 $(A \to B)$ 所对应的态射
(\ref{cotangent-equation-comparison-map})。
\end{lemma}

\begin{proof}
我们将在下文利用 $A \to B$ 的满射性证明：在 $\mathcal{C}_{B/A}$ 上
存在短正合层列
$$
0 \to \pi^{-1}(I/I^2) \to \mathcal{J}/\mathcal{J}^2 \to \Omega \to 0
$$
。取 $L\pi_!$ 及其相应的同调长正合列，并使用评注
\ref{cotangent-remark-make-map} 所示的
$H_1(\mathcal{C}_{B/A}, \mathcal{J}/\mathcal{J}^2)$ 与
$H_0(\mathcal{C}_{B/A}, \mathcal{J}/\mathcal{J}^2)$ 的消失，再由引理
\ref{cotangent-lemma-pi-lower-shriek-constant-sheaf} 即得所需结论。

\medskip\noindent
剩下的是验证上面提到的局部陈述。对 $\mathcal{C}_{B/A}$ 的每个对象
$U = (P \to B)$，我们可以选取同构 $P = A[E]$，使得同态
$P \to B$ 将每个 $e \in E$ 映到零。于是
$J = \mathcal{J}(U) \subset P = \mathcal{O}(U)$
等于 $J = IP + (e; e \in E)$。上述短层列在 $U$ 上的值为
$$
0 \to I/I^2 \to J/J^2 \to \Omega_{P/A} \otimes_P B \to 0
$$
验证从略（提示：唯一棘手之处是 $IP \cap J^2 = IJ$；例如可由
《代数续篇》引理 \ref{more-algebra-lemma-conormal-sequence-H1-regular} 推出）。
\end{proof}

\begin{lemma}
\label{cotangent-lemma-relation-with-naive-cotangent-complex}
设 $A \to B$ 为环同态。则 $\tau_{\geq -1}L_{B/A}$
与朴素余切复形典范拟同构。
\end{lemma}

\begin{proof}
考虑核为 $I$ 的 $P = A[B] \to B$。$B$ 在 $A$ 上的朴素余切复形
$\NL_{B/A}$ 是复形 $I/I^2 \to \Omega_{P/A} \otimes_P B$，参见
《代数》定义 \ref{algebra-definition-naive-cotangent-complex}。
注意，我们已在 (\ref{cotangent-equation-comparison-map}) 中构造了典范态射
$$
c : \NL_{B/A} \longrightarrow \tau_{\geq -1}L_{B/A}
$$
考虑与环同态 $A \to A[B] \to B$ 相关的可区别三角
(\ref{cotangent-equation-triangle})
$$
L_{P/A} \otimes_P^\mathbf{L} B \to L_{B/A} \to L_{B/P} \to 
(L_{P/A} \otimes_P^\mathbf{L} B)[1]
$$
。我们知道
$L_{P/A} = \Omega_{P/A}[0] = \NL_{P/A}$ 在 $D(P)$ 中成立
（引理 \ref{cotangent-lemma-cotangent-complex-polynomial-algebra} 与
《代数》引理 \ref{algebra-lemma-NL-polynomial-algebra}），并且
$\tau_{\geq -1}L_{B/P} = I/I^2[1] = \NL_{B/P}$ 在 $D(B)$ 中成立
（引理 \ref{cotangent-lemma-surjection} 与
《代数》引理 \ref{algebra-lemma-NL-surjection}）。由
《代数》引理 \ref{algebra-lemma-exact-sequence-NL}
及与该可区别三角相关的上同调长正合列，要证明 $c$ 是拟同构，
只需证明态射 $L_{P/A} \to L_{B/A} \to L_{B/P}$ 在上同调群上
与朴素余切复形的相应态射
$\NL_{P/A} \to \NL_{B/A} \to \NL_{B/P}$
相容。验证从略。
\end{proof}

\begin{remark}
\label{cotangent-remark-explicit-comparison-map}
我们可以如下明确写出引理
\ref{cotangent-lemma-relation-with-naive-cotangent-complex} 的比较态射。
令 $P_\bullet$ 为 $B$ 在 $A$ 上的标准预解，并令
$I = \Ker(A[B] \to B)$。回忆 $P_0 = A[B]$。该引理中的态射
由交换图
$$
\xymatrix{
L_{B/A} \ar[d] & \ldots \ar[r] &
\Omega_{P_2/A} \otimes_{P_2} B
\ar[r] \ar[d] &
\Omega_{P_1/A} \otimes_{P_1} B
\ar[r] \ar[d] &
\Omega_{P_0/A} \otimes_{P_0} B
\ar[d] \\
\NL_{B/A} & \ldots \ar[r] &
0 \ar[r] & 
I/I^2 \ar[r] &
\Omega_{P_0/A} \otimes_{P_0} B
}
$$
给出。我们通过将 $\text{d}f \otimes b$ 映到
$(d_0(f) - d_1(f))b$ 在 $I/I^2$ 中的类，构造以 $I/I^2$ 为靶的向下箭头。
这里 $d_i : P_1 \to P_0$（$i = 0, 1$）是单纯形结构的两个面态射。
这是有意义的，因为 $d_0 - d_1$ 将 $P_1$ 映入
$I = \Ker(P_0 \to B)$。此规则良定义的验证从略。
我们的态射与微分
$\Omega_{P_1/A} \otimes_{P_1} B \to \Omega_{P_0/A} \otimes_{P_0} B$
相容，因为该微分将 $\text{d}f \otimes b$ 映到
$\text{d}(d_0(f) - d_1(f)) \otimes b$。此外，微分
$\Omega_{P_2/A} \otimes_{P_2} B \to \Omega_{P_1/A} \otimes_{P_1} B$
将 $\text{d}f \otimes b$ 映到
$\text{d}(d_0(f) - d_1(f) + d_2(f)) \otimes b$，后者被我们的向下箭头
消去。因此得到了复形态射。它与引理
\ref{cotangent-lemma-relation-with-naive-cotangent-complex} 中的态射相同这一验证从略。
\end{remark}

\begin{remark}
\label{cotangent-remark-surjection}
采用评注 \ref{cotangent-remark-make-map} 的记号。那里给出的论证表明，谱序列的微分
$$
H_2(\mathcal{C}_{B/A}, \mathcal{J}/\mathcal{J}^2)
\longrightarrow
H_0(\mathcal{C}_{B/A}, \text{Tor}_1^\mathcal{O}(\mathcal{J}, \underline{B}))
$$
是同构。令 $\mathcal{C}'_{B/A}$ 表示 $\mathcal{C}_{B/A}$ 中由满射
$P \to B$ 构成的全子范畴。余切复形与朴素余切复形相一致（引理
\ref{cotangent-lemma-relation-with-naive-cotangent-complex}），这表明我们在
$\mathcal{C}'_{B/A}$ 上有正合层列
$$
0 \to \underline{H_1(L_{B/A})} \to
\mathcal{J}/\mathcal{J}^2 \xrightarrow{\text{d}} \Omega \to
\underline{H_2(L_{B/A})} \to 0
$$
。由此，整个范畴 $\mathcal{C}_{B/A}$ 上的 $\Ker(d)$ 与 $\Coker(d)$
具有消失的高阶同调群，因为由引理 \ref{cotangent-lemma-identify-pi-shriek}，
这些群由常值单纯形阿贝尔群的同调群计算。因此我们推出
$$
H_n(\mathcal{C}_{B/A}, \mathcal{J}/\mathcal{J}^2) \to H_n(L_{B/A})
$$
对一切 $n \geq 2$ 都是同构。结合上面的评注，我们得到公式
$H_2(L_{B/A}) =
H_0(\mathcal{C}_{B/A}, \text{Tor}_1^\mathcal{O}(\mathcal{J}, \underline{B}))$。
\end{remark}





\section{Quillen 谱序列}
\label{cotangent-section-spectral-sequence}

\noindent
本节讨论一个联系导出张量积与余切复形的谱序列。

\begin{lemma}
\label{cotangent-lemma-vanishing-symmetric-powers}
记号和假设同《站点上的上同调》例
\ref{sites-cohomology-example-category-to-point}。假设 $\mathcal{C}$
具有《站点上的上同调》引理
\ref{sites-cohomology-lemma-compute-by-cosimplicial-resolution}
中的余单纯形对象。设 $\mathcal{F}$ 为满足
$H_0(\mathcal{C}, \mathcal{F}) = 0$ 的平坦 $\underline{B}$-模。
则当 $l < k$ 时，
$H_l(\mathcal{C}, \text{Sym}_{\underline{B}}^k(\mathcal{F})) = 0$。
\end{lemma}

\begin{proof}
在张量积、外幂和对称幂中省略下标 ${}_{\underline{B}}$。
我们对 $k$ 归纳证明该引理。$k = 0, 1$ 的情形由假设即得。
若 $k > 1$，考虑正合复形
$$
\ldots \to
\wedge^2\mathcal{F} \otimes \text{Sym}^{k - 2}\mathcal{F} \to
\mathcal{F} \otimes \text{Sym}^{k - 1}\mathcal{F} \to
\text{Sym}^k\mathcal{F} \to 0
$$
，其微分如 Koszul 复形中所定义。若将其视为
$\text{Sym}^k\mathcal{F}$ 的预解，便得到第一象限谱序列
$$
E_1^{p, q} =
H_p(\mathcal{C},
\wedge^{q + 1}\mathcal{F} \otimes \text{Sym}^{k - q - 1}\mathcal{F})
\Rightarrow
H_{p + q}(\mathcal{C}, \text{Sym}^k(\mathcal{F}))
$$
由《站点上的上同调》引理 \ref{sites-cohomology-lemma-eilenberg-zilber}，有
$$
L\pi_!(\wedge^{q + 1}\mathcal{F} \otimes \text{Sym}^{k - q - 1}\mathcal{F}) =
L\pi_!(\wedge^{q + 1}\mathcal{F}) \otimes_B^\mathbf{L}
L\pi_!(\text{Sym}^{k - q - 1}\mathcal{F}))
$$
% P11-E0029：冻结源在上一公式末尾保留一个疑似多余的右圆括号。
由导出张量积的构造可知，将归纳假设与
$H_0(\mathcal{C}, \wedge^{q + 1}(\mathcal{F})) = 0$ 的消失结合起来，
便可证明所需结论。该消失成立，是因为
$\wedge^{q + 1}(\mathcal{F})$ 是 $\mathcal{F}^{\otimes q + 1}$ 的商，且
$H_0(\mathcal{C}, \mathcal{F}^{\otimes q + 1})$ 是
$H_0(\mathcal{C}, \mathcal{F})^{\otimes q + 1}$ 的商，而后者为零。
\end{proof}

\begin{remark}
\label{cotangent-remark-first-homology-symmetric-power}
在引理 \ref{cotangent-lemma-vanishing-symmetric-powers} 的情形下，可以证明
$H_k(\mathcal{C}, \text{Sym}^k(\mathcal{F})) =
\wedge^k_B(H_1(\mathcal{C}, \mathcal{F}))$。事实上，从证明可以推出，
$H_k(\mathcal{C}, \text{Sym}^k(\mathcal{F}))$ 是下式的 $S_k$-余不变量：
$$
H^{-k}(L\pi_!(\mathcal{F}) \otimes_B^\mathbf{L}
L\pi_!(\mathcal{F}) \otimes_B^\mathbf{L}
\ldots \otimes_B^\mathbf{L} L\pi_!(\mathcal{F})) =
H_1(\mathcal{C}, \mathcal{F})^{\otimes k}
$$
因此，我们的断言是：这一作用由 $S_k$ 在张量积上的通常作用乘以符号特征给出。
要证明这一点，必须逐一核查与多重复形相关的全复形定义中的符号约定。
验证从略。
\end{remark}

\begin{lemma}
\label{cotangent-lemma-map-tors-zero}
设 $A$ 为环，$P = A[E]$ 为多项式环。令
$I = (e; e \in E) \subset P$。则态射
$\text{Tor}_i^P(A, I^{n + 1}) \to \text{Tor}_i^P(A, I^n)$
对一切 $i$ 和 $n$ 都为零。
\end{lemma}

\begin{proof}
以 $x_e \in P$ 表示对应于 $e \in E$ 的变量。$A$ 在 $P$ 上的一个自由预解
由 $x_e$ 上的 Koszul 复形 $K_\bullet$ 给出。这里 $K_i$ 的基由外积
$e_1 \wedge \ldots \wedge e_i$（$e_1, \ldots, e_i \in E$）给出，且
$d(e) = x_e$。因此 $K_\bullet \otimes_P I^n = I^nK_\bullet$ 计算
$\text{Tor}_i^P(A, I^n)$。注意一切均分次，其中
$\deg(x_e) = 1$、$\deg(e) = 1$，而对 $a \in A$ 有 $\deg(a) = 0$。
设 $\xi \in I^{n + 1}K_i$ 是次数为 $m$ 的齐次余圈。注意
$m \geq i + 1 + n$。由于 $K_\bullet$ 在次数 $>0$ 正合，存在
$\eta \in K_{i + 1}$ 使 $\xi = \text{d}\eta$（$i = 0$ 的情形留给读者）。
现在 $\deg(\eta) = m \geq i + 1 + n$。因而用上述基展开 $\eta$，
可见各坐标都属于 $I^n$。所以 $\xi$ 在 $I^nK_\bullet$ 的同调中映为零，
正如所需。
\end{proof}

\begin{theorem}[Quillen 谱序列]
\label{cotangent-theorem-quillen-spectral-sequence}
设 $A \to B$ 为满环同态。考虑 $\mathcal{C}_{B/A}$ 上的
$\underline{B}$-模层
$\Omega = \Omega_{\mathcal{O}/A} \otimes_\mathcal{O} \underline{B}$，
参见第 \ref{cotangent-section-compute-L-pi-shriek} 节。则存在 $E_1$ 页为
$$
E_1^{p, q} =
H_{- p - q}(\mathcal{C}_{B/A}, \text{Sym}^p_{\underline{B}}(\Omega))
\Rightarrow \text{Tor}^A_{- p - q}(B, B)
$$
的谱序列，其中 $d_r$ 的双次数为 $(r, -r + 1)$。此外，当 $i < k$ 时，
$H_i(\mathcal{C}_{B/A}, \text{Sym}^k_{\underline{B}}(\Omega)) = 0$。
\end{theorem}

\begin{proof}
令 $I \subset A$ 为 $A \to B$ 的核，并令
$\mathcal{J} \subset \mathcal{O}$ 为
$\mathcal{O} \to \underline{B}$ 的核。于是
$I\mathcal{O} \subset \mathcal{J}$。令
$\mathcal{K} = \mathcal{J}/I\mathcal{O}$ 且
$\overline{\mathcal{O}} = \mathcal{O}/I\mathcal{O}$。

\medskip\noindent
对 $\mathcal{C}_{B/A}$ 的每个对象 $U = (P \to B)$，我们可以选取同构
$P = A[E]$，使同态 $P \to B$ 将每个 $e \in E$ 映到零。于是
$J = \mathcal{J}(U) \subset P = \mathcal{O}(U)$
等于 $J = IP + (e; e \in E)$。此外，
$\overline{\mathcal{O}}(U) = B[E]$ 且 $K = \mathcal{K}(U) = (e; e \in E)$
，而且 $K$ 是多项式环 $B[E]$ 中由这些变量生成的理想。特别地，显然
$$
K/K^2 \xrightarrow{\text{d}} \Omega_{P/A} \otimes_P B
$$
是双射。换言之，
$\Omega = \mathcal{K}/\mathcal{K}^2$ 且
$\text{Sym}_B^k(\Omega) = \mathcal{K}^k/\mathcal{K}^{k + 1}$。
注意到 $\pi_!(\Omega) = \Omega_{B/A} = 0$（引理
\ref{cotangent-lemma-identify-H0}），因为 $A \to B$ 是满射（《代数》引理
\ref{algebra-lemma-trivial-differential-surjective}）。
由引理 \ref{cotangent-lemma-vanishing-symmetric-powers}，我们推出
$$
H_i(\mathcal{C}_{B/A}, \mathcal{K}^k/\mathcal{K}^{k + 1}) =
H_i(\mathcal{C}_{B/A}, \text{Sym}^k_{\underline{B}}(\Omega)) = 0
$$
对 $i < k$ 成立。这证明了定理的最后一个陈述。

\medskip\noindent
证明该定理的方法是注意到
$$
B \otimes_A^\mathbf{L} B = L\pi_!(\mathcal{O}) \otimes_A^\mathbf{L} B =
L\pi_!(\mathcal{O} \otimes_{\underline{A}}^\mathbf{L} \underline{B}) =
L\pi_!(\overline{\mathcal{O}})
$$
第一个等号由引理 \ref{cotangent-lemma-apply-O-B-comparison} 得到，第二个等号由
《站点上的上同调》引理 \ref{sites-cohomology-lemma-change-of-rings} 得到，
第三个等号则因为 $\mathcal{O}$ 在 $\underline{A}$ 上平坦。层
$\overline{\mathcal{O}}$ 有滤过
$$
\ldots \subset
\mathcal{K}^3 \subset
\mathcal{K}^2 \subset
\mathcal{K} \subset
\overline{\mathcal{O}}
$$
。这在表示 $L\pi_!(\overline{\mathcal{O}})$ 的复形 $C$ 上诱导滤过
$F$，其中 $F^pC$ 表示 $L\pi_!(\mathcal{K}^p)$（$C$ 与 $F$ 的构造从略）。
考虑《同调》第 \ref{homology-section-filtered-complex} 节中与
$(C, F)$ 相关的谱序列。其 $E_1$ 页为
$$
E_1^{p, q} = H_{- p - q}(\mathcal{C}_{B/A}, \mathcal{K}^p/\mathcal{K}^{p + 1})
\quad\Rightarrow\quad
H_{- p - q}(\mathcal{C}_{B/A}, \overline{\mathcal{O}}) = 
\text{Tor}_{- p - q}^A(B, B)
$$
，微分为 $E_r^{p, q} \to E_r^{p + r, q - r + 1}$。为了证明收敛性，
我们将证明：对每个 $k$，存在 $c$ 使得
$H_i(\mathcal{C}_{B/A}, \mathcal{K}^n) = 0$
对 $i < k$ 与 $n > c$ 成立\footnote{事后可以推出“正确的”消失：
当 $i < n$ 时，$H_i(\mathcal{C}_{B/A}, \mathcal{K}^n) = 0$。}。

\medskip\noindent
给定 $k \geq 0$，令 $c = k^2$。我们断言
$$
H_i(\mathcal{C}_{B/A}, \mathcal{K}^{n + c}) \to
H_i(\mathcal{C}_{B/A}, \mathcal{K}^n)
$$
对 $i < k$ 及一切 $n \geq 0$ 都是零。注意
$\mathcal{K}^n/\mathcal{K}^{n + c}$ 有有限滤过，其相继商为
$\mathcal{K}^m/\mathcal{K}^{m + 1}$（$n \leq m < n + c$），且当
$i < n$ 时，
$H_i(\mathcal{C}_{B/A}, \mathcal{K}^m/\mathcal{K}^{m + 1}) = 0$
（见上文）。因此，该断言蕴含：对 $i < k$ 及一切 $n \geq k$，
$H_i(\mathcal{C}_{B/A}, \mathcal{K}^{n + c}) = 0$，这正是我们需要证明的。

\medskip\noindent
证明该断言。回忆，对任一 $\mathcal{O}$-模 $\mathcal{F}$，态射
$\mathcal{F} \to \mathcal{F} \otimes_\mathcal{O}^\mathbf{L} B$
在施用 $L\pi_!$ 后诱导同构，参见引理
\ref{cotangent-lemma-O-homology-B-homology}。考虑态射
$$
\mathcal{K}^{n + k} \otimes_\mathcal{O}^\mathbf{L} B
\longrightarrow
\mathcal{K}^n \otimes_\mathcal{O}^\mathbf{L} B
$$
我们断言，此态射在次数 $0, -1, \ldots, - k + 1$ 的上同调层上诱导零态射。
若这一第二断言成立，则 $k$ 重复合
$$
\mathcal{K}^{n + c} \otimes_\mathcal{O}^\mathbf{L} B
\longrightarrow
\mathcal{K}^n \otimes_\mathcal{O}^\mathbf{L} B
$$
分解经过
$\tau_{\leq -k}\mathcal{K}^n \otimes_\mathcal{O}^\mathbf{L} B$，从而当
$i < k$ 时在 $H_i(\mathcal{C}_{B/A}, -) = L_i\pi_!( - )$ 上诱导零，
参见《导出范畴》引理
\ref{derived-lemma-trick-vanishing-composition}。由上面的评注，这意味着
$H_i(\mathcal{C}_{B/A}, \mathcal{K}^{n + c}) \to
H_i(\mathcal{C}_{B/A}, \mathcal{K}^n)$
同样为零，因而证明了第一个断言。

\medskip\noindent
证明第二个断言。该陈述是局部的，故可在上述对象
$U = (P \to B)$ 上工作。我们必须证明态射
$$
\text{Tor}_i^P(B, K^{n + k}) \to \text{Tor}_i^P(B, K^n)
$$
当 $i < k$ 时为零。存在谱序列
$$
\text{Tor}_a^P(P/IP, \text{Tor}_b^{P/IP}(B, K^n))
\Rightarrow
\text{Tor}_{a + b}^P(B, K^n),
$$
，参见《代数续篇》例 \ref{more-algebra-example-tor-change-rings}。
因此，只需证明态射
$$
\text{Tor}_i^{P/IP}(B, K^{n + 1}) \to \text{Tor}_i^{P/IP}(B, K^n)
$$
对一切 $i$ 都为零。这正是引理 \ref{cotangent-lemma-map-tors-zero}。
\end{proof}

\begin{remark}
\label{cotangent-remark-elucidate-ss}
在定理 \ref{cotangent-theorem-quillen-spectral-sequence} 的情形下，令
$I = \Ker(A \to B)$。则
$H^{-1}(L_{B/A}) = H_1(\mathcal{C}_{B/A}, \Omega) = I/I^2$，参见
引理 \ref{cotangent-lemma-surjection}。因此由评注
\ref{cotangent-remark-first-homology-symmetric-power}，
$H_k(\mathcal{C}_{B/A}, \text{Sym}^k(\Omega)) = \wedge^k_B(I/I^2)$。
于是 $E_1$ 页形如
$$
\begin{matrix}
B \\
0 \\
0 & I/I^2 \\
0 & H^{-2}(L_{B/A}) \\
0 & H^{-3}(L_{B/A}) & \wedge^2(I/I^2) \\
0 & H^{-4}(L_{B/A}) & H_3(\mathcal{C}_{B/A}, \text{Sym}^2(\Omega)) \\
0 & H^{-5}(L_{B/A}) & H_4(\mathcal{C}_{B/A}, \text{Sym}^2(\Omega)) &
\wedge^3(I/I^2)
\end{matrix}
$$
，其微分沿水平方向。因此得到边缘态射
$\text{Tor}_i^A(B, B) \to H^{-i}(L_{B/A})$（$i > 0$）及
$\wedge^i_B(I/I^2) \to \text{Tor}_i^A(B, B)$。最后，
$\text{Tor}_1^A(B, B) = I/I^2$，并且有五项正合列
$$
\text{Tor}_3^A(B, B) \to H^{-3}(L_{B/A}) \to \wedge^2_B(I/I^2) \to
\text{Tor}_2^A(B, B) \to H^{-2}(L_{B/A}) \to 0
$$
，它由低次数项组成。
\end{remark}

\begin{remark}
\label{cotangent-remark-elucidate-degree-two}
设 $A \to B$ 为环同态，$P_\bullet$ 为 $B$ 在 $A$ 上的预解
（评注 \ref{cotangent-remark-resolution}）。令
$J_n = \Ker(P_n \to B)$。注意到
$$
\text{Tor}_2^{P_n}(B, B) = 
\text{Tor}_1^{P_n}(J_n, B) =
\Ker(J_n \otimes_{P_n} J_n \to J_n^2).
$$
因此，由评注 \ref{cotangent-remark-surjection}，$H_2(L_{B/A})$ 典范地等于
$$
\Coker(\text{Tor}_2^{P_1}(B, B) \to \text{Tor}_2^{P_0}(B, B))
$$
。为使其更为明确，我们依照例 \ref{cotangent-example-resolution-length-two}
选取 $P_2$、$P_1$、$P_0$。我们断言
$$
\text{Tor}_2^{P_1}(B, B) =
\wedge^2(\bigoplus\nolimits_{t \in T} B)\ \oplus
\ \bigoplus\nolimits_{t \in T} J_0\ \oplus
\ \text{Tor}_2^{P_0}(B, B)
$$
确切地说，第一直和项的基元素 $x_t \wedge x_{t'}$ 对应于
$J_1 \otimes_{P_1} J_1$ 的元素
$x_t \otimes x_{t'} - x_{t'} \otimes x_t$。对 $f \in J_0$，
第二直和项的元素 $x_t \otimes f$ 对应于
$J_1 \otimes_{P_1} J_1$ 的元素
$x_t \otimes s_0(f) - s_0(f) \otimes x_t$。最后，态射
$\text{Tor}_2^{P_0}(B, B) \to \text{Tor}_2^{P_1}(B, B)$ 由 $s_0$ 给出。
态射
$d_0 - d_1 : \text{Tor}_2^{P_1}(B, B) \to \text{Tor}_2^{P_0}(B, B)$
在最后一个直和项上为零，将 $x_t \otimes f$ 映到
$f \otimes f_t - f_t \otimes f$，并将 $x_t \wedge x_{t'}$ 映到
$f_t \otimes f_{t'} - f_{t'} \otimes f_t$。综上，我们推出有正合列
$$
\wedge^2_B(J_0/J_0^2) \to \text{Tor}_2^{P_0}(B, B) \to H^{-2}(L_{B/A}) \to 0
$$
这样，我们便直接证明了评注 \ref{cotangent-remark-elucidate-ss}
中所讨论的 Quillen 谱序列的一个推论。
\end{remark}






\section{与 Lichtenbaum--Schlessinger 复形的比较}
\label{cotangent-section-compare-higher}

\noindent
设 $A \to B$ 为环同态。文献 \cite{Lichtenbaum-Schlessinger}
对 $\tau_{\geq -2}L_{B/A}$ 给出了相当明确的刻画；它常用于计算奇点的
半普遍形变空间。其构造如下。选取 $A$ 上的多项式代数 $P$ 以及核为
$I$ 的满射 $P \to B$。选取 $I$ 的生成元 $f_t$（$t \in T$），
从而诱导满射 $F = \bigoplus_{t \in T} P \to I$，其中 $F$ 是自由
$P$-模。令 $Rel \subset F$ 为 $F \to I$ 的核；换言之，$Rel$ 是
$f_t$ 之间的关系之集。令 $TrivRel \subset Rel$ 为平凡关系的子模，
即 $Rel$ 中由元素
$(\ldots, f_{t'}, 0, \ldots, 0, -f_t, 0, \ldots)$ 生成的子模。
考虑 $B$-模复形
\begin{equation}
\label{cotangent-equation-lichtenbaum-schlessinger}
Rel/TrivRel \longrightarrow
F \otimes_P B \longrightarrow
\Omega_{P/A} \otimes_P B
\end{equation}
，其中最后一项置于次数 $0$。第一个态射是显然的，第二个态射将对应于
$t \in T$ 的基元素映到 $\text{d}f_t \otimes 1$。

\begin{definition}
\label{cotangent-definition-biderivation}
设 $A \to B$ 为环同态，$M$ 为 $A$ 上的 $(B, B)$-双模。
所谓 {\it $A$-双导子}，是指满足
$\lambda(xy) = x\lambda(y) + \lambda(x)y$ 的 $A$-线性映射
$\lambda : B \to M$。
\end{definition}

\noindent
多项式代数上的双导子很容易描述。

\begin{lemma}
\label{cotangent-lemma-polynomial-ring-unique}
设 $P = A[S]$ 为 $A$ 上的多项式环，$M$ 为 $A$ 上的 $(P, P)$-双模。
给定每个 $s \in S$ 所对应的 $m_s \in M$，存在唯一的 $A$-双导子
$\lambda : P \to M$，它将 $s$ 映到 $m_s$。
\end{lemma}

\begin{proof}
我们在 $M$ 中令
$$
\lambda(s_1 \ldots s_t) =
\sum s_1 \ldots s_{i - 1} m_{s_i} s_{i + 1} \ldots s_t
$$
。按 $A$-线性延拓便得到一个双导子。
\end{proof}

\noindent
下面是比较陈述。读者也可参阅
\cite[page 206, Proposition 12]{Andre-Homologie}，或参阅文献
\cite{Doncel}；后者将复形 (\ref{cotangent-equation-lichtenbaum-schlessinger})
扩充一项，并将比较推广到 $\tau_{\geq -3}$。

\begin{lemma}
\label{cotangent-lemma-compare-higher}
在上述情形中，以 $L$ 表示复形
(\ref{cotangent-equation-lichtenbaum-schlessinger})。在 $D(B)$ 中存在典范态射
$L_{B/A} \to L$，它在 $D(B)$ 中诱导同构
$\tau_{\geq -2}L_{B/A} \to L$。
\end{lemma}

\begin{proof}
令 $P_\bullet \to B$ 为 $B$ 在 $A$ 上的预解（评注
\ref{cotangent-remark-resolution}）。我们将 $L_{B/A}$ 与
$\Omega_{P_\bullet/A} \otimes B$ 等同。为了构造该态射，需要作出若干选择。

\medskip\noindent
选取与给定态射 $P_0 \to B$ 和 $P \to B$ 相容的 $A$-代数同态
$\psi : P_0 \to P$。

\medskip\noindent
对某个集合 $S$ 写成 $P_1 = A[S]$。对 $s \in S$，可写成
$$
\psi(d_0(s) - d_1(s)) = \sum p_{s, t} f_t
$$
，其中 $p_{s, t} \in P$。通过态射
$(\psi \circ d_0, \psi \circ d_1)$，将
$F = \bigoplus_{t \in T} P$ 视为 $(P_1, P_1)$-双模。
由引理 \ref{cotangent-lemma-polynomial-ring-unique}，得到唯一的 $A$-双导子
$\lambda : P_1 \to F$，它将 $s$ 映到坐标为 $p_{s,t}$ 的向量。
按构造，复合
$$
P_1 \longrightarrow F \longrightarrow P
$$
将 $f \in P_1$ 映到 $\psi(d_0(f) - d_1(f))$，因为映射
$f \mapsto \psi(d_0(f) - d_1(f))$ 是一个 $A$-双导子，且在生成元上
与该复合一致。

\medskip\noindent
对 $g \in P_2$，我们断言
$\lambda(d_0(g) - d_1(g) + d_2(g))$ 是 $Rel$ 的元素。事实上，由上一段
最后的说明，它在 $P$ 中的像为
$$
\psi((d_0 - d_1)(d_0(g) - d_1(g) + d_2(g)))
$$
；由《单纯形》第 \ref{simplicial-section-complexes} 节，它为零）。
% P11-E0030：冻结源在上一句末尾保留一个疑似无匹配的右圆括号。

\medskip\noindent
$\psi$ 的选择确定态射
$$
\text{d}\psi \otimes 1 :
\Omega_{P_0/A} \otimes B
\longrightarrow
\Omega_{P/A} \otimes B
$$
将 $\lambda$ 与态射 $F \to F \otimes B$ 复合，得到通常的
$A$-导子，因为 $F \otimes B$ 上的两个 $P_1$-模结构相同。因此
$\lambda$ 确定态射
$$
\overline{\lambda} :
\Omega_{P_1/A} \otimes B
\longrightarrow
F \otimes B
$$
最后，我们得到 $B$-线性态射
$$
q :
\Omega_{P_2/A} \otimes B
\longrightarrow
Rel/TrivRel
$$
，它将 $\text{d}g$ 映到
$\lambda(d_0(g) - d_1(g) + d_2(g))$ 在该商中的类。

\medskip\noindent
图
$$
\xymatrix{
\Omega_{P_3/A} \otimes B \ar[r] \ar[d] &
\Omega_{P_2/A} \otimes B \ar[r] \ar[d]_q &
\Omega_{P_1/A} \otimes B \ar[r] \ar[d]_{\overline{\lambda}} &
\Omega_{P_0/A} \otimes B \ar[d]_{\text{d}\psi \otimes 1} \\
0 \ar[r] &
Rel/TrivRel \ar[r] &
F \otimes B \ar[r] &
\Omega_{P/A} \otimes B
}
$$
交换（计算从略），从而得到该引理的态射。由评注
\ref{cotangent-remark-explicit-comparison-map} 和引理
\ref{cotangent-lemma-relation-with-naive-cotangent-complex}，可见此态射诱导同构
$H_1(L_{B/A}) \to H_1(L)$ 与 $H_0(L_{B/A}) \to H_0(L)$。

\medskip\noindent
还须证明态射 $L_{B/A} \to L$ 诱导同构
$H_2(L_{B/A}) \to H_2(L)$。选取 $B$ 在 $A$ 上的预解，其中
$P_0 = P = A[u_i]$，再依照例 \ref{cotangent-example-resolution-length-two}
选取 $P_1$ 和 $P_2$。在评注 \ref{cotangent-remark-elucidate-degree-two}
中，我们构造了正合列
$$
\wedge^2_B(J_0/J_0^2) \to \text{Tor}_2^{P_0}(B, B) \to H^{-2}(L_{B/A}) \to 0
$$
，其中 $P_0 = P$ 且 $J_0 = \Ker(P \to B) = I$。利用短正合列
$0 \to I \to P \to B \to 0$ 与
$0 \to Rel \to F \to I \to 0$ 计算 Tor 群，得到
$\text{Tor}_2^P(B, B) = \Ker(Rel \otimes B \to F \otimes B)$。
在这一等同下，态射
$\wedge^2_B(I/I^2) \to \text{Tor}_2^P(B, B)$ 的像恰好是
$TrivRel \otimes B$ 的像。因此 $H_2(L_{B/A}) \cong H_2(L)$。

\medskip\noindent
最后还须核查，我们的态射 $L_{B/A} \to L$ 确实诱导上述同构。
以下将直接使用例 \ref{cotangent-example-resolution-length-two} 及评注
\ref{cotangent-remark-elucidate-degree-two}、\ref{cotangent-remark-surjection}
中讨论的记号和结果，不再另行说明。在
$\text{Tor}_2^{P_0}(B, B) = \Ker(I \otimes_P I \to I^2)$
中选取元素 $\xi$。对某些 $h_{t',t} \in P$，写成
$\xi = \sum h_{t', t}f_{t'} \otimes f_t$。沿上述正合列追踪可知，
$\xi$ 对应于元素
$r \in Rel \subset F = \bigoplus_{t \in T} P$ 在 $Rel \otimes B$
中的像，其中第 $t$ 个坐标为
$r_t = \sum_{t' \in T} h_{t', t}f_{t'}$。另一方面，$\xi$ 对应于
$H_2(L_{B/A}) = H_2(\Omega)$ 的如下元素：在二重扩张
$$
0 \to
\text{Tor}_2^\mathcal{O}(\underline{B}, \underline{B})
\to 
\mathcal{J} \otimes_\mathcal{O} \mathcal{J} \to \mathcal{J}
\to
\mathcal{J}/\mathcal{J}^2 \to 0
$$
下取 $\xi$ 的边界，再经
$\text{d} : H_2(\mathcal{J}/\mathcal{J}^2) \to H_2(\Omega)$
所得的像。我们计算该元素的相继超越映射。首先有
$$
\xi = (d_0 - d_1)(- \sum s_0(h_{t', t} f_{t'}) \otimes x_t)
$$
，接着有
$$
\sum s_0(h_{t', t} f_{t'}) x_t = d_0(v_r) - d_1(v_r) + d_2(v_r)
$$
，这是由于我们在例 \ref{cotangent-example-resolution-length-two} 中对变量 $v$
所作的选择。可以选取上面的态射 $\lambda$，使
$\lambda(u_i) = 0$ 且 $\lambda(x_t) = - e_t$，其中 $e_t \in F$
表示对应于 $t \in T$ 的基向量。因此，上述态射 $q$ 的构造将
$\text{d}v_r$ 映到
$$
\lambda(\sum s_0(h_{t', t} f_{t'}) x_t) =
\sum\nolimits_t \left(\sum\nolimits_{t'} h_{t', t}f_{t'}\right) e_t
$$
，这与 $\xi$ 在 $Rel \otimes B$ 中的像相符（上面出现的两个负号相消）。
这一相符性完成了证明。
\end{proof}

\begin{remark}[Lichtenbaum--Schlessinger 复形的函子性]
\label{cotangent-remark-functoriality-lichtenbaum-schlessinger}
考虑环同态的交换方块
$$
\xymatrix{
A' \ar[r] & B' \\
A \ar[u] \ar[r] & B \ar[u]
}
$$
。选取分解
$$
\xymatrix{
A' \ar[r] & P' \ar[r] & B' \\
A \ar[u] \ar[r] & P \ar[u] \ar[r] & B \ar[u]
}
$$
，其中 $P$ 是 $A$ 上的多项式代数，$P'$ 是 $A'$ 上的多项式代数。
选取 $\Ker(P \to B)$ 的生成元 $f_t$（$t \in T$）。
对 $t \in T$，以 $f'_t$ 表示 $f_t$ 在 $P'$ 中的像。选取
$f'_s \in P'$，使得当 $t \in T' = T \amalg S$ 时，元素 $f'_t$
生成 $P' \to B'$ 的核。令
$F = \bigoplus_{t \in T} P$ 且
$F' = \bigoplus_{t' \in T'} P'$。令
$Rel = \Ker(F \to P)$ 与 $Rel' = \Ker(F' \to P')$，其中各坐标上的
态射分别由乘以 $f_t$ 和 $f'_t$ 给出。最后，分别令 $TrivRel$ 与
$TrivRel'$ 等于 $Rel$ 与 $TrivRel$ 中由元素
% P11-E0031：冻结源把第二个承载子模写成 TrivRel，而不是 Rel'。
$(\ldots, f_{t'}, 0, \ldots, 0, -f_t, 0, \ldots)$
生成的子模，其中 $t,t'$ 分别取自 $T$ 与 $T'$。作出这些选择后，得到典范交换图
$$
\xymatrix{
L' : &
Rel'/TrivRel' \ar[r] &
F' \otimes_{P'} B' \ar[r] &
\Omega_{P'/A'} \otimes_{P'} B' \\
L : \ar[u] &
Rel/TrivRel \ar[r] \ar[u] &
F \otimes_P B \ar[r] \ar[u] &
\Omega_{P/A} \otimes_P B \ar[u]
}
$$
此外，追踪引理 \ref{cotangent-lemma-compare-higher} 的证明中所作的选择，
读者可见会得到交换图
$$
\xymatrix{
L_{B'/A'} \ar[r] & L' \\
L_{B/A} \ar[r] \ar[u] & L \ar[u]
}
$$
\end{remark}





\section{局部完全交的余切复形}
\label{cotangent-section-lci}

\noindent
若 $A \to B$ 是局部完全交同态，则 $L_{B/A}$ 是完美复形。
证明这一点的关键是下面的引理。

\begin{lemma}
\label{cotangent-lemma-special-case}
设 $A = \mathbf{Z}[x_1, \ldots, x_n] \to B = \mathbf{Z}$ 为将
$x_i$ 映到 $0$（$i = 1, \ldots, n$）的环同态。令
$I = (x_1, \ldots, x_n) \subset A$。则 $L_{B/A}$ 与 $I/I^2[1]$
拟同构。
\end{lemma}

\begin{proof}
有若干种证明方法。例如，可以明确构造 $B$ 在 $A$ 上的预解并作计算。
我们将使用 (\ref{cotangent-equation-triangle})。确切地说，考虑可区别三角
$$
L_{\mathbf{Z}[x_1, \ldots, x_n]/\mathbf{Z}}
\otimes_{\mathbf{Z}[x_1, \ldots, x_n]} \mathbf{Z} \to
L_{\mathbf{Z}/\mathbf{Z}} \to
L_{\mathbf{Z}/\mathbf{Z}[x_1, \ldots, x_n]}\to
L_{\mathbf{Z}[x_1, \ldots, x_n]/\mathbf{Z}}
\otimes_{\mathbf{Z}[x_1, \ldots, x_n]} \mathbf{Z}[1]
$$
由引理 \ref{cotangent-lemma-cotangent-complex-polynomial-algebra}，复形
$L_{\mathbf{Z}[x_1, \ldots, x_n]/\mathbf{Z}}$ 与
$\Omega_{\mathbf{Z}[x_1, \ldots, x_n]/\mathbf{Z}}$ 拟同构。
由引理 \ref{cotangent-lemma-when-zero}，复形 $L_{\mathbf{Z}/\mathbf{Z}}$
在 $D(\mathbf{Z})$ 中为零。因此，$L_{B/A}$ 只有一个非零上同调群；
由引理 \ref{cotangent-lemma-surjection}，它正如引理中所描述。
\end{proof}

\begin{lemma}
\label{cotangent-lemma-mod-regular-sequence}
设 $A \to B$ 为满环同态，其核 $I$ 由 Koszul 正则序列（例如正则序列）
生成。则 $L_{B/A}$ 与 $I/I^2[1]$ 拟同构。
\end{lemma}

\begin{proof}
设 $f_1, \ldots, f_r \in I$ 是生成 $I$ 的 Koszul 正则序列。
考虑将 $x_i$ 映到 $f_i$ 的环同态
$\mathbf{Z}[x_1, \ldots, x_r] \to A$。由于
$x_1, \ldots, x_r$ 是 $\mathbf{Z}[x_1, \ldots, x_r]$ 中的正则序列，
可见 $x_1, \ldots, x_r$ 上的 Koszul 复形是
$\mathbf{Z} = \mathbf{Z}[x_1, \ldots, x_r]/(x_1, \ldots, x_r)$
在 $\mathbf{Z}[x_1, \ldots, x_r]$ 上的自由预解（参见《代数续篇》引理
\ref{more-algebra-lemma-regular-koszul-regular}）。因此，
$f_1, \ldots, f_r$ 为 Koszul 正则序列这一假设恰好意味着
$B = A \otimes_{\mathbf{Z}[x_1, \ldots, x_r]}^\mathbf{L} \mathbf{Z}$。
所以，由引理 \ref{cotangent-lemma-flat-base-change-cotangent-complex} 与
\ref{cotangent-lemma-special-case}，
$L_{B/A} = L_{\mathbf{Z}/\mathbf{Z}[x_1, \ldots, x_r]}
\otimes_\mathbf{Z}^\mathbf{L} B$。
\end{proof}

\begin{lemma}
\label{cotangent-lemma-mod-Koszul-regular-ideal}
设 $A \to B$ 为满环同态，其核 $I$ 是 Koszul 理想。
则 $L_{B/A}$ 与 $I/I^2[1]$ 拟同构。
\end{lemma}

\begin{proof}
在 $\Spec(A)$ 上局部地，理想 $I$ 由 Koszul 正则序列生成，参见
《代数续篇》定义 \ref{more-algebra-definition-regular-ideal}。
因此结论由引理 \ref{cotangent-lemma-flat-base-change-cotangent-complex} 得到。
\end{proof}

\begin{proposition}
\label{cotangent-proposition-cotangent-complex-local-complete-intersection}
设 $A \to B$ 为局部完全交同态。则 $L_{B/A}$ 是 Tor 振幅位于
$[-1, 0]$ 的完美复形。
\end{proposition}

\begin{proof}
选取核为 $J$ 的满射 $P = A[x_1, \ldots, x_n] \to B$。
由引理 \ref{cotangent-lemma-relation-with-naive-cotangent-complex}，可见
$J/J^2 \to \bigoplus B\text{d}x_i$ 与
$\tau_{\geq -1}L_{B/A}$ 拟同构。注意 $J/J^2$ 是有限投射模
（《代数续篇》引理
\ref{more-algebra-lemma-quasi-regular-ideal-finite-projective}），
故 $\tau_{\geq -1}L_{B/A}$ 是 Tor 振幅位于 $[-1, 0]$ 的完美复形。
因此只须证明当 $i \not \in [-1, 0]$ 时 $H^i(L_{B/A}) = 0$。
这由 (\ref{cotangent-equation-triangle})
$$
L_{P/A} \otimes_P^\mathbf{L} B \to L_{B/A} \to L_{B/P} \to
L_{P/A} \otimes_P^\mathbf{L} B[1]
$$
及引理 \ref{cotangent-lemma-mod-Koszul-regular-ideal} 得到：除非
$i \in \{-1, 0\}$，否则 $H^i(L_{B/P})$ 为零。（对左端项还使用了
引理 \ref{cotangent-lemma-cotangent-complex-polynomial-algebra}。）
\end{proof}








\section{张量积与余切复形}
\label{cotangent-section-tensor-product}

\noindent
设 $R$ 为环，$A$、$B$ 为 $R$-代数。本节讨论
$L_{A \otimes_R B/R}$。我们所需的大部分信息都包含在下图中：
\begin{equation}
\label{cotangent-equation-tensor-product}
\vcenter{
\xymatrix{
L_{A/R} \otimes_A^\mathbf{L} (A \otimes_R B) \ar[r] &
L_{A \otimes_R B/B} \ar[r] &
E \\
L_{A/R} \otimes_A^\mathbf{L} (A \otimes_R B) \ar[r] \ar@{=}[u] &
L_{A \otimes_R B/R} \ar[r] \ar[u] &
L_{A \otimes_R B/A} \ar[u] \\
 &
L_{B/R} \otimes_B^\mathbf{L} (A \otimes_R B) \ar[u] \ar@{=}[r] &
L_{B/R} \otimes_B^\mathbf{L} (A \otimes_R B) \ar[u]
}
}
\end{equation}
说明：中间一行是环同态 $R \to A \to A \otimes_R B$ 的基本三角
(\ref{cotangent-equation-triangle})。中间一列是环同态
$R \to B \to A \otimes_R B$ 的基本三角
(\ref{cotangent-equation-triangle})。其次，$E$ 是 $D(A \otimes_R B)$ 中
“嵌入”右上角的对象，也就是说，它使最上面一行和最右边一列都成为
可区别三角。将《导出范畴》命题 \ref{derived-proposition-9}
应用于左下方块（在缺失位置放置 $0$），可知这样的 $E$ 存在。
更明确地说，例如可以把 $E$ 定义为复形态射
$$
L_{A/R} \otimes_A^\mathbf{L} (A \otimes_R B) \oplus
L_{B/R} \otimes_B^\mathbf{L} (A \otimes_R B)
\longrightarrow
L_{A \otimes_R B/R}
$$
的锥（《导出范畴》定义 \ref{derived-definition-cone}），再应用 TR3
得到以 $E$ 为靶的两个态射。在 Tor 独立的情形下，对象 $E$ 为零。

\begin{lemma}
\label{cotangent-lemma-tensor-product-tor-independent}
若 $A$ 与 $B$ 是 Tor 独立的 $R$-代数，则
(\ref{cotangent-equation-tensor-product}) 中的对象 $E$ 为零。在此情形下有
$$
L_{A \otimes_R B/R} =
L_{A/R} \otimes_A^\mathbf{L} (A \otimes_R B) \oplus
L_{B/R} \otimes_B^\mathbf{L} (A \otimes_R B)
$$
，它由 $A \otimes_R B$-模复形
$L_{A/R} \otimes_R B \oplus L_{B/R} \otimes_R A $
表示。
\end{lemma}

\begin{proof}
前两个陈述直接由引理 \ref{cotangent-lemma-flat-base-change-cotangent-complex} 得到。
最后一个陈述成立，是因为 $L_{A/R}$ 是自由 $A$-模的复形，故
$L_{A/R} \otimes_A^\mathbf{L} (A \otimes_R B)$ 由
$L_{A/R} \otimes_A (A \otimes_R B) = L_{A/R} \otimes_R B$
表示。
\end{proof}

\noindent
一般地，对象 $E$ 有如下性质。

\begin{lemma}
\label{cotangent-lemma-tensor-product}
设 $R$ 为环，$A$、$B$ 为 $R$-代数。对象 $E$ 在
(\ref{cotangent-equation-tensor-product}) 中满足
$$
H^i(E) =
\left\{
\begin{matrix}
0 & \text{若} & i \geq -1 \\
\text{Tor}_1^R(A, B) & \text{若} & i = -2
\end{matrix}
\right.
$$
\end{lemma}

\begin{proof}
我们采用 $E$ 是态射
$L_{B/R} \otimes_B^\mathbf{L} (A \otimes_R B) \to L_{A \otimes_R B/A}$
之锥的描述。由引理 \ref{cotangent-lemma-compare-higher}，典范截断
$\tau_{\geq -2}L_{B/R}$ 与 $\tau_{\geq -2}L_{A \otimes_R B/A}$
由 Lichtenbaum--Schlessinger 复形
(\ref{cotangent-equation-lichtenbaum-schlessinger}) 计算。这些同构与函子性相容
（评注 \ref{cotangent-remark-functoriality-lichtenbaum-schlessinger}）。
因此，本证明使用 Lichtenbaum--Schlessinger 复形。

\medskip\noindent
选取 $R$ 上的多项式代数 $P$ 和满射 $P \to B$。选取该满射之核的
生成元 $f_t \in P$（$t \in T$）。令
$Rel \subset F = \bigoplus_{t \in T} P$ 为态射 $F \to P$ 的核，
该态射将对应于 $t$ 的基向量映到 $f_t$。令
$P_A = A \otimes_R P$ 且
$F_A = A \otimes_R F = P_A \otimes_P F$。令 $Rel_A$ 为
$F_A \to P_A$ 的核。利用正合列
$$
0 \to Rel \to F \to P \to B \to 0
$$
及 Tor 的标准短正合列，得到正合列
$$
A \otimes_R Rel \to Rel_A \to \text{Tor}_1^R(A, B) \to 0
$$
注意 $P_A \to A \otimes_R B$ 是满射，其核由 $P_A$ 中的元素
$1 \otimes f_t$ 生成。以 $TrivRel_A \subset Rel_A$ 表示由元素
$(\ldots, 1 \otimes f_{t'}, 0, \ldots,
0, - 1 \otimes f_t \otimes 1, 0, \ldots)$
% P11-E0032：冻结源在上一坐标中保留疑似多余的 \otimes 1。
生成的 $P_A$-子模。由于
$TrivRel \otimes_R A \to TrivRel_A$ 是满射，得到典范正合列
$$
A \otimes_R (Rel/TrivRel) \to Rel_A/TrivRel_A \to \text{Tor}_1^R(A, B) \to 0
$$
Lichtenbaum--Schlessinger 复形之间的态射由下图给出：
$$
\xymatrix{
Rel_A/TrivRel_A \ar[r] &
F_A \otimes_{P_A} (A \otimes_R B) \ar[r] &
\Omega_{P_A/A \otimes_R B} \otimes_{P_A} (A \otimes_R B) \\
Rel/TrivRel \ar[r] \ar[u]_{-2} &
F \otimes_P B \ar[r] \ar[u]_{-1} &
\Omega_{P/A} \otimes_P B \ar[u]_0
}
$$
% P11-E0033：冻结源在上图右上项保留与该复形基环不符的微分模下标。
注意，对源施用函子 $A \otimes_R - = P_A \otimes_P -$ 后，纵向态射
$-1$ 与 $-0$ 诱导同构，而纵向态射 $-2$ 恰好给出以上所见的那个态射，
它的余核就是所需的 Tor 模。
% P11-E0034：冻结源把图中标为 0 的纵向态射称作 -0。
\end{proof}












\section{环同态的形变与余切复形}
\label{cotangent-section-deformations}

\noindent
本节续接《形变理论》第 \ref{defos-section-deformations} 节，
我们强烈建议读者先阅读该节。首先取满环同态 $A' \to A$，其核是平方为零的
理想 $I$。再假设给定环同态 $A \to B$、$B$-模 $N$ 以及 $A$-模态射
$c : I \to N$。本节要问的是，能否找到适合下图中问号位置的对象：
\begin{equation}
\label{cotangent-equation-to-solve}
\vcenter{
\xymatrix{
0 \ar[r] & N \ar[r] & {?} \ar[r] & B \ar[r] & 0 \\
0 \ar[r] & I \ar[u]^c \ar[r] & A' \ar[u] \ar[r] & A \ar[u] \ar[r] & 0
}
}
\end{equation}
并且还要问该解（若存在）在何种意义下唯一。更确切地说，我们寻找
$A'$-代数的满射 $B' \to B$，其核为平方零理想并与 $N$ 等同，且
$A' \to B'$ 诱导给定态射 $c$。我们称 $B'$ 是
(\ref{cotangent-equation-to-solve}) 的一个{\it 解}。

\begin{lemma}
\label{cotangent-lemma-find-obstruction}
在上述情形中，有：
\begin{enumerate}
\item 存在典范元素 $\xi \in \Ext^2_B(L_{B/A}, N)$；它为零是
(\ref{cotangent-equation-to-solve}) 存在解的充分必要条件。
\item 若存在解，则解的同构类之集在 $\Ext^1_B(L_{B/A}, N)$
的作用下为主齐性空间。
\item 给定一个解 $B'$，所有适合 (\ref{cotangent-equation-to-solve}) 的
$B'$ 自同构之集典范同构于 $\Ext^0_B(L_{B/A}, N)$。
\end{enumerate}
\end{lemma}

\begin{proof}
借助等同 $\NL_{B/A} = \tau_{\geq -1}L_{B/A}$（引理
\ref{cotangent-lemma-relation-with-naive-cotangent-complex}）与
$H^0(L_{B/A}) = \Omega_{B/A}$（引理 \ref{cotangent-lemma-identify-H0}），
我们已在《形变理论》引理 \ref{defos-lemma-huge-diagram} 与
\ref{defos-lemma-choices} 中见过第 (2)、(3) 项。

\medskip\noindent
证明 (1)。粗略地说，将朴素余切复形替换为完整余切复形后，结论便由
《形变理论》评注 \ref{defos-remark-parametrize-solutions} 中的讨论得到。
下面给出更详细的说明。由《形变理论》引理
\ref{defos-lemma-parametrize-solutions} 和评注
\ref{defos-remark-parametrize-solutions}，存在元素
$$
\xi' \in
\Ext^1_A(\NL_{A/A'}, N) =
\Ext^1_B(\NL_{A/A'} \otimes_A^\mathbf{L} B, N) =
\Ext^1_B(L_{A/A'} \otimes_A^\mathbf{L} B, N)
$$
（关于这些等号，参见《形变理论》评注
\ref{defos-remark-parametrize-solutions}，并使用
$\NL_{A'/A} = \tau_{\geq -1} L_{A'/A}$）
% P11-E0035：冻结源在上一等式中把 A/A' 反写为 A'/A。
，使得存在解当且仅当此元素属于态射
$$
\Ext^1_B(\NL_{B/A'}, N) = \Ext^1_B(L_{B/A'}, N)
\longrightarrow
\Ext^1_B(L_{A/A'} \otimes_A^\mathbf{L} B, N)
$$
的像。$A' \to A \to B$ 的可区别三角
(\ref{cotangent-equation-triangle}) 给出长正合列
$$
\ldots \to
\Ext^1_B(L_{B/A'}, N) \to
\Ext^1_B(L_{A/A'} \otimes_A^\mathbf{L} B, N) \to
\Ext^2_B(L_{B/A}, N) \to \ldots
$$
因此，取 $\xi$ 为 $\xi'$ 的像即可。
\end{proof}





\section{模的 Atiyah 类}
\label{cotangent-section-atiyah}

\noindent
设 $A \to B$ 为环同态，$M$ 为 $B$-模。设 $P \to B$ 为
$\mathcal{C}_{B/A}$ 的对象（第 \ref{cotangent-section-compute-L-pi-shriek} 节）。
考虑主部的扩张
$$
0 \to \Omega_{P/A} \otimes_P M \to P^1_{P/A}(M) \to M \to 0
$$
，参见《代数》引理 \ref{algebra-lemma-sequence-of-principal-parts}。
由《代数》评注 \ref{algebra-remark-functoriality-principal-parts}，
此列对 $P$ 具有函子性。因此在 $\mathcal{C}_{B/A}$ 上得到
$\mathcal{O}$-模层的短正合列
$$
0 \to \Omega_{\mathcal{O}/\underline{A}} \otimes_\mathcal{O} \underline{M} \to
P^1_{\mathcal{O}/\underline{A}}(M) \to \underline{M} \to 0
$$
。由引理 \ref{cotangent-lemma-pi-shriek-standard} 及 $L_{B/A}$ 各项的平坦性，有
$L\pi_!(\Omega_{\mathcal{O}/\underline{A}} \otimes_\mathcal{O} \underline{M})
= L_{B/A} \otimes_B M = L_{B/A} \otimes_B^\mathbf{L} M$
。由引理 \ref{cotangent-lemma-pi-lower-shriek-constant-sheaf}，有
$L\pi_!(\underline{M}) = M$。因此得到可区别三角
\begin{equation}
\label{cotangent-equation-atiyah}
L_{B/A} \otimes_B^\mathbf{L} M \to
L\pi_!\left(P^1_{\mathcal{O}/\underline{A}}(M)\right) \to M
\to L_{B/A} \otimes_B^\mathbf{L} M [1]
\end{equation}
；它位于 $D(B)$ 中。这里使用《站点上的上同调》评注
\ref{sites-cohomology-remark-O-homology-B-homology-general}，
以便在 $D(B)$ 中而不只是在 $D(A)$ 中得到可区别三角。

\begin{definition}
\label{cotangent-definition-atiyah-class}
设 $A \to B$ 为环同态，$M$ 为 $B$-模。
(\ref{cotangent-equation-atiyah}) 中的态射
$M \to L_{B/A} \otimes_B^\mathbf{L} M[1]$
称为 $M$ 的{\it Atiyah 类}。
\end{definition}




\section{余切复形}
\label{cotangent-section-cotangent-complex}

\noindent
本节讨论站点上环层同态的余切复形。在后续各节中，我们将其特殊化，
得到赋环拓扑斯态射、赋环空间态射、概形态射、代数空间态射等的余切复形。

\medskip\noindent
设 $\mathcal{C}$ 为站点，以 $\Sh(\mathcal{C})$ 表示相应的拓扑斯。
设 $\mathcal{A}$ 为 $\mathcal{C}$ 上的环层，以
$\mathcal{A}\textit{-Alg}$ 表示 $\mathcal{A}$-代数范畴。
考虑一对伴随函子 $(U, V)$，其中
$V : \mathcal{A}\textit{-Alg} \to \Sh(\mathcal{C})$ 是遗忘函子，而
$U : \Sh(\mathcal{C}) \to \mathcal{A}\textit{-Alg}$ 将集合层
$\mathcal{E}$ 送到 $\mathcal{A}$ 上以 $\mathcal{E}$ 为变量集的多项式代数
$\mathcal{A}[\mathcal{E}]$。令 $X_\bullet$ 为
$\text{Fun}(\mathcal{A}\textit{-Alg}, \mathcal{A}\textit{-Alg})$
中的单纯形对象，它构造于《单纯形》第 \ref{simplicial-section-standard} 节。

\medskip\noindent
现在假设 $\mathcal{A} \to \mathcal{B}$ 是环层同态。于是
$\mathcal{B}$ 是范畴 $\mathcal{A}\textit{-Alg}$ 的对象。以
$\mathcal{P}_\bullet = X_\bullet(\mathcal{B})$ 表示所得的单纯形
$\mathcal{A}$-代数。回忆
$\mathcal{P}_0 = \mathcal{A}[\mathcal{B}]$,
$\mathcal{P}_1 = \mathcal{A}[\mathcal{A}[\mathcal{B}]]$，依此类推。
还要回忆，存在增广
$$
\epsilon : \mathcal{P}_\bullet \longrightarrow \mathcal{B}
$$
，其中把 $\mathcal{B}$ 视为常值单纯形 $\mathcal{A}$-代数。

\begin{definition}
\label{cotangent-definition-standard-resolution-sheaves-rings}
设 $\mathcal{C}$ 为站点，$\mathcal{A} \to \mathcal{B}$ 为
$\mathcal{C}$ 上的环层同态。$\mathcal{B}$ 在 $\mathcal{A}$ 上的
{\it 标准预解}是增广
$\epsilon : \mathcal{P}_\bullet \to \mathcal{B}$
，其各项为
$$
\mathcal{P}_0 = \mathcal{A}[\mathcal{B}],\quad
\mathcal{P}_1 = \mathcal{A}[\mathcal{A}[\mathcal{B}]],\quad \ldots
$$
，态射如上所构造。
\end{definition}

\noindent
有了这一定义，环层同态的余切复形定义如下。我们将使用《站点上的模》
第 \ref{sites-modules-section-differentials} 节中定义的微分模。

\begin{definition}
\label{cotangent-definition-cotangent-complex-morphism-sheaves-rings}
设 $\mathcal{C}$ 为站点，$\mathcal{A} \to \mathcal{B}$ 为
$\mathcal{C}$ 上的环层同态。{\it 余切复形}
$L_{\mathcal{B}/\mathcal{A}}$ 是与单纯形模
$$
\Omega_{\mathcal{P}_\bullet/\mathcal{A}}
\otimes_{\mathcal{P}_\bullet, \epsilon} \mathcal{B}
$$
相伴的 $\mathcal{B}$-模复形，其中
$\epsilon : \mathcal{P}_\bullet \to \mathcal{B}$ 是
$\mathcal{B}$ 在 $\mathcal{A}$ 上的标准预解。我们通常把
$L_{\mathcal{B}/\mathcal{A}}$ 看作 $D(\mathcal{B})$ 的对象。
\end{definition}

\noindent
这些构造满足与第 \ref{cotangent-section-functoriality} 节所讨论者类似的函子性。
确切地说，给定 $\mathcal{C}$ 上环层的交换图
\begin{equation}
\label{cotangent-equation-commutative-square-sheaves}
\vcenter{
\xymatrix{
\mathcal{B} \ar[r] & \mathcal{B}' \\
\mathcal{A} \ar[u] \ar[r] & \mathcal{A}' \ar[u]
}
}
\end{equation}
，存在典范的 $\mathcal{B}$-线性复形态射
$$
L_{\mathcal{B}/\mathcal{A}} \longrightarrow L_{\mathcal{B}'/\mathcal{A}'}
$$
，其构造如下。若 $\mathcal{P}_\bullet \to \mathcal{B}$ 是
$\mathcal{B}$ 在 $\mathcal{A}$ 上的标准预解，而
$\mathcal{P}'_\bullet \to \mathcal{B}'$ 是
$\mathcal{B}'$ 在 $\mathcal{A}'$ 上的标准预解，则存在单纯形
$\mathcal{A}$-代数的典范态射
$\mathcal{P}_\bullet \to \mathcal{P}'_\bullet$，它与增广
$\mathcal{P}_\bullet \to \mathcal{B}$ 与
$\mathcal{P}'_\bullet \to \mathcal{B}'$ 相容。态射
$$
\mathcal{P}_0 = \mathcal{A}[\mathcal{B}]
\longrightarrow
\mathcal{A}'[\mathcal{B}'] = \mathcal{P}'_0,
\quad
\mathcal{P}_1 = \mathcal{A}[\mathcal{A}[\mathcal{B}]]
\longrightarrow
\mathcal{A}'[\mathcal{A}'[\mathcal{B}']] = \mathcal{P}'_1
$$
等由给定态射 $\mathcal{A} \to \mathcal{A}'$ 与
$\mathcal{B} \to \mathcal{B}'$ 给出。所需态射
$L_{\mathcal{B}/\mathcal{A}} \to L_{\mathcal{B}'/\mathcal{A}'}$
便来自微分层上的相伴态射。

\begin{lemma}
\label{cotangent-lemma-pullback-cotangent-morphism-topoi}
设 $f : \Sh(\mathcal{D}) \to \Sh(\mathcal{C})$ 为拓扑斯态射，
$\mathcal{A} \to \mathcal{B}$ 为 $\mathcal{C}$ 上的环层同态。则
$f^{-1}L_{\mathcal{B}/\mathcal{A}} = L_{f^{-1}\mathcal{B}/f^{-1}\mathcal{A}}$。
\end{lemma}

\begin{proof}
图
$$
\xymatrix{
\mathcal{A}\textit{-Alg} \ar[d]_{f^{-1}} \ar[r] &
\Sh(\mathcal{C}) \ar@<1ex>[l] \ar[d]^{f^{-1}} \\
f^{-1}\mathcal{A}\textit{-Alg} \ar[r] & \Sh(\mathcal{D}) \ar@<1ex>[l]
}
$$
交换。
\end{proof}

\begin{lemma}
\label{cotangent-lemma-compute-L-morphism-sheaves-rings}
设 $\mathcal{C}$ 为站点，$\mathcal{A} \to \mathcal{B}$ 为
$\mathcal{C}$ 上的环层同态。则
$H^i(L_{\mathcal{B}/\mathcal{A}})$ 是与预层
$U \mapsto H^i(L_{\mathcal{B}(U)/\mathcal{A}(U)})$ 相伴的层。
\end{lemma}

\begin{proof}
在 $\mathcal{C}$ 上赋予混沌拓扑（此时预层都是层），所得站点记为
$\mathcal{C}'$。存在拓扑斯态射
$f : \Sh(\mathcal{C}) \to \Sh(\mathcal{C}')$，其中 $f_*$
是层到预层的包含，而 $f^{-1}$ 是层化。由引理
\ref{cotangent-lemma-pullback-cotangent-morphism-topoi}，只须对
$\mathcal{C}'$ 证明结论，也就是只须处理 $\mathcal{C}$ 具有混沌拓扑的情形。

\medskip\noindent
若 $\mathcal{C}$ 带有混沌拓扑，则
$L_{\mathcal{B}/\mathcal{A}}(U)$ 等于
$L_{\mathcal{B}(U)/\mathcal{A}(U)}$，因为图
$$
\xymatrix{
\mathcal{A}\textit{-Alg} \ar[d]_{\text{$U$ 上的截面}} \ar[r] &
\Sh(\mathcal{C}) \ar@<1ex>[l] \ar[d]^{\text{$U$ 上的截面}} \\
\mathcal{A}(U)\textit{-Alg} \ar[r] & \textit{Sets} \ar@<1ex>[l]
}
$$
交换。
\end{proof}

\begin{remark}
\label{cotangent-remark-map-sections-over-U}
由引理 \ref{cotangent-lemma-compute-L-morphism-sheaves-rings} 的证明显然可知，
对任意 $U \in \Ob(\mathcal{C})$，存在 $\mathcal{B}(U)$-模复形的典范态射
$L_{\mathcal{B}(U)/\mathcal{A}(U)} \to L_{\mathcal{B}/\mathcal{A}}(U)$
。此外，这些态射与限制态射相容，而复形
$L_{\mathcal{B}/\mathcal{A}}$ 是规则
$U \mapsto L_{\mathcal{B}(U)/\mathcal{A}(U)}$ 的层化。
\end{remark}

\begin{lemma}
\label{cotangent-lemma-H0-L-morphism-sheaves-rings}
设 $\mathcal{C}$ 为站点，$\mathcal{A} \to \mathcal{B}$ 为
$\mathcal{C}$ 上的环层同态。则
$H^0(L_{\mathcal{B}/\mathcal{A}}) = \Omega_{\mathcal{B}/\mathcal{A}}$。
\end{lemma}

\begin{proof}
由引理 \ref{cotangent-lemma-compute-L-morphism-sheaves-rings}、
\ref{cotangent-lemma-identify-H0} 及《站点上的模》引理
\ref{sites-modules-lemma-differentials-sheafify} 得到。
\end{proof}

\begin{lemma}
\label{cotangent-lemma-compute-L-product-sheaves-rings}
设 $\mathcal{C}$ 为站点，$\mathcal{A} \to \mathcal{B}$ 与
$\mathcal{A} \to \mathcal{B}'$ 为 $\mathcal{C}$ 上的环层同态。则
$$
L_{\mathcal{B} \times \mathcal{B}'/\mathcal{A}}
\longrightarrow
L_{\mathcal{B}/\mathcal{A}} \oplus L_{\mathcal{B}'/\mathcal{A}}
$$
在 $D(\mathcal{B} \times \mathcal{B}')$ 中是同构。
\end{lemma}

\begin{proof}
由引理 \ref{cotangent-lemma-compute-L-morphism-sheaves-rings}，只须对环同态证明。
对环而言，这正是引理 \ref{cotangent-lemma-cotangent-complex-product}。
\end{proof}

\noindent
环层余切复形的基本三角是环同态情形之结果的直接推论。

\begin{lemma}
\label{cotangent-lemma-triangle-sheaves-rings}
设 $\mathcal{D}$ 为站点，$\mathcal{A} \to \mathcal{B} \to \mathcal{C}$
为 $\mathcal{D}$ 上的环层同态。存在典范可区别三角
$$
L_{\mathcal{B}/\mathcal{A}} \otimes_\mathcal{B}^\mathbf{L} \mathcal{C}
\to L_{\mathcal{C}/\mathcal{A}} \to L_{\mathcal{C}/\mathcal{B}} \to
L_{\mathcal{B}/\mathcal{A}} \otimes_\mathcal{B}^\mathbf{L} \mathcal{C}[1]
$$
；它位于 $D(\mathcal{C})$ 中。
\end{lemma}

\begin{proof}
我们将使用评注 \ref{cotangent-remark-triangle} 与 \ref{cotangent-remark-explicit-map}
中描述的方法构造该三角，并直接使用其中提到的结果。依照这些评注，
先在 $\mathcal{B} \to \mathcal{C}$ 是环层单射的情形下构造三角。
在此情形下，令
\begin{enumerate}
\item $\mathcal{P}_\bullet$ 为 $\mathcal{B}$ 在 $\mathcal{A}$ 上的标准预解；
\item $\mathcal{Q}_\bullet$ 为 $\mathcal{C}$ 在 $\mathcal{A}$ 上的标准预解；
\item $\mathcal{R}_\bullet$ 为 $\mathcal{C}$ 在 $\mathcal{B}$ 上的标准预解；
\item $\mathcal{S}_\bullet$ 为 $\mathcal{B}$ 在 $\mathcal{B}$ 上的标准预解；
\item $\overline{\mathcal{Q}}_\bullet =
\mathcal{Q}_\bullet \otimes_{\mathcal{P}_\bullet} \mathcal{B}$；并且
\item $\overline{\mathcal{R}}_\bullet =
\mathcal{R}_\bullet \otimes_{\mathcal{S}_\bullet} \mathcal{B}$。
\end{enumerate}
所求可区别三角是与单纯形 $\mathcal{C}$-模的短正合列
$$
0 \to
\Omega_{\mathcal{P}_\bullet/\mathcal{A}}
\otimes_{\mathcal{P}_\bullet} \mathcal{C} \to
\Omega_{\mathcal{Q}_\bullet/\mathcal{A}}
\otimes_{\mathcal{Q}_\bullet} \mathcal{C} \to
\Omega_{\overline{\mathcal{Q}}_\bullet/\mathcal{B}}
\otimes_{\overline{\mathcal{Q}}_\bullet} \mathcal{C} \to 0
$$
相伴的可区别三角。前两项等于引理陈述中三角的前两项。
将最后一项与 $L_{\mathcal{C}/\mathcal{B}}$ 等同，使用复形的拟同构
$$
L_{\mathcal{C}/\mathcal{B}} =
\Omega_{\mathcal{R}_\bullet/\mathcal{B}}
\otimes_{\mathcal{R}_\bullet} \mathcal{C}
\longrightarrow
\Omega_{\overline{\mathcal{R}}_\bullet/\mathcal{B}}
\otimes_{\overline{\mathcal{R}}_\bullet} \mathcal{C}
\longleftarrow
\Omega_{\overline{\mathcal{Q}}_\bullet/\mathcal{B}}
\otimes_{\overline{\mathcal{Q}}_\bullet} \mathcal{C}
$$
上述所有构造都可以先在预层层次完成，再进行层化。因此，要证明列正合，
或证明态射是拟同构，只须对环同态
% P11-E0036：冻结源英文把复数主语 maps 误写为单数 map。
$\mathcal{A}(U) \to \mathcal{B}(U) \to \mathcal{C}(U)$
证明相应陈述；这些结论已经知道。这样便完成了
$\mathcal{B} \to \mathcal{C}$ 为单射时的证明。

\medskip\noindent
一般地，必要时以 $\mathcal{B} \times \mathcal{C}$ 替换
$\mathcal{C}$，从而归约到 $\mathcal{B} \to \mathcal{C}$ 为单射的情形。
由评注 \ref{cotangent-remark-triangle} 中的论证和引理
\ref{cotangent-lemma-compute-L-product-sheaves-rings}，这是可行的。
\end{proof}

\begin{lemma}
\label{cotangent-lemma-stalk-cotangent-complex}
设 $\mathcal{C}$ 为站点，$\mathcal{A} \to \mathcal{B}$ 为
$\mathcal{C}$ 上的环层同态。若 $p$ 是 $\mathcal{C}$ 的点，则
$(L_{\mathcal{B}/\mathcal{A}})_p = L_{\mathcal{B}_p/\mathcal{A}_p}$。
\end{lemma}

\begin{proof}
这是引理 \ref{cotangent-lemma-pullback-cotangent-morphism-topoi} 的特殊情形。
\end{proof}

\noindent
关于朴素余切复形的构造及其性质，参见《站点上的模》
第 \ref{sites-modules-section-netherlander} 节。

\begin{lemma}
\label{cotangent-lemma-compare-cotangent-complex-with-naive}
设 $\mathcal{C}$ 为站点，$\mathcal{A} \to \mathcal{B}$ 为
$\mathcal{C}$ 上的环层同态。存在典范态射
$L_{\mathcal{B}/\mathcal{A}} \to \NL_{\mathcal{B}/\mathcal{A}}$
，它将朴素余切复形与截断
$\tau_{\geq -1}L_{\mathcal{B}/\mathcal{A}}$ 等同。
\end{lemma}

\begin{proof}
令 $\mathcal{P}_\bullet$ 为 $\mathcal{B}$ 在 $\mathcal{A}$ 上的标准预解，
并令 $\mathcal{I} = \Ker(\mathcal{A}[\mathcal{B}] \to \mathcal{B})$。
回忆 $\mathcal{P}_0 = \mathcal{A}[\mathcal{B}]$。该引理中的态射
由交换图
$$
\xymatrix{
L_{\mathcal{B}/\mathcal{A}} \ar[d] & \ldots \ar[r] &
\Omega_{\mathcal{P}_2/\mathcal{A}} \otimes_{\mathcal{P}_2} \mathcal{B}
\ar[r] \ar[d] &
\Omega_{\mathcal{P}_1/\mathcal{A}} \otimes_{\mathcal{P}_1} \mathcal{B}
\ar[r] \ar[d] &
\Omega_{\mathcal{P}_0/\mathcal{A}} \otimes_{\mathcal{P}_0} \mathcal{B}
\ar[d] \\
\NL_{\mathcal{B}/\mathcal{A}} & \ldots \ar[r] &
0 \ar[r] & 
\mathcal{I}/\mathcal{I}^2 \ar[r] &
\Omega_{\mathcal{P}_0/\mathcal{A}} \otimes_{\mathcal{P}_0} \mathcal{B}
}
$$
给出。我们将局部截面 $\text{d}f \otimes b$ 映到
$(d_0(f) - d_1(f))b$ 在 $\mathcal{I}/\mathcal{I}^2$ 中的类，
由此构造以 $\mathcal{I}/\mathcal{I}^2$ 为靶的向下箭头。
这里 $d_i : \mathcal{P}_1 \to \mathcal{P}_0$（$i = 0, 1$）
是单纯形结构的两个面态射。这是有意义的，因为 $d_0-d_1$ 将
$\mathcal{P}_1$ 映入
$\mathcal{I} = \Ker(\mathcal{P}_0 \to \mathcal{B})$。
此规则良定义的验证从略。我们的态射与微分
$\Omega_{\mathcal{P}_1/\mathcal{A}} \otimes_{\mathcal{P}_1} \mathcal{B}
\to \Omega_{\mathcal{P}_0/\mathcal{A}} \otimes_{\mathcal{P}_0} \mathcal{B}$
相容，因为该微分将局部截面 $\text{d}f \otimes b$ 映到
$\text{d}(d_0(f) - d_1(f)) \otimes b$。此外，微分
$\Omega_{\mathcal{P}_2/\mathcal{A}} \otimes_{\mathcal{P}_2} \mathcal{B}
\to \Omega_{\mathcal{P}_1/\mathcal{A}} \otimes_{\mathcal{P}_1} \mathcal{B}$
将局部截面 $\text{d}f \otimes b$ 映到
$\text{d}(d_0(f) - d_1(f) + d_2(f)) \otimes b$，后者被我们的向下箭头
消去。因此得到了复形态射。

\medskip\noindent
为了证明我们的态射在上同调层 $H^0$ 与 $H^{-1}$ 上诱导同构，
作如下论证。令 $\mathcal{C}'$ 为与 $\mathcal{C}$ 具有相同底范畴、
但赋予混沌拓扑的站点。令
$f : \Sh(\mathcal{C}) \to \Sh(\mathcal{C}')$ 为拉回函子是层化的
拓扑斯态射。把给定态射 $\mathcal{A}' \to \mathcal{B}'$
看作 $\mathcal{C}'$ 上的环层同态。上述构造在 $\mathcal{C}'$ 上给出态射
$L_{\mathcal{B}'/\mathcal{A}'} \to \NL_{\mathcal{B}'/\mathcal{A}'}$，
它在 $\mathcal{C}'$ 任一对象 $U$ 上的值正是态射
$$
L_{\mathcal{B}(U)/\mathcal{A}(U)} \to \NL_{\mathcal{B}(U)/\mathcal{A}(U)}
$$
；这是评注 \ref{cotangent-remark-explicit-comparison-map} 中的态射，并在
$H^0$ 与 $H^{-1}$ 上诱导同构。由于
$f^{-1}L_{\mathcal{B}'/\mathcal{A}'} = L_{\mathcal{B}/\mathcal{A}}$
（引理 \ref{cotangent-lemma-pullback-cotangent-morphism-topoi}），且
$f^{-1}\NL_{\mathcal{B}'/\mathcal{A}'} = \NL_{\mathcal{B}/\mathcal{A}}$
（《站点上的模》引理 \ref{sites-modules-lemma-pullback-NL}），
该引理得证。
\end{proof}











\section{模层的 Atiyah 类}
\label{cotangent-section-atiyah-general}

\noindent
设 $\mathcal{C}$ 为站点，$\mathcal{A} \to \mathcal{B}$ 为环层同态，
$\mathcal{F}$ 为 $\mathcal{B}$-模层。令
$\mathcal{P}_\bullet \to \mathcal{B}$ 为 $\mathcal{B}$ 在
$\mathcal{A}$ 上的标准预解（第 \ref{cotangent-section-cotangent-complex} 节）。
对每个 $n \geq 0$，考虑主部的扩张
\begin{equation}
\label{cotangent-equation-atiyah-extension}
0 \to
\Omega_{\mathcal{P}_n/\mathcal{A}} \otimes_{\mathcal{P}_n} \mathcal{F} \to
\mathcal{P}^1_{\mathcal{P}_n/\mathcal{A}}(\mathcal{F}) \to
\mathcal{F} \to 0
\end{equation}
，参见《站点上的模》引理
\ref{sites-modules-lemma-sequence-of-principal-parts}。此构造的函子性
（《站点上的模》评注
\ref{sites-modules-remark-functoriality-principal-parts}）告诉我们，
(\ref{cotangent-equation-atiyah-extension}) 是单纯形
$\mathcal{P}_\bullet$-模短正合列的第 $n$ 次部分（《站点上的上同调》
第 \ref{sites-cohomology-section-simplicial-modules} 节）。
使用《站点上的上同调》评注
\ref{sites-cohomology-remark-homology-augmentation} 中的函子
$L\pi_! : D(\mathcal{P}_\bullet) \to D(\mathcal{B})$
（这里使用 $\mathcal{P}_\bullet \to \mathcal{A}$ 是预解），
% P11-E0037：冻结源把标准预解的增广靶写成 A，而不是 B。
得到可区别三角
\begin{equation}
\label{cotangent-equation-atiyah-general}
L_{\mathcal{B}/\mathcal{A}} \otimes_\mathcal{B}^\mathbf{L} \mathcal{F} \to
L\pi_!\left(\mathcal{P}^1_{\mathcal{P}_\bullet/\mathcal{A}}(\mathcal{F})\right)
\to \mathcal{F} \to
L_{\mathcal{B}/\mathcal{A}} \otimes_\mathcal{B}^\mathbf{L} \mathcal{F} [1]
\end{equation}
；它位于 $D(\mathcal{B})$ 中。

\begin{definition}
\label{cotangent-definition-atiyah-class-general}
设 $\mathcal{C}$ 为站点，$\mathcal{A} \to \mathcal{B}$ 为环层同态，
$\mathcal{F}$ 为 $\mathcal{B}$-模层。
(\ref{cotangent-equation-atiyah-general}) 中的态射 $\mathcal{F} \to
L_{\mathcal{B}/\mathcal{A}} \otimes_\mathcal{B}^\mathbf{L} \mathcal{F}[1]$
称为 $\mathcal{F}$ 的{\it Atiyah 类}。
\end{definition}









\section{赋环空间态射的余切复形}
\label{cotangent-section-cotangent-morphism-ringed-spaces}

\noindent
赋环空间态射的余切复形用上面定义的余切复形来定义。

\begin{definition}
\label{cotangent-definition-cotangent-complex-morphism-ringed-spaces}
设 $f : (X, \mathcal{O}_X) \to (S, \mathcal{O}_S)$ 为赋环空间态射。
$f$ 的{\it 余切复形} $L_f$ 是
$L_f = L_{\mathcal{O}_X/f^{-1}\mathcal{O}_S}$。
我们也使用记号
$L_f = L_{X/S} = L_{\mathcal{O}_X/\mathcal{O}_S}$。
\end{definition}

\noindent
更确切地说，这意味着考虑环层同态
$f^\sharp : f^{-1}\mathcal{O}_S \to \mathcal{O}_X$ 的余切复形
（定义 \ref{cotangent-definition-cotangent-complex-morphism-sheaves-rings}）；
这些环层位于拓扑空间 $X$ 所对应的站点上（《站点》例
\ref{sites-example-site-topological}）。

\begin{lemma}
\label{cotangent-lemma-H0-L-morphism-ringed-spaces}
设 $f : (X, \mathcal{O}_X) \to (S, \mathcal{O}_S)$ 为赋环空间态射。
则 $H^0(L_{X/S}) = \Omega_{X/S}$。
\end{lemma}

\begin{proof}
这是引理 \ref{cotangent-lemma-H0-L-morphism-sheaves-rings} 的特殊情形。
\end{proof}

\begin{lemma}
\label{cotangent-lemma-triangle-ringed-spaces}
设 $f : X \to Y$ 与 $g : Y \to Z$ 为赋环空间态射。则存在典范可区别三角
$$
Lf^* L_{Y/Z} \to L_{X/Z} \to L_{X/Y} \to Lf^*L_{Y/Z}[1]
$$
；它位于 $D(\mathcal{O}_X)$ 中。
\end{lemma}

\begin{proof}
令 $h = g \circ f$，于是
$h^{-1}\mathcal{O}_Z = f^{-1}g^{-1}\mathcal{O}_Z$。由引理
\ref{cotangent-lemma-pullback-cotangent-morphism-topoi}，有
$f^{-1}L_{Y/Z} = L_{f^{-1}\mathcal{O}_Y/h^{-1}\mathcal{O}_Z}$
，且这是平坦 $f^{-1}\mathcal{O}_Y$-模的复形。因此，上述可区别三角
是引理 \ref{cotangent-lemma-triangle-sheaves-rings} 中可区别三角的一个实例，
其中 $\mathcal{A} = h^{-1}\mathcal{O}_Z$、
$\mathcal{B} = f^{-1}\mathcal{O}_Y$，且 $\mathcal{C} = \mathcal{O}_X$。
\end{proof}

\begin{lemma}
\label{cotangent-lemma-compare-cotangent-complex-with-naive-ringed-spaces}
设 $f : (X, \mathcal{O}_X) \to (Y, \mathcal{O}_Y)$ 为赋环空间态射。
存在典范态射 $L_{X/Y} \to \NL_{X/Y}$，它将朴素余切复形与截断
$\tau_{\geq -1}L_{X/Y}$ 等同。
\end{lemma}

\begin{proof}
这是引理 \ref{cotangent-lemma-compare-cotangent-complex-with-naive} 的特殊情形。
\end{proof}






\section{赋环空间的形变与余切复形}
\label{cotangent-section-deformations-ringed-spaces}

\noindent
本节续接《形变理论》第
\ref{defos-section-deformations-ringed-spaces} 节，我们强烈建议读者先阅读该节。
先简要回顾设置。给定赋环空间的一阶加厚
$t : (S, \mathcal{O}_S) \to (S', \mathcal{O}_{S'})$，其中
$\mathcal{J} = \Ker(t^\sharp)$；给定赋环空间态射
$f : (X, \mathcal{O}_X) \to (S, \mathcal{O}_S)$、$\mathcal{O}_X$-模
$\mathcal{G}$，以及模层的 $f$-态射
$c : \mathcal{J} \to \mathcal{G}$。我们要问的是，能否找到适合下图问号位置的对象：
\begin{equation}
\label{cotangent-equation-to-solve-ringed-spaces}
\vcenter{
\xymatrix{
0 \ar[r] & \mathcal{G} \ar[r] & {?} \ar[r] & \mathcal{O}_X \ar[r] & 0 \\
0 \ar[r] & \mathcal{J} \ar[u]^c \ar[r] & \mathcal{O}_{S'} \ar[u] \ar[r] &
\mathcal{O}_S \ar[u] \ar[r] & 0
}
}
\end{equation}
并且还要问该解（若存在）在何种意义下唯一。更确切地说，我们寻找一阶加厚
$i : (X, \mathcal{O}_X) \to (X', \mathcal{O}_{X'})$
和《形变理论》方程 (\ref{defos-equation-morphism-thickenings})
中的加厚态射 $(f, f')$，其中 $\Ker(i^\sharp)$ 与 $\mathcal{G}$ 等同，
且 $(f')^\sharp$ 诱导给定态射 $c$。我们称 $X'$ 是
(\ref{cotangent-equation-to-solve-ringed-spaces}) 的一个{\it 解}。

\begin{lemma}
\label{cotangent-lemma-find-obstruction-ringed-spaces}
在上述情形中，有：
\begin{enumerate}
\item 存在典范元素
$\xi \in \Ext^2_{\mathcal{O}_X}(L_{X/S}, \mathcal{G})$
；它为零是 (\ref{cotangent-equation-to-solve-ringed-spaces}) 存在解的充分必要条件。
\item 若存在解，则解的同构类之集在
$\Ext^1_{\mathcal{O}_X}(L_{X/S}, \mathcal{G})$ 的作用下为主齐性空间。
\item 给定一个解 $X'$，所有适合
(\ref{cotangent-equation-to-solve-ringed-spaces}) 的 $X'$ 自同构之集典范同构于
$\Ext^0_{\mathcal{O}_X}(L_{X/S}, \mathcal{G})$。
\end{enumerate}
\end{lemma}

\begin{proof}
借助等同 $\NL_{X/S} = \tau_{\geq -1}L_{X/S}$（引理
\ref{cotangent-lemma-compare-cotangent-complex-with-naive-ringed-spaces}）与
$H^0(L_{X/S}) = \Omega_{X/S}$（引理
\ref{cotangent-lemma-H0-L-morphism-ringed-spaces}），我们已在《形变理论》引理
\ref{defos-lemma-huge-diagram-ringed-spaces} 与
\ref{defos-lemma-choices-ringed-spaces} 中见过第 (2)、(3) 项。

\medskip\noindent
证明 (1)。粗略地说，将朴素余切复形替换为完整余切复形后，结论便由
《形变理论》评注 \ref{defos-remark-parametrize-solutions-ringed-spaces}
中的讨论得到。下面给出更详细的说明。由《形变理论》引理
\ref{defos-lemma-parametrize-solutions-ringed-spaces}，存在元素
$$
\xi' \in
\Ext^1_{\mathcal{O}_X}(Lf^*\NL_{S/S'}, \mathcal{G}) =
\Ext^1_{\mathcal{O}_X}(Lf^*L_{S/S'}, \mathcal{G})
$$
，使得存在解当且仅当此元素属于态射
$$
\Ext^1_{\mathcal{O}_X}(NL_{X/S'}, \mathcal{G}) =
% P11-E0038：冻结源在上一项用普通字母 NL，而不是本章宏 \NL。
\Ext^1_{\mathcal{O}_X}(L_{X/S'}, \mathcal{G})
\longrightarrow
\Ext^1_{\mathcal{O}_X}(Lf^*L_{S/S'}, \mathcal{G})
$$
的像。引理 \ref{cotangent-lemma-triangle-ringed-spaces} 中
$X \to S \to S'$ 的可区别三角给出长正合列
$$
\ldots \to
\Ext^1_{\mathcal{O}_X}(L_{X/S'}, \mathcal{G}) \to
\Ext^1_{\mathcal{O}_X}(Lf^*L_{S/S'}, \mathcal{G}) \to
\Ext^2_{\mathcal{O}_X}(L_{X/S}, \mathcal{G}) \to \ldots
$$
因此，取 $\xi$ 为 $\xi'$ 的像即可。
\end{proof}










\section{赋环拓扑斯态射的余切复形}
\label{cotangent-section-cotangent-morphism-ringed-topoi}

\noindent
赋环拓扑斯态射的余切复形用上面定义的余切复形来定义。

\begin{definition}
\label{cotangent-definition-cotangent-complex-morphism-ringed-topoi}
设 $(f, f^\sharp) : (\Sh(\mathcal{C}), \mathcal{O}_\mathcal{C}) \to
(\Sh(\mathcal{D}), \mathcal{O}_\mathcal{D})$ 为赋环拓扑斯态射。
$f$ 的{\it 余切复形} $L_f$ 是
$L_f = L_{\mathcal{O}_\mathcal{C}/f^{-1}\mathcal{O}_\mathcal{D}}$。
有时也写作 $L_f = L_{\mathcal{O}_\mathcal{C}/\mathcal{O}_\mathcal{D}}$。
\end{definition}

\noindent
此定义适用于许多情形，但并不总能给出预期的对象。例如，若
$f : X \to Y$ 是概形态射，则 $f$ 诱导大 \'etale 站点的态射
$f_{big} : (\Sch/X)_\etale \to (\Sch/Y)_\etale$
，后者是赋环拓扑斯态射（《下降》评注
\ref{descent-remark-change-topologies-ringed}）。然而，由于
$(f_{big})^\sharp$ 是同构，$L_{f_{big}} = 0$。另一方面，若在把 $f$
看作底层 Zariski 赋环拓扑斯之间态射的意义下取 $L_f$，那么 $L_f$
确实与余切复形 $L_{X/Y}$（定义见下文）一致，其零次上同调层为
$\Omega_{X/Y}$。

\begin{lemma}
\label{cotangent-lemma-H0-L-morphism-ringed-topoi}
设 $f : (\Sh(\mathcal{C}), \mathcal{O}) \to
(\Sh(\mathcal{B}), \mathcal{O}_\mathcal{B})$ 为赋环拓扑斯态射。
则 $H^0(L_f) = \Omega_f$。
\end{lemma}

\begin{proof}
这是引理 \ref{cotangent-lemma-H0-L-morphism-sheaves-rings} 的特殊情形。
\end{proof}

\begin{lemma}
\label{cotangent-lemma-triangle-ringed-topoi}
设 $f : (\Sh(\mathcal{C}_1), \mathcal{O}_1) \to
(\Sh(\mathcal{C}_2), \mathcal{O}_2)$ 与
$g : (\Sh(\mathcal{C}_2), \mathcal{O}_2) \to
(\Sh(\mathcal{C}_3), \mathcal{O}_3)$ 为赋环拓扑斯态射。
则存在典范可区别三角
$$
Lf^* L_g \to L_{g \circ f} \to L_f \to Lf^*L_g[1]
$$
；它位于 $D(\mathcal{O}_1)$ 中。
\end{lemma}

\begin{proof}
令 $h = g \circ f$，于是
$h^{-1}\mathcal{O}_3 = f^{-1}g^{-1}\mathcal{O}_3$。由引理
\ref{cotangent-lemma-pullback-cotangent-morphism-topoi}，有
$f^{-1}L_g = L_{f^{-1}\mathcal{O}_2/h^{-1}\mathcal{O}_3}$
，且这是平坦 $f^{-1}\mathcal{O}_2$-模的复形。因此，上述可区别三角
是引理 \ref{cotangent-lemma-triangle-sheaves-rings} 中可区别三角的一个实例，
其中 $\mathcal{A} = h^{-1}\mathcal{O}_3$、
$\mathcal{B} = f^{-1}\mathcal{O}_2$，且 $\mathcal{C} = \mathcal{O}_1$。
\end{proof}

\begin{lemma}
\label{cotangent-lemma-compare-cotangent-complex-with-naive-ringed-topoi}
设 $f : (\Sh(\mathcal{C}), \mathcal{O}) \to
(\Sh(\mathcal{B}), \mathcal{O}_\mathcal{B})$ 为赋环拓扑斯态射。
存在典范态射 $L_f \to \NL_f$，它将朴素余切复形与截断
$\tau_{\geq -1}L_f$ 等同。
\end{lemma}

\begin{proof}
这是引理 \ref{cotangent-lemma-compare-cotangent-complex-with-naive} 的特殊情形。
\end{proof}








\section{赋环拓扑斯的形变与余切复形}
\label{cotangent-section-deformations-ringed-topoi}

\noindent
本节续接《形变理论》第
\ref{defos-section-deformations-ringed-topoi} 节，我们强烈建议读者先阅读该节。
先简要回顾设置。给定赋环拓扑斯的一阶加厚
$t : (\Sh(\mathcal{B}), \mathcal{O}_\mathcal{B}) \to
(\Sh(\mathcal{B}'), \mathcal{O}_{\mathcal{B}'})$，其中
$\mathcal{J} = \Ker(t^\sharp)$；给定赋环拓扑斯态射
$f : (\Sh(\mathcal{C}), \mathcal{O}) \to
(\Sh(\mathcal{B}), \mathcal{O}_\mathcal{B})$、$\mathcal{O}$-模
$\mathcal{G}$，以及 $f^{-1}\mathcal{O}_\mathcal{B}$-模层的态射
$f^{-1}\mathcal{J} \to \mathcal{G}$。我们要问的是，能否找到适合下图问号位置的对象：
\begin{equation}
\label{cotangent-equation-to-solve-ringed-topoi}
\vcenter{
\xymatrix{
0 \ar[r] & \mathcal{G} \ar[r] & {?} \ar[r] & \mathcal{O} \ar[r] & 0 \\
0 \ar[r] & f^{-1}\mathcal{J} \ar[u]^c \ar[r] &
f^{-1}\mathcal{O}_{\mathcal{B}'} \ar[u] \ar[r] &
f^{-1}\mathcal{O}_\mathcal{B} \ar[u] \ar[r] & 0
}
}
\end{equation}
并且还要问该解（若存在）在何种意义下唯一。更确切地说，我们寻找一阶加厚
$i : (\Sh(\mathcal{C}), \mathcal{O}) \to (\Sh(\mathcal{C}'), \mathcal{O}')$
和《形变理论》方程
(\ref{defos-equation-morphism-thickenings-ringed-topoi}) 中的加厚态射
$(f,f')$，其中 $\Ker(i^\sharp)$ 与 $\mathcal{G}$ 等同，且
$(f')^\sharp$ 诱导给定态射 $c$。我们称
$(\Sh(\mathcal{C}'), \mathcal{O}')$ 是
(\ref{cotangent-equation-to-solve-ringed-topoi}) 的一个{\it 解}。

\begin{lemma}
\label{cotangent-lemma-find-obstruction-ringed-topoi}
在上述情形中，有：
\begin{enumerate}
\item 存在典范元素
$\xi \in \Ext^2_\mathcal{O}(L_f, \mathcal{G})$
；它为零是 (\ref{cotangent-equation-to-solve-ringed-topoi}) 存在解的充分必要条件。
\item 若存在解，则解的同构类之集在
$\Ext^1_\mathcal{O}(L_f, \mathcal{G})$ 的作用下为主齐性空间。
\item 给定一个解 $X'$，所有适合
(\ref{cotangent-equation-to-solve-ringed-topoi}) 的 $X'$ 自同构之集典范同构于
$\Ext^0_\mathcal{O}(L_f, \mathcal{G})$。
% P11-E0039：冻结源在拓扑斯情形中沿用了未定义的 X' 记号。
\end{enumerate}
\end{lemma}

\begin{proof}
借助等同 $\NL_f = \tau_{\geq -1}L_f$（引理
\ref{cotangent-lemma-compare-cotangent-complex-with-naive-ringed-topoi}）与
$H^0(L_f) = \Omega_f$（引理 \ref{cotangent-lemma-H0-L-morphism-ringed-topoi}），
我们已在《形变理论》引理 \ref{defos-lemma-huge-diagram-ringed-topoi}
与 \ref{defos-lemma-choices-ringed-topoi} 中见过第 (2)、(3) 项。

\medskip\noindent
证明 (1)。为使记号与《形变理论》第
\ref{defos-section-deformations-ringed-topoi} 节一致，我们写
$\NL_f = \NL_{\mathcal{O}/\mathcal{O}_\mathcal{B}}$ 及
$L_f = L_{\mathcal{O}/\mathcal{O}_\mathcal{B}}$，并对态射 $t$ 与
$t \circ f$ 使用类似记号。由《形变理论》引理
\ref{defos-lemma-parametrize-solutions-ringed-topoi}，存在元素
$$
\xi' \in
\Ext^1_\mathcal{O}(
Lf^*\NL_{\mathcal{O}_\mathcal{B}/\mathcal{O}_{\mathcal{B}'}}, \mathcal{G}) =
\Ext^1_\mathcal{O}(
Lf^*L_{\mathcal{O}_\mathcal{B}/\mathcal{O}_{\mathcal{B}'}}, \mathcal{G})
$$
，使得存在解当且仅当此元素属于态射
$$
\Ext^1_\mathcal{O}(
\NL_{\mathcal{O}/\mathcal{O}_{\mathcal{B}'}}, \mathcal{G}) =
\Ext^1_\mathcal{O}(
L_{\mathcal{O}/\mathcal{O}_{\mathcal{B}'}}, \mathcal{G})
\longrightarrow
\Ext^1_\mathcal{O}(
Lf^*L_{\mathcal{O}_\mathcal{B}/\mathcal{O}_{\mathcal{B}'}}, \mathcal{G})
$$
的像。引理 \ref{cotangent-lemma-triangle-ringed-topoi} 中 $f$ 与 $t$
的可区别三角给出长正合列
$$
\ldots \to
\Ext^1_\mathcal{O}(
L_{\mathcal{O}/\mathcal{O}_{\mathcal{B}'}}, \mathcal{G}) \to
\Ext^1_\mathcal{O}(
Lf^*L_{\mathcal{O}_\mathcal{B}/\mathcal{O}_{\mathcal{B}'}}, \mathcal{G})
\to
\Ext^1_\mathcal{O}(
L_{\mathcal{O}/\mathcal{O}_\mathcal{B}}, \mathcal{G})
% P11-E0040：冻结源把障碍项的次数写成 Ext^1，而不是 Ext^2。
$$
因此，取 $\xi$ 为 $\xi'$ 的像即可。
\end{proof}












\section{概形态射的余切复形}
\label{cotangent-section-cotangent-morphism-schemes}

\noindent
如上所述，我们如下定义概形态射的余切复形。

\begin{definition}
\label{cotangent-definition-cotangent-morphism-schemes}
设 $f : X \to Y$ 为概形态射。$X$ 在 $Y$ 上的{\it 余切复形
$L_{X/Y}$} 是把 $f$ 视为赋环空间态射时的余切复形
（定义 \ref{cotangent-definition-cotangent-complex-morphism-ringed-spaces}）。
\end{definition}

\noindent
特别地，第 \ref{cotangent-section-cotangent-morphism-ringed-spaces} 节的结果
适用于概形态射的余切复形。下一个引理表明，此定义与环同态的定义相容，
并且还蕴含 $L_{X/Y}$ 是 $D_\QCoh(\mathcal{O}_X)$ 的对象。

\begin{lemma}
\label{cotangent-lemma-morphism-affine-schemes}
设 $f : X \to Y$ 为概形态射。设 $U = \Spec(B) \subset X$ 与
$V = \Spec(A) \subset Y$ 为满足 $f(U) \subset V$ 的仿射开集。
存在复形的典范态射
$$
\widetilde{L_{B/A}} \longrightarrow L_{X/Y}|_U
$$
，它在 $D(\mathcal{O}_U)$ 中是同构。此态射与限制到
$X$ 和 $Y$ 中更小的仿射开集相容。
\end{lemma}

\begin{proof}
由评注 \ref{cotangent-remark-map-sections-over-U}，存在 $B=\mathcal{O}_X(U)$-模
复形的典范态射
$L_{\mathcal{O}_X(U)/f^{-1}\mathcal{O}_Y(U)} \to L_{X/Y}(U)$
，它与进一步的限制相容。利用典范态射
$A \to f^{-1}\mathcal{O}_Y(U)$，得到 $B$-模复形的典范态射
$L_{B/A} \to L_{\mathcal{O}_X(U)/f^{-1}\mathcal{O}_Y(U)}$
。利用 $\widetilde{\ }$ 函子的泛性质（参见《概形》引理
\ref{schemes-lemma-compare-constructions}），得到引理陈述中的态射。
要检验此态射在上同调层上为同构，只须检验它在茎上诱导同构。
这立即由引理 \ref{cotangent-lemma-stalk-cotangent-complex} 与
\ref{cotangent-lemma-localize} 得到（还使用了如下茎的描述：
$\mathcal{O}_X$ 与 $f^{-1}\mathcal{O}_Y$
在点 $\mathfrak p \in \Spec(B)$ 处分别为 $B_\mathfrak p$ 与
$A_\mathfrak q$，其中 $\mathfrak q = A \cap \mathfrak p$；
所用参考结果为《概形》引理 \ref{schemes-lemma-spec-sheaves} 和
《层》引理 \ref{sheaves-lemma-stalk-pullback}）。
\end{proof}

\begin{lemma}
\label{cotangent-lemma-scheme-over-ring}
设 $\Lambda$ 为环，$X$ 为 $\Lambda$ 上的概形。则
$$
L_{X/\Spec(\Lambda)} = L_{\mathcal{O}_X/\underline{\Lambda}}
$$
，其中 $\underline{\Lambda}$ 是 $X$ 上取值为 $\Lambda$ 的常值层。
\end{lemma}

\begin{proof}
令 $p : X \to \Spec(\Lambda)$ 为结构态射，并令
$q : \Spec(\Lambda) \to (*, \Lambda)$ 为显然态射。由引理
\ref{cotangent-lemma-triangle-ringed-spaces} 的可区别三角，只须证明 $L_q=0$。
为此，只须对 $\mathfrak p \in \Spec(\Lambda)$ 证明
$$
(L_q)_\mathfrak p =
L_{\mathcal{O}_{\Spec(\Lambda), \mathfrak p}/\Lambda} =
L_{\Lambda_\mathfrak p/\Lambda}
$$
为零（引理 \ref{cotangent-lemma-stalk-cotangent-complex}），而这由引理
\ref{cotangent-lemma-when-zero} 得到。
\end{proof}






\section{环上概形的余切复形}
\label{cotangent-section-cotangent-schemes-variant}

\noindent
设 $\Lambda$ 为环，$X$ 为 $\Lambda$ 上的概形。写作
$L_{X/\Spec(\Lambda)} = L_{X/\Lambda}$，其合理性由引理
\ref{cotangent-lemma-scheme-over-ring} 保证。本节给出一个类似于引理
\ref{cotangent-lemma-compute-cotangent-complex} 的 $L_{X/\Lambda}$ 描述。
确切地说，我们构造纤维化于 $X_{Zar}$ 上的范畴
$\mathcal{C}_{X/\Lambda}$，并在其上赋予一个（多项式）
$\Lambda$-代数层 $\mathcal{O}$，使得
$$
L_{X/\Lambda} =
L\pi_!(\Omega_{\mathcal{O}/\underline{\Lambda}} \otimes_\mathcal{O}
\underline{\mathcal{O}}_X).
$$
稍后将利用范畴 $\mathcal{C}_{X/\Lambda}$ 构造凝聚层叠的朴素障碍理论。

\medskip\noindent
设 $\Lambda$ 为环，$X$ 为 $\Lambda$ 上的概形。令
$\mathcal{C}_{X/\Lambda}$ 为以下交换图所构成的范畴：
\begin{equation}
\label{cotangent-equation-object}
\vcenter{
\xymatrix{
X \ar[d] & U \ar[l] \ar[d] \\
\Spec(\Lambda) & \mathbf{A} \ar[l]
}
}
\end{equation}
，其中
\begin{enumerate}
\item $U$ 是 $X$ 的开子概形；
\item 存在同构 $\mathbf{A} = \Spec(P)$，其中 $P$ 是 $\Lambda$
上以某个集合为变量集的多项式代数。
\end{enumerate}
换言之，$\mathbf{A}$ 是 $\Spec(\Lambda)$ 上的（无限维）仿射空间。
态射由交换图给出。回忆 $X_{Zar}$ 表示 $X$ 的小 Zariski 站点。
存在遗忘函子
$$
u : \mathcal{C}_{X/\Lambda} \to X_{Zar},\ (U \to \mathbf{A}) \mapsto U
$$
注意，$U$ 上的纤维范畴典范等价于第
\ref{cotangent-section-compute-L-pi-shriek} 节引入的范畴
$\mathcal{C}_{\mathcal{O}_X(U)/\Lambda}$。

\begin{lemma}
\label{cotangent-lemma-category-fibred}
在上述情形中，范畴 $\mathcal{C}_{X/\Lambda}$ 纤维化于 $X_{Zar}$ 上。
\end{lemma}

\begin{proof}
给定 $\mathcal{C}_{X/\Lambda}$ 的对象 $U \to \mathbf{A}$ 和
$X_{Zar}$ 的态射 $U' \to U$，考虑 $\mathcal{C}_{X/\Lambda}$ 的对象
$U' \to \mathbf{A}$，其中 $U' \to \mathbf{A}$ 是
$U' \to U$ 与 $U \to \mathbf{A}$ 的复合。$\mathcal{C}_{X/\Lambda}$ 中的态射
$(U' \to \mathbf{A}) \to (U \to \mathbf{A})$
在 $X_{Zar}$ 上是强笛卡儿的。
\end{proof}

\noindent
在 $\mathcal{C}_{X/\Lambda}$ 上赋予从 $X_{Zar}$ 继承的拓扑
（参见《叠》第 \ref{stacks-section-topology} 节）。函子 $u$ 定义拓扑斯态射
$\pi : \Sh(\mathcal{C}_{X/\Lambda}) \to \Sh(X_{Zar})$。
站点 $\mathcal{C}_{X/\Lambda}$ 带有若干环层：
\begin{enumerate}
\item 由规则
$(U \to \mathbf{A}) \mapsto \Gamma(\mathbf{A}, \mathcal{O}_\mathbf{A})$。
给出的层 $\mathcal{O}$；
\item 由规则 $(U \to \mathbf{A}) \mapsto \mathcal{O}_X(U)$ 给出的层
$\underline{\mathcal{O}}_X = \pi^{-1}\mathcal{O}_X$；
\item 常值层 $\underline{\Lambda}$。
\end{enumerate}
得到赋环拓扑斯态射
\begin{equation}
\label{cotangent-equation-pi-schemes}
\vcenter{
\xymatrix{
(\Sh(\mathcal{C}_{X/\Lambda}), \underline{\mathcal{O}}_X) \ar[r]_i \ar[d]_\pi &
(\Sh(\mathcal{C}_{X/\Lambda}), \mathcal{O}) \\
(\Sh(X_{Zar}), \mathcal{O}_X)
}
}
\end{equation}
态射 $i$ 在底层拓扑斯上是恒等态射，而
$i^\sharp : \mathcal{O} \to \underline{\mathcal{O}}_X$ 是显然同态。
态射 $\pi$ 是《站点上的上同调》情形
\ref{sites-cohomology-situation-fibred-category} 的特殊情形。
以下导出函子将在后文发挥重要作用：
$
Li^* : D(\mathcal{O}) \longrightarrow D(\underline{\mathcal{O}}_X)
$
左伴随于 $Ri_* = i_* : D(\underline{\mathcal{O}}_X) \to D(\mathcal{O})$；以及
$
L\pi_! : D(\underline{\mathcal{O}}_X) \longrightarrow D(\mathcal{O}_X)
$
左伴随于
$\pi^* = \pi^{-1} : D(\mathcal{O}_X) \to D(\underline{\mathcal{O}}_X)$。
借助前面的工作，可以计算 $L\pi_!$。

\begin{remark}
\label{cotangent-remark-compute-L-pi-shriek}
在上述情形中，对每个开集 $U \subset X$，令 $P_{\bullet,U}$ 为
$\mathcal{O}_X(U)$ 在 $\Lambda$ 上的标准预解，并令
$\mathbf{A}_{n,U}=\Spec(P_{n,U})$。于是
$\mathbf{A}_{\bullet, U}$
是 $\mathcal{C}_{X/\Lambda}$ 在 $U$ 上的纤维范畴
$\mathcal{C}_{\mathcal{O}_X(U)/\Lambda}$ 的余单纯形对象。此外，
如评注 \ref{cotangent-remark-resolution} 所讨论，
$\mathbf{A}_{\bullet,U}$ 是《站点上的上同调》引理
\ref{sites-cohomology-lemma-compute-by-cosimplicial-resolution}
中那样的 $\mathcal{C}_{\mathcal{O}_X(U)/\Lambda}$ 余单纯形对象。
由于构造 $U \mapsto \mathbf{A}_{\bullet,U}$ 对 $U$ 具有函子性，
给定 $\mathcal{C}_{X/\Lambda}$ 上的任一（阿贝尔）层 $\mathcal{F}$，
得到预层复形
$$
U \longmapsto \mathcal{F}(\mathbf{A}_{\bullet, U})
$$
，其上同调群计算 $\mathcal{F}$ 在该纤维范畴上的同调。由
《站点上的上同调》引理
\ref{sites-cohomology-lemma-compute-left-derived-pi-shriek}，可知其层化
计算 $L_n\pi_!(\mathcal{F})$。换言之，若一个层复形在次数 $-n$ 的项
是 $U \mapsto \mathcal{F}(\mathbf{A}_{n,U})$ 的层化，则该复形计算
$L\pi_!(\mathcal{F})$。
\end{remark}

\noindent
有了这一评注，现在可以陈述本节的主要结果。

\begin{lemma}
\label{cotangent-lemma-cotangent-morphism-schemes}
在上述情形中，$D(\mathcal{O}_X)$ 中存在典范同构
$$
L_{X/\Lambda} = 
L\pi_!(Li^*\Omega_{\mathcal{O}/\underline{\Lambda}}) =
L\pi_!(i^*\Omega_{\mathcal{O}/\underline{\Lambda}}) =
L\pi_!(\Omega_{\mathcal{O}/\underline{\Lambda}}
\otimes_\mathcal{O} \underline{\mathcal{O}}_X)
$$
\end{lemma}

\begin{proof}
首先注意，对 $\mathcal{C}_{X/\Lambda}$ 的任一对象
$(U \to \mathbf{A})$，层 $\mathcal{O}$ 的值都是 $\Lambda$ 上的多项式代数。
因此 $\Omega_{\mathcal{O}/\underline{\Lambda}}$ 是平坦
$\mathcal{O}$-模，从而引理陈述中的第二、第三个等号成立。

\medskip\noindent
由评注 \ref{cotangent-remark-compute-L-pi-shriek}，对象
$L\pi_!(\Omega_{\mathcal{O}/\underline{\Lambda}}
\otimes_\mathcal{O} \underline{\mathcal{O}}_X)$
由预层复形
$$
U \mapsto
\left(\Omega_{\mathcal{O}/\underline{\Lambda}}
\otimes_\mathcal{O} \underline{\mathcal{O}}_X\right)(\mathbf{A}_{\bullet, U})
=
\Omega_{P_{\bullet, U}/\Lambda} \otimes_{P_{\bullet, U}} \mathcal{O}_X(U) =
L_{\mathcal{O}_X(U)/\Lambda}
$$
的层化计算，其中使用评注 \ref{cotangent-remark-compute-L-pi-shriek} 的记号。
现在评注 \ref{cotangent-remark-map-sections-over-U} 表明
$L\pi_!(\Omega_{\mathcal{O}/\underline{\Lambda}}
\otimes_\mathcal{O} \underline{\mathcal{O}}_X)$
计算 $X$ 上环层同态
$\underline{\Lambda} \to \mathcal{O}_X$ 的余切复形。
由引理 \ref{cotangent-lemma-scheme-over-ring}，这正是所需结论。
\end{proof}








\section{代数空间态射的余切复形}
\label{cotangent-section-cotangent-morphism-spaces}

\noindent
我们利用相应的小 \'etale 站点之间的态射定义代数空间态射的余切复形。

\begin{definition}
\label{cotangent-definition-cotangent-morphism-spaces}
设 $S$ 为概形，$f : X \to Y$ 为 $S$ 上代数空间的态射。
$X$ 在 $Y$ 上的{\it 余切复形 $L_{X/Y}$}，是 $X$ 与 $Y$ 的小
\'etale 站点之间赋环拓扑斯态射 $f_{small}$ 的余切复形
（参见《空间的性质》引理
\ref{spaces-properties-lemma-morphism-ringed-topoi} 及定义
\ref{cotangent-definition-cotangent-complex-morphism-ringed-topoi}）。
\end{definition}

\noindent
特别地，第 \ref{cotangent-section-cotangent-morphism-ringed-topoi} 节的结果
适用于代数空间态射的余切复形。以下引理表明，此定义与环同态及概形的定义
相容，并且 $L_{X/Y}$ 是 $D_\QCoh(\mathcal{O}_X)$ 的对象。

\begin{lemma}
\label{cotangent-lemma-etale-localization}
设 $S$ 为概形。考虑 $S$ 上代数空间的交换图
$$
\xymatrix{
U \ar[d]_p \ar[r]_g & V \ar[d]^q \\
X \ar[r]^f & Y
}
$$
，其中 $p$ 与 $q$ 是 \'etale 的。则在 $D(\mathcal{O}_U)$ 中有典范等同
$L_{X/Y}|_{U_\etale} = L_{U/V}$。
\end{lemma}

\begin{proof}
余切复形的形成与拉回交换（引理
\ref{cotangent-lemma-pullback-cotangent-morphism-topoi}），并且
$p_{small}^{-1}\mathcal{O}_X = \mathcal{O}_U$，且
$g_{small}^{-1}\mathcal{O}_{V_\etale} =
p_{small}^{-1}f_{small}^{-1}\mathcal{O}_{Y_\etale}$
，因为 $q_{small}^{-1}\mathcal{O}_{Y_\etale} =
\mathcal{O}_{V_\etale}$
（《空间的性质》引理
\ref{spaces-properties-lemma-etale-exact-pullback}）。追踪各定义即得
$L_{X/Y}|_{U_\etale} = L_{U/V}$。
\end{proof}

\begin{lemma}
\label{cotangent-lemma-compare-spaces-schemes}
设 $S$ 为概形，$f : X \to Y$ 为 $S$ 上代数空间的态射。假设 $X$ 与
$Y$ 可分别由概形 $X_0$ 与 $Y_0$ 表示。则在 $D(\mathcal{O}_X)$ 中
有典范等同 $L_{X/Y} = \epsilon^*L_{X_0/Y_0}$，其中 $\epsilon$
如《空间的导出范畴》第
\ref{spaces-perfect-section-derived-quasi-coherent-etale} 节所定义，而
$L_{X_0/Y_0}$ 如定义 \ref{cotangent-definition-cotangent-morphism-schemes} 所定义。
% P11-E0041：冻结源英文在 representable 前遗漏系动词 are。
\end{lemma}

\begin{proof}
令 $f_0 : X_0 \to Y_0$ 为对应于 $f$ 的概形态射。存在典范态射
$\epsilon^{-1}f_0^{-1}\mathcal{O}_{Y_0} \to f_{small}^{-1}\mathcal{O}_Y$
，它与
$\epsilon^\sharp : \epsilon^{-1}\mathcal{O}_{X_0} \to \mathcal{O}_X$
相容，因为有交换图
$$
\xymatrix{
X_{0, Zar} \ar[d]_{f_0} & X_\etale \ar[l]^\epsilon \ar[d]^f \\
Y_{0, Zar} & Y_\etale \ar[l]_\epsilon
}
$$
，参见《空间的导出范畴》评注
\ref{spaces-perfect-remark-match-total-direct-images}。因此，由第
\ref{cotangent-section-cotangent-complex} 节所讨论的函子性及引理
\ref{cotangent-lemma-pullback-cotangent-morphism-topoi}，得到典范态射
$$
\epsilon^{-1}L_{X_0/Y_0} =
\epsilon^{-1}L_{\mathcal{O}_{X_0}/f_0^{-1}\mathcal{O}_{Y_0}} =
L_{\epsilon^{-1}\mathcal{O}_{X_0}/\epsilon^{-1}f_0^{-1}\mathcal{O}_{Y_0}}
\longrightarrow
L_{\mathcal{O}_X/f^{-1}_{small}\mathcal{O}_Y} = L_{X/Y}
$$
。为证明诱导态射 $\epsilon^*L_{X_0/Y_0} \to L_{X/Y}$ 是同构，
可以在几何点处的茎上检验（《空间的性质》定理
\ref{spaces-properties-theorem-exactness-stalks}）。我们将用引理
\ref{cotangent-lemma-stalk-cotangent-complex} 计算各茎。设
$\overline{x} : \Spec(k) \to X_0$ 为位于 $x \in X_0$ 上方的几何点，
并设 $\overline{y} = f \circ \overline{x}$ 位于 $y \in Y_0$ 上方。则
$$
L_{X/Y, \overline{x}} =
L_{\mathcal{O}_{X, \overline{x}}/\mathcal{O}_{Y, \overline{y}}}
$$
且
$$
(\epsilon^*L_{X_0/Y_0})_{\overline{x}} =
L_{X_0/Y_0, x} \otimes_{\mathcal{O}_{X_0, x}}
\mathcal{O}_{X, \overline{x}} =
L_{\mathcal{O}_{X_0, x}/\mathcal{O}_{Y_0, y}}
\otimes_{\mathcal{O}_{X_0, x}} \mathcal{O}_{X, \overline{x}}
$$
省略一些细节（提示：使用拉回的茎等于在像点处的茎，参见《站点》引理
\ref{sites-lemma-point-morphism-sites}；还要使用模的相应结果，参见
《站点上的模》引理 \ref{sites-modules-lemma-pullback-stalk}）。
注意，$\mathcal{O}_{X,\overline{x}}$ 是
$\mathcal{O}_{X_0,x}$ 的严格 Hensel 化；对
$\mathcal{O}_{Y,\overline{y}}$ 亦然（《空间的性质》引理
\ref{spaces-properties-lemma-describe-etale-local-ring}）。
因此结论由引理 \ref{cotangent-lemma-cotangent-complex-henselization} 得到。
\end{proof}

\begin{lemma}
\label{cotangent-lemma-space-over-ring}
设 $\Lambda$ 为环，$X$ 为 $\Lambda$ 上的代数空间。则
$$
L_{X/\Spec(\Lambda)} = L_{\mathcal{O}_X/\underline{\Lambda}}
$$
，其中 $\underline{\Lambda}$ 是 $X_\etale$ 上取值为 $\Lambda$ 的常值层。
\end{lemma}

\begin{proof}
令 $p : X \to \Spec(\Lambda)$ 为结构态射，并令
$q : \Spec(\Lambda)_\etale \to (*, \Lambda)$ 为显然态射。由引理
\ref{cotangent-lemma-triangle-ringed-topoi} 的可区别三角，只须证明 $L_q=0$。
为此，根据《空间的性质》定理
\ref{spaces-properties-theorem-exactness-stalks}，只须对几何点
$\overline{t} : \Spec(k) \to \Spec(\Lambda)$ 证明
$$
(L_q)_{\overline{t}} =
L_{\mathcal{O}_{\Spec(\Lambda)_\etale, \overline{t}}/\Lambda}
$$
为零（引理 \ref{cotangent-lemma-stalk-cotangent-complex}）。由于
$\mathcal{O}_{\Spec(\Lambda)_\etale, \overline{t}}$ 是
$\Lambda$ 的某个局部环的严格 Hensel 化（《空间的性质》引理
\ref{spaces-properties-lemma-describe-etale-local-ring}），
结论由引理 \ref{cotangent-lemma-when-zero} 得到。
\end{proof}










\section{环上代数空间的余切复形}
\label{cotangent-section-cotangent-spaces-variant}

\noindent
设 $\Lambda$ 为环，$X$ 为 $\Lambda$ 上的代数空间。写作
$L_{X/\Spec(\Lambda)} = L_{X/\Lambda}$，其合理性由引理
\ref{cotangent-lemma-space-over-ring} 保证。本节给出一个类似于引理
\ref{cotangent-lemma-compute-cotangent-complex} 的 $L_{X/\Lambda}$ 描述。
确切地说，我们构造纤维化于 $X_\etale$ 上的范畴
$\mathcal{C}_{X/\Lambda}$，并在其上赋予一个（多项式）
$\Lambda$-代数层 $\mathcal{O}$，使得
$$
L_{X/\Lambda} =
L\pi_!(\Omega_{\mathcal{O}/\underline{\Lambda}} \otimes_\mathcal{O}
\underline{\mathcal{O}}_X).
$$
稍后将利用范畴 $\mathcal{C}_{X/\Lambda}$ 构造凝聚层叠的朴素障碍理论。

\medskip\noindent
设 $\Lambda$ 为环，$X$ 为 $\Lambda$ 上的代数空间。令
$\mathcal{C}_{X/\Lambda}$ 为以下交换图所构成的范畴：
\begin{equation}
\label{cotangent-equation-object-space}
\vcenter{
\xymatrix{
X \ar[d] & U \ar[l] \ar[d] \\
\Spec(\Lambda) & \mathbf{A} \ar[l]
}
}
\end{equation}
，其中
\begin{enumerate}
\item $U$ 是概形；
\item $U \to X$ 是 \'etale 态射；
\item 存在同构 $\mathbf{A}=\Spec(P)$，其中 $P$ 是 $\Lambda$
上以某个集合为变量集的多项式代数。
\end{enumerate}
换言之，$\mathbf{A}$ 是 $\Spec(\Lambda)$ 上的（无限维）仿射空间。
态射由交换图给出。回忆 $X_\etale$ 表示 $X$ 的小 \'etale 站点，
其对象是 \'etale 于 $X$ 的概形。存在遗忘函子
$$
u : \mathcal{C}_{X/\Lambda} \to X_\etale,
\quad
(U \to \mathbf{A}) \mapsto U
$$
注意，$U$ 上的纤维范畴典范等价于第
\ref{cotangent-section-compute-L-pi-shriek} 节引入的范畴
$\mathcal{C}_{\mathcal{O}_X(U)/\Lambda}$。

\begin{lemma}
\label{cotangent-lemma-category-fibred-space}
在上述情形中，范畴 $\mathcal{C}_{X/\Lambda}$ 纤维化于 $X_\etale$ 上。
\end{lemma}

\begin{proof}
给定 $\mathcal{C}_{X/\Lambda}$ 的对象 $U \to \mathbf{A}$ 和
$X_\etale$ 的态射 $U' \to U$，考虑 $\mathcal{C}_{X/\Lambda}$ 的对象
$U' \to \mathbf{A}$，其中 $U' \to \mathbf{A}$ 是
$U' \to U$ 与 $U \to \mathbf{A}$ 的复合。态射
$(U' \to \mathbf{A}) \to (U \to \mathbf{A})$
在 $X_\etale$ 上是强笛卡儿的。
\end{proof}

\noindent
在 $\mathcal{C}_{X/\Lambda}$ 上赋予从 $X_\etale$ 继承的拓扑
（参见《叠》第 \ref{stacks-section-topology} 节）。函子 $u$ 定义拓扑斯态射
$\pi : \Sh(\mathcal{C}_{X/\Lambda}) \to \Sh(X_\etale)$。
站点 $\mathcal{C}_{X/\Lambda}$ 带有若干环层：
\begin{enumerate}
\item 由规则
$(U \to \mathbf{A}) \mapsto \Gamma(\mathbf{A}, \mathcal{O}_\mathbf{A})$。
给出的层 $\mathcal{O}$；
\item 由规则 $(U \to \mathbf{A}) \mapsto \mathcal{O}_X(U)$ 给出的层
$\underline{\mathcal{O}}_X = \pi^{-1}\mathcal{O}_X$；
\item 常值层 $\underline{\Lambda}$。
\end{enumerate}
得到赋环拓扑斯态射
\begin{equation}
\label{cotangent-equation-pi-spaces}
\vcenter{
\xymatrix{
(\Sh(\mathcal{C}_{X/\Lambda}), \underline{\mathcal{O}}_X) \ar[r]_i \ar[d]_\pi &
(\Sh(\mathcal{C}_{X/\Lambda}), \mathcal{O}) \\
(\Sh(X_\etale), \mathcal{O}_X)
}
}
\end{equation}
态射 $i$ 在底层拓扑斯上是恒等态射，而
$i^\sharp : \mathcal{O} \to \underline{\mathcal{O}}_X$ 是显然同态。
态射 $\pi$ 是《站点上的上同调》情形
\ref{sites-cohomology-situation-fibred-category} 的特殊情形。
以下导出函子将在后文发挥重要作用：
$
Li^* : D(\mathcal{O}) \longrightarrow D(\underline{\mathcal{O}}_X)
$
左伴随于 $Ri_* = i_* : D(\underline{\mathcal{O}}_X) \to D(\mathcal{O})$；以及
$
L\pi_! : D(\underline{\mathcal{O}}_X) \longrightarrow D(\mathcal{O}_X)
$
左伴随于
$\pi^* = \pi^{-1} : D(\mathcal{O}_X) \to D(\underline{\mathcal{O}}_X)$。
借助前面的工作，可以计算 $L\pi_!$。

\begin{remark}
\label{cotangent-remark-compute-L-pi-shriek-spaces}
在上述情形中，对 $X_\etale$ 的每个对象 $U \to X$，令
$P_{\bullet,U}$ 为 $\mathcal{O}_X(U)$ 在 $\Lambda$ 上的标准预解，并令
$\mathbf{A}_{n,U}=\Spec(P_{n,U})$。于是 $\mathbf{A}_{\bullet,U}$
是 $\mathcal{C}_{X/\Lambda}$ 在 $U$ 上的纤维范畴
$\mathcal{C}_{\mathcal{O}_X(U)/\Lambda}$ 的余单纯形对象。此外，
如评注 \ref{cotangent-remark-resolution} 所讨论，
$\mathbf{A}_{\bullet,U}$ 是《站点上的上同调》引理
\ref{sites-cohomology-lemma-compute-by-cosimplicial-resolution}
中那样的 $\mathcal{C}_{\mathcal{O}_X(U)/\Lambda}$ 余单纯形对象。
由于构造 $U \mapsto \mathbf{A}_{\bullet,U}$ 对 $U$ 具有函子性，
给定 $\mathcal{C}_{X/\Lambda}$ 上的任一（阿贝尔）层 $\mathcal{F}$，
得到预层复形
$$
U \longmapsto \mathcal{F}(\mathbf{A}_{\bullet, U})
$$
，其上同调群计算 $\mathcal{F}$ 在该纤维范畴上的同调。由
《站点上的上同调》引理
\ref{sites-cohomology-lemma-compute-left-derived-pi-shriek}，可知其层化
计算 $L_n\pi_!(\mathcal{F})$。换言之，若一个层复形在次数 $-n$ 的项
是 $U \mapsto \mathcal{F}(\mathbf{A}_{n,U})$ 的层化，则该复形计算
$L\pi_!(\mathcal{F})$。
\end{remark}

\noindent
有了这一评注，现在可以陈述本节的主要结果。

\begin{lemma}
\label{cotangent-lemma-cotangent-morphism-spaces}
在上述情形中，$D(\mathcal{O}_X)$ 中存在典范同构
$$
L_{X/\Lambda} = 
L\pi_!(Li^*\Omega_{\mathcal{O}/\underline{\Lambda}}) =
L\pi_!(i^*\Omega_{\mathcal{O}/\underline{\Lambda}}) =
L\pi_!(\Omega_{\mathcal{O}/\underline{\Lambda}}
\otimes_\mathcal{O} \underline{\mathcal{O}}_X)
$$
\end{lemma}

\begin{proof}
首先注意，对 $\mathcal{C}_{X/\Lambda}$ 的任一对象
$(U \to \mathbf{A})$，层 $\mathcal{O}$ 的值都是 $\Lambda$ 上的多项式代数。
因此 $\Omega_{\mathcal{O}/\underline{\Lambda}}$ 是平坦
$\mathcal{O}$-模，从而引理陈述中的第二、第三个等号成立。

\medskip\noindent
由评注 \ref{cotangent-remark-compute-L-pi-shriek-spaces}，对象
$L\pi_!(\Omega_{\mathcal{O}/\underline{\Lambda}}
\otimes_\mathcal{O} \underline{\mathcal{O}}_X)$
由预层复形
$$
U \mapsto
\left(\Omega_{\mathcal{O}/\underline{\Lambda}}
\otimes_\mathcal{O} \underline{\mathcal{O}}_X\right)(\mathbf{A}_{\bullet, U})
=
\Omega_{P_{\bullet, U}/\Lambda} \otimes_{P_{\bullet, U}} \mathcal{O}_X(U) =
L_{\mathcal{O}_X(U)/\Lambda}
$$
的层化计算，其中使用评注 \ref{cotangent-remark-compute-L-pi-shriek-spaces} 的记号。
现在评注 \ref{cotangent-remark-map-sections-over-U} 表明
$L\pi_!(\Omega_{\mathcal{O}/\underline{\Lambda}}
\otimes_\mathcal{O} \underline{\mathcal{O}}_X)$
计算 $X_\etale$ 上环层同态
$\underline{\Lambda} \to \mathcal{O}_X$ 的余切复形。
由引理 \ref{cotangent-lemma-space-over-ring}，这正是所需结论。
\end{proof}






\section{代数空间的纤维积与余切复形}
\label{cotangent-section-fibre-product}

\noindent
设 $S$ 为概形，$X \to B$ 与 $Y \to B$ 为 $S$ 上代数空间的态射。
考虑纤维积 $X \times_B Y$，其投影态射为
$p : X \times_B Y \to X$ 与 $q : X \times_B Y \to Y$。
本节讨论 $L_{X \times_B Y/B}$。我们所需的大部分信息都包含在下图中：
\begin{equation}
\label{cotangent-equation-fibre-product}
\vcenter{
\xymatrix{
Lp^*L_{X/B} \ar[r] &
L_{X \times_B Y/Y} \ar[r] &
E \\
Lp^*L_{X/B} \ar[r] \ar@{=}[u] &
L_{X \times_B Y/B} \ar[r] \ar[u] &
L_{X \times_B Y/X} \ar[u] \\
 &
Lq^*L_{Y/B} \ar[u] \ar@{=}[r] &
Lq^*L_{Y/B} \ar[u]
}
}
\end{equation}
说明：中间一行是引理 \ref{cotangent-lemma-triangle-ringed-topoi} 中态射
$X \times_B Y \to X \to B$ 的基本三角。中间一列是态射
$X \times_B Y \to Y \to B$ 的基本三角。其次，$E$ 是
$D(\mathcal{O}_{X \times_B Y})$ 中“嵌入”右上角的对象，也就是说，
它使最上面一行和最右边一列都成为可区别三角。将《导出范畴》命题
\ref{derived-proposition-9} 应用于左下方块（在缺失位置放置 $0$），
可知这样的 $E$ 存在。更明确地说，例如可以把 $E$ 定义为复形态射
$$
Lp^*L_{X/B} \oplus Lq^*L_{Y/B} \longrightarrow L_{X \times_B Y/B}
$$
的锥（《导出范畴》定义 \ref{derived-definition-cone}），再应用 TR3
得到以 $E$ 为靶的两个态射。在 Tor 独立的情形下，对象 $E$ 为零。

\begin{lemma}
\label{cotangent-lemma-fibre-product-tor-independent}
在上述情形中，若 $X$ 与 $Y$ 在 $B$ 上 Tor 独立，则
(\ref{cotangent-equation-fibre-product}) 中的对象 $E$ 为零。在此情形下有
$$
L_{X \times_B Y/B} = Lp^*L_{X/B} \oplus Lq^*L_{Y/B}
$$
\end{lemma}

\begin{proof}
选取概形 $W$ 和满 \'etale 态射 $W \to B$；选取概形 $U$ 和满
\'etale 态射 $U \to X \times_B W$；选取概形 $V$ 和满
\'etale 态射 $V \to Y \times_B W$。则
$U \times_W V \to X \times_B Y$ 也是满 \'etale 态射。
因此只须证明 $E$ 在 $U \times_W V$ 上的限制为零。由引理
\ref{cotangent-lemma-compare-spaces-schemes} 及《空间的导出范畴》引理
\ref{spaces-perfect-lemma-tor-independent}，归约到概形情形。
取适当的仿射开集后，进一步归约到仿射概形情形。再利用引理
\ref{cotangent-lemma-morphism-affine-schemes}，归约到环的张量积情形，
也就是引理 \ref{cotangent-lemma-tensor-product-tor-independent}。
\end{proof}

\noindent
一般地，对象 $E$ 有如下性质。

\begin{lemma}
\label{cotangent-lemma-fibre-product}
设 $S$ 为概形，$X \to B$ 与 $Y \to B$ 为 $S$ 上代数空间的态射。
(\ref{cotangent-equation-fibre-product}) 中的对象 $E$ 对 $i=0,-1$ 满足
$H^i(E)=0$，而对几何点
$(\overline{x}, \overline{y}) : \Spec(k) \to X \times_B Y$ 有
$$
H^{-2}(E)_{(\overline{x}, \overline{y})} =
\text{Tor}_1^R(A, B) \otimes_{A \otimes_R B} C
$$
，其中 $R = \mathcal{O}_{B, \overline{b}}$、
$A = \mathcal{O}_{X, \overline{x}}$、
$B = \mathcal{O}_{Y, \overline{y}}$，且
$C = \mathcal{O}_{X \times_B Y, (\overline{x}, \overline{y})}$。
\end{lemma}

\begin{proof}
余切复形的形成与取茎及拉回交换，参见引理
\ref{cotangent-lemma-stalk-cotangent-complex} 与
\ref{cotangent-lemma-pullback-cotangent-morphism-topoi}。注意，$C$ 是
$A \otimes_R B$ 的 Hensel 化。由第 \ref{cotangent-section-localization} 节的结果，
$L_{C/R} = L_{A \otimes_R B/R} \otimes_{A \otimes_R B} C$
。因此，$E$ 在该几何点处的茎是态射
$L_{A/R} \otimes C \to L_{A \otimes_R B/R} \otimes C$ 的锥。
所以，引理的结论由环的情形，即引理 \ref{cotangent-lemma-tensor-product} 得到。
\end{proof}
